



	
	\chapter{Aprendizaje Supervisado}{\label{Cap2}}

	      %  Capitulo 2 	
	

	%%  Applied Spatial Data Analysis with R  de Roger S. Bivand (Edzer Pebesma y Virgilio Gómez-Rubio) como  
	%%    Analysing  spatial  point patterns in R de  Adrian Baddeley.\\
	

	%% Sistema de información geografica los ordenadores, como software R y GVSIG,
	

\section{Preliminares}            


El  aprendizaje estadístico  supervisado es uno de los más comunes y exitosos tipos de aprendizaje estadístico  (o machine learning). En este capítulo se describe el aprendizaje supervisado con más detalle y se explican algunos algoritmos comunes mismo. En el Capítulo \ref{Cap1}  vimos una introducción al machine learning, con  la clasificación de las Flores Iris en varias especies, de acuerdo con las medidas físicas de las flores.\\

Recordemos que el aprendizaje supervisado es usado cada vez que queremos predecir cierta salida de una entrada dada, siempre te tengamos ejemplos suficientes de pares de entrada y salida  (input{\scriptsize/}output). Se construye el modelo de aprendizaje automático con estos pares de entrada{\scriptsize/}salida, que se encuentran en  la data de entrenamiento o capacitación. El objetivo es hacer predicciones exactas para datos nuevos, nunca antes vistos o tratados.   El aprendizaje supervisado frecuentemente requiere de un esfuerzo humano para construir un conjunto de entrenamiento,  pero ahora  es automatizado y generalmente facilita y acelera el trabajo, que de otra manera sería muy laboriosa y  en algunos casos inviable.


\subsection{Clasificación y Regresión}


Generalmente en el aprendizaje supervisado se consideran dos tipos de problemas más importantes llamados clasificación y regresión.\\

En el problema de clasificación, el objetivo es predecir una clase etiquetada  seleccionada  de una lista predefinidas de posibilidades. En el  Capítulo \ref{Cap1} se presentó la clasificación de las flores iris dentro de tres tipos de  especie  posible. El problema de clasificación algunas veces es separada en dos tipos, clasificación binaria y clasificación multiclase. La clasificación binaria consta exactamente de dos clases, mientras que la multiclase tiene más de dos clases.  Se puede pensar en una clasificación binaria como una prueba para responder una pregunta con dos respuestas, si o no. Un ejemplo práctico de la clasificación binaria es clasificar un correo como spam o no spam. En esta clasificación la tarea en responder, si o no a la pregunta.\\
	


\begin{minipage}[t]{0.2\textwidth}
	\centering\raisebox{\dimexpr \topskip-\height}{%
		
\includegraphics[width=\textwidth]{Imagen_Machine/Pajaro.png}}
\end{minipage}%\hfill
\begin{minipage}[t]{0.75\textwidth} 
	{\scriptsize	
		En la clasificación binaria frecuentemente se habla de una  clase positiva y una clase		
		negativa.	Aquí, la clase positiva no representa beneficio o valor, sólo es objeto		
		de estudio. Así, que cuando vemos  "positivo"  ( spam),  significa clase spam. Designar		
		una de las dos clases como  positiva  es a menudo una materia subjetiva, y específico		
		para el dominio.}
\end{minipage}




El problema de las flores iris, es un ejemplo de  clasificación multiclase.  Otro típico ejemplo de multiclase, es predecir el lenguaje o idioma en qué  está escrito un sitio Web; las clases en este caso son una lista predefinida de posibles idiomas.  \\

En el caso del problema de regresión, el objetivo es predecir un número de tipo punto-flotante  (floating-point number)  en términos de programación y en términos matemático, un número real. Como ejemplo podemos citar, el ingreso anual de una persona por su educación, por su edad, o por el lugar  donde vive. Cuando se predice el ingreso, el valor es una cantidad, y puede ser cualquier número en un rango dado. Otro ejemplo de regresión es predecir el rendimiento de la cosecha de maíz en una granja, dado los atributos previos como rendimiento previo, tiempo (clima) y número de  trabajadores en la granja. El rendimiento nuevamente es un número tipo punto flotante.\\


Una forma fácil de distinguir entre clasificación y regresión es  preguntar si hay algún  tipo de continuidad en la salida. Si hay continuidad en los posibles resultados de salidas, entonces el problema es regresión.  Pensando en el ejemplo de la predicción del ingreso anual. Aquí, claramente hay continuidad en la salida. Si una persona gana 40.000 o 40.001 en un año no hace una diferencia tangible, aun, cuando esas son cantidades diferentes en dinero; si nuestro algoritmo predice 39.999 o 40.001 cuando debe predecir 40.000, no es tan importante.  \\


Por el contrario, para la tarea de reconocimiento del idioma en el sitio web (problema de clasificación), no hay continuidad, entre idiomas, no hay un idioma entre dos. Un sitio web tiene un sólo idioma{\footnote{Pedimos a los lingüistas que disculpen la presentación simplificada de las lenguas como entidades distintas y fijas.}.   

\subsection{Generalización, Sobreajuste,  y Subajuste}  

En el aprendizaje supervisado se  construye un modelo sobre los datos de entrenamiento  que sea capaz de hacer predicciones exactas sobre datos nuevos, que no se han usado en la construcción del modelo pero que   tienen la misma característica de los datos de entrenamiento.

 Si un modelo es capaz de hacer predicciones exactas sobre 	datos nuevos, decimos que es capaz de generalizar\index{capacidad de generalizar} del conjunto de entrenamiento al conjunto de prueba. Así que lo deseado es construir un modelo capaz de generalizar con la mayor exactitud posible.


Por lo general, creamos un modelo de tal manera que pueda hacer predicciones exactas sobre el conjunto de entrenamiento. Si los conjuntos de entrenamiento  y  prueba  tienen  sus  características  comunes  y  semejantes,  se espera que el modelo  también  sea  exacto  en el conjunto de prueba.  Sin embargo, hay algunos casos donde esto puede salir mal. Por ejemplo, si construimos modelos muy complejos y exactos en el conjunto de entrenamiento, el modelo tal vez puede no generalizar bien. \\


Para ilustrar esta situación usaremos el siguiente ejemplo ideado: Supongamos que un aficionado científico de datos quiere predecir si un cliente comprará o no un bote, basándose en los registros de los compradores  de  botes  y  los clientes que sabemos que no están interesados comprar un bote. El objetivo es enviar correos electrónicos promocionales a las personas que probablemente hagan una compra, pero no molestar a aquellos clientes que no estarán interesados.\\


Supongamos que los registros de clientes son los que se muestran en la Tabla \ref{Clientes}.\\



 %                 Example data about customers
 
\begin{table}[h!]
	\centering
	\caption{Ejemplo de datos sobre clientes}
	\includegraphics[scale=0.65]{../Imagen_Machine/Clientes}
	\label{Clientes}
\end{table} 


Después de analizar los datos, el científico aficionado se le ocurre la siguiente regla:
"Si el cliente es mayor de 45 años y tiene menos de 3 hijos o no es divorciado, entonces 
quiere comprar un bote". Cuando se le preguntó qué tan bien funciona su regla,
el científico de datos responde:  "¡Es 100 por ciento exacta!" {}   Y de hecho, en los datos de 
la  tabla, la regla es perfectamente exacta. Hay muchas reglas posibles que podríamos
generar  que explicaría perfectamente si alguien en este conjunto de datos quiere comprar o no un bote.  La edad no aparece dos veces en los datos, por lo que podríamos decir que para personas con 66, 52, 53 o 58 años de edad quiere comprar un bote, mientras  que los demás no. Si bien podemos inventar muchas reglas que funcionan bien en estos datos, es importante  resaltar que  no estamos interesados en hacer predicciones para este conjunto de datos; pues, ya sabemos las respuestas para estos clientes.   Queremos saber si los  nuevos clientes  probablemente compran  un bote. Por lo tanto, queremos encontrar una regla que funcione bien para nuevos clientes y que logre una exactitud alta en el conjunto de entrenamiento. Sin embargo, una exactitud del 100 por ciento en el conjunto de entrenamiento generalmente  no nos ayuda en esta situación.  No podemos esperar que la regla propuesta por el aficionado científico de datos, funcionara muy bien con datos nuevos.  Parece demasiado complejo y está  basado  en  muy pocos datos.   Por ejemplo, la parte de la regla  "{}o no está divorciado"{}   se observa en un solo cliente.\\


La única medida para conocer  si un algoritmo funcionará  bien  en nuevos datos es la evaluación en el conjunto de prueba.  Sin embargo, intuitivamente esperamos que los modelos simples  generalicen mejor en los  datos nuevos.  Si la regla era que "las personas mayores de 50 años quieren comprar un bote", y esto explicaría el comportamiento de todos los clientes, confiaríamos más en la regla que involucra niños, estado civil, y además la edad. Por lo tanto, siempre queremos
encontrar el modelo más simple.  Construir un modelo que sea demasiado complejo para la cantidad de información que tenemos, como hizo aficionado el científico, se llama sobreajuste\index{sobreajuste}.  El sobreajuste ocurre cuando se ajusta un modelo demasiado  cerca de las particularidades del conjunto de entrenamiento y se obtiene un modelo que funciona muy bien en el conjunto de entrenamiento pero no es capaz de generalizar para nuevos datos. 
Por otro lado,  si su modelo es demasiado simple, digamos: "Todo cliente que poseen un casa propia,  comprará un bote”,  entonces es posible que no sea capaz de capturar todos los aspectos y variaciones  en los datos, y el modelo funcionará mal incluso en el conjunto de entrenamiento. Elegir un modelo demasiado simple se llama subajuste\index{subajuste}.\\


Cuanto más complejo  sea un modelo  mejor será la capacidad de  predecir sobre  los datos de entrenamiento; sin embargo, si el modelo se hace demasiado complejo, comenzamos a enfocarnos más  en cada punto individual de los datos del conjunto de entrenamiento, y el modelo no funcionará bien sobre los nuevos datos (datos de prueba).\\

Hay un modelo dentro de esta  situación (punto justo  o  Sweet spot)  que producirá el mejor rendimiento de generalización; y este es el modelo que queremos encontrar.\\



En la Figura \ref{complejidad} se ilustra la  compensación entre sobreajuste y subajuste  modelo

\begin{figure}[h!]
	\centering
	\includegraphics[scale=0.50]{../Imagen_Machine/complejidad}
	\caption{Compensación de la complejidad del modelo contra  la exactitud  de los conjuntos de  entrenamiento  y  la prueba}
	\label{complejidad}
\end{figure} 



\subsubsection{Relación de la complejidad del modelo con el tamaño del conjunto de datos}

Es importante tener en cuenta que la complejidad del modelo está íntimamente ligada a la variación del número
entradas contenidas en el cojunto de entrenamiento: cuando la mayor sea cantidad de puntos en el conjunto de  datos, un modelo más complejo puede generarse sin sobreajuste. Por lo general, coleccionando más puntos en la  data producirá más información, por lo que los conjuntos con más datos permiten construir modelos más complejos. Sin embargo, simplemente duplicar los mismos puntos de datos o recopilar  datos muy similares no ayudará.\\
 

Volviendo al ejemplo de venta de botes, si vemos 10,000 filas más de datos de clientes,
y todos cumplen con la regla: "Si el cliente es mayor de 45 años y tiene menos
de 3 hijos o no está divorciado, entonces quieren comprar un bote", estaríamos mucho más seguros que esta es una buena regla, que cuando se desarrolló utilizando solo las 12 filas de la tabla \ref{Clientes}.\\

  %   wonders =  marabillas

Tener más datos y construir modelos apropiadamente más complejos a menudo puede funcionar de
maravilla para tareas de aprendizaje supervisado. En este libro, nos centraremos en trabajar con
conjuntos de datos de tamaños fijos. Pero en el mundo real, generalmente hay posibilidad de decidir cuantos datos a usar, lo que podría ser más beneficioso para ajustar y afinar el modelo. Por último, le sugerimos que  "Nunca subestimes el poder de más datos".
                   %  tweaking   = ajustar


\section{Ejemplos de algunos conjuntos de datos}


Para ilustrar la aplicación de los diferentes algoritmos, se utilizarán  algunos conjuntos de datos. Conjuntos de datos pequeños y sintéticos (es decir, indeados), diseñado para resaltar
aspectos muy particulares de los algoritmos; y otros conjuntos de datos grandes, como ejemplos del mundo real.\\

Un ejemplo de un conjunto  de clasificación  con dos clases, es la data  sintética    de falsificación de datos (forge dataset), que tiene dos características.


El siguiente código crea un diagrama de dispersión (Figura \ref{DispersionForge}) que visualiza todos de los puntos de datos de este conjunto. El gráfico tiene la primera característica en el eje x y la segunda característica sobre el eje y. Como siempre  en los diagramas de dispersión, cada punto de datos es representado como un punto. El color y la forma del punto indican su clase:\\

\begin{lstlisting}
In[2]:  # generate dataset
        X, y = mglearn.datasets.make_forge()
        # plot dataset
        mglearn.discrete_scatter(X[:, 0], X[:, 1], y)
        plt.legend(["Class 0", "Class 1"], loc=4)
        plt.xlabel("First feature")
        plt.ylabel("Second feature")
        print("X.shape: {}".format(X.shape))

Out[2]:

    X.shape: (26, 2)
\end{lstlisting}



\begin{figure}[h!]
	\centering
	\includegraphics[scale=0.50]{../Imagen_Machine/DispersionForge}
	\caption{Diagrama de dispersión de la dataset  de falsificación de datos (forge dataset)}
	\label{DispersionForge}
\end{figure} 


Como puede ver en  {\ttfamily X.shape}, este conjunto de datos consta de 26 datos, con 2 características.
Para ilustrar el algoritmo de regresión, utilizaremos el conjunto sintético de datos de Onda.
El conjunto de datos de Onda (wave) tiene una única característica de entrada y una variable objetivo  continua (o respuesta)  que queremos modelar. El diagrama en la figura \ref{DispersionForge}  muestra la característica única en el eje x,  y el objetivo de regresión (la salida) en el eje y:\\

\begin{lstlisting}
In[3]:  X, y = mglearn.datasets.make_wave(n_samples=40)
        plt.plot(X, y, 'o')
        plt.ylim(-3, 3)
        plt.xlabel("Feature")
        plt.ylabel("Target")
\end{lstlisting}


\begin{figure}[h!]
	\centering
	\includegraphics[scale=0.50]{../Imagen_Machine/regresion}
	\caption{Gráfico del conjunto de datos de onda, en el eje $X$ se muestra la
		característica y en  el eje $Y$ el objetivo (respuesta de la regresión).}
	\label{regresion}
\end{figure} 



Utilizamos conjuntos de datos muy simples y de baja dimensión porque se pueden visualizar fácilmente en  dos dimensiones, sin embargo, para base de datos con más de dos características es difícil visualizar. Cualquier intuición derivada de conjuntos de datos con pocas características (también llamada los conjuntos de datos de baja dimensión\index{conjuntos de baja dimensión}) podrían no ayudar en conjuntos de datos con muchas características (conjuntos de datos de alta dimensión\index{conjuntos de alta dimensión}). No obstante, tenga en cuenta, que la inspección de algoritmos en los conjuntos de datos de baja dimensión pueden ser muy instructiva.\\


\begin{minipage}[t]{0.2\textwidth}
	\centering\raisebox{\dimexpr \topskip-\height}{%
		
\includegraphics[width=\textwidth]{Imagen_Machine/Pajaro.png}}
\end{minipage}%\hfill
\begin{minipage}[t]{0.75\textwidth} 
	{\scriptsize	
		Los conjuntos de datos que se incluyen en {\ttfamily scikit-learn} generalmente se almacenan como		
		manojos de objetos, que contienen información sobre el conjunto de datos así como los 		
		datos reales. Todo lo que necesitas saber sobre los objetos tipo manojo de objetos, es que		
		se comportan como diccionarios, con el beneficio adicional de que usted puede acceder a los 		
		valores usando un punto  (como en {\ttfamily bunch.key} en lugar de {\ttfamily bunch['key']}).}
\end{minipage}  \\
%\vspace{0.5cm}

Complementamos estos conjuntos sintéticos y de baja dimensión, con dos conjuntos de datos del mundo real que
están incluidos en {\ttfamily scikit-learn}. Uno de estos es el conjunto de datos del Cáncer de Mama de Wisconsin (Wisconsin Breast Cancer), que registra mediciones clínicas de tumores de cáncer de mama.  Cada tumor está etiquetado como {}"benigno"{}  (para tumores inofensivos) o  {}"maligno"{} (para tumores cancerosos), en este caso la tarea es aprender a predecir si un tumor es maligno basado en las medidas del tejido.\\

Los datos  puede cargarse utilizando la función {\ttfamily{load\_breast\_cancer}} de
 {\ttfamily scikit-learn}:\\

\begin{lstlisting}
In[4]:  from sklearn.datasets import load_breast_cancer
        cancer = load_breast_cancer()
        print("cancer.keys(): \n{}".format(cancer.keys()))

Out[4]:
 
    cancer.keys():
    dict_keys(['feature_names', 'data', 'DESCR', 'target', 'target_names'])

\end{lstlisting}


\noindent La data consiste en 569 puntos o datos, con 30 características cada una:\\


\begin{lstlisting}

In[5]:  print("Forma de la data del cancer: {}".format(cancer.data.shape))

Out[5]:

   Forma de la data del cancer: (569, 30)

 \end{lstlisting} 
 
\noindent  De esos  569 datos, 212  son etiquetados como malignos y 357 como benignos:\\

\begin{lstlisting}

In[6]:  print("Conteos de muestra por clase:\n{}".format(
        {n: v for n, v in zip(cancer.target_names, np.bincount(
               cancer.target))}))
               
Out[6]:

    Conteos de muestra por clase:
    {'benign': 357, 'malignant': 212}
\end{lstlisting}



Para obtener una descripcion semantica del significado de cada característica, podemos verlo con 
el atributo  {\ttfamily feature\_names}:\\

\begin{lstlisting}
In[7]:  print("Nombres de las carateristica:\n{}".format(
                cancer.feature_names))
Out[7]:

    Nombres de las carateristica:
    ['mean radius' 'mean texture' 'mean perimeter' 'mean area'
     'mean smoothness' 'mean compactness' 'mean concavity'
     'mean concave points' 'mean symmetry' 'mean fractal dimension'
     'radius error' 'texture error' 'perimeter error' 'area error'
     'smoothness error' 'compactness error' 'concavity error'
     'concave points error' 'symmetry error' 'fractal dimension error'
     'worst radius' 'worst texture' 'worst perimeter' 'worst area'
     'worst smoothness' 'worst compactness' 'worst concavity'
     'worst concave points' 'worst symmetry' 'worst fractal dimension']
\end{lstlisting}


Si estas interesado en obtener más información al respecto puede ver en 
{\ttfamily{cancer.DESCR}} .\\

También usaremos el conjunto de datos mundo real, el conjunto de datos de Vivienda de Boston (Boston Housing), para la regresión.  La tarea asociada con este conjunto de datos es predecir el valor promedio de las viviendas en algunos barrios de Boston en la década de 1970, utilizando información como la tasa de criminalidad, proximidad al río Charles, accesibilidad a la autopista, etc.  El conjunto de datos contiene 506 puntos de datos, descritos por 13 características:\\

\begin{lstlisting}
In[8]:  from sklearn.datasets import load_boston
        boston = load_boston()
        print("Data shape: {}".format(boston.data.shape))

Out[8]:

    Data shape: (506, 13)
\end{lstlisting}


Nuevamente, puede obtener más información sobre el conjunto de datos en  el atributo {\ttfamily{DESCR}}
de boston. Para nuestros propósitos aquí, expandiremos este conjunto de datos no solo
considera estas 13 medidas como características de entrada, sino  también plantea  todos los productos
entre características (también llamadas interacciones). En otras palabras, no solo consideraremos
tasa de criminalidad y accesibilidad a la autopista como características, sino también el producto de la tasa de criminalidad por la accesibilidad a la autopista.  La inclusión de caractarísticas derivadas o extendidas como estas se denomina ingeniería de características, se discutá en el Capítulo \ref{Cap4} con más detalle. 

El conjunto de datos extendido esta disponible en {\ttfamily mglearn.datasets} y puede ser cargado mediante  la función {\ttfamily load\_extended\_boston}:\\

\begin{lstlisting}
In[9]:  X, y = mglearn.datasets.load_extended_boston()
        print("X.shape: {}".format(X.shape))

Out[9]:

    X.shape: (506, 104)

\end{lstlisting}

 Las 104 características resultantes son las 13 características originales junto con las 91  combinaciones\footnote {Esto se llama coeficiente binomial, que es el número de combinaciones de k elementos que se pueden seleccionar 	de un conjunto de n elementos. A menudo, esto se escribe como $(\frac{n}{k})$	y se expresa como "n elige k"; en este caso, " de 13 se  elige 2".} posibles de las características  de las 13 tomadas de dos en dos.  Utilizaremos estos conjuntos de datos para explicar e ilustrar las propiedades de los diferentes algoritmos de aprendizaje  estadístico. \\

A continuación se hace una descripción de los algoritmos de aprendizaje estadístico que se estudian en este texto. Comenzaremos con el algoritmo k-vecinos más cercanos (k-NN) que introducimos  en el capítulo anterior, y  seguimos con los algoritmos lineales, hasta llegar a una breve introducción de redes neuronales. 


\section{Algoritmos de aprendizaje supervisados}
	

En esta sección  revisaremos los algoritmos de aprendizaje automátizado más populares,  explicaremos cómo
aprenden de los datos y cómo hacen predicciones. También discutiremos el concepto de complejidad del modelo que se presenta en cada uno, y se proporciona una visión general de cómo cada algoritmo construye el modelo. Examinaremos las fortalezas y las debilidades de cada algoritmo y a qué tipo de datos se pueden aplicar mejor. También se explica el significado de algunos los parámetros y opciones más importantes\footnote{Hablar de todos ellos está más allá del alcance del libro, y lo remitimos a la \href{https://scikit-learn.org/stable/index.html}{\color{vino} documentación} de {\ttfamily scikit-learn}
para más detalles.}.  Además, se    describen algunas  variantes de los algoritmos de  clasificación y regresión.\\



No es necesario leer las descripciones de cada algoritmo en detalle, pero
entender los modelos te dará una mejor idea de las diferentes formas que los algoritmos de aprendizaje automático pueden funcionar. Este capítulo también puede ser usado como guía de referencia, y puede volver a él cuando no esté seguro sobre el funcionamiento de cualquiera de los algoritmos.



\subsection{K-Vecinos más cercanos (k-NN\index{k-NN})} % k-Nearest Neighbors

El algoritmo k-NN es posiblemente el algoritmo de aprendizaje automático más simple. La construcción del modelo consiste sólo en almacenar el conjunto de datos de entrenamiento, luego  para hacer una predicción sobre  un nuevo punto de datos, el algoritmo encuentra los puntos (datos) más cercanos en el conjunto  de entrenamiento al nuevo dato (sus "vecinos más cercanos").



\subsubsection{Clasificación del K-Vecino}


En su versión más simple, el algoritmo k-NN\index{k-vecinos más cercanos} solo considera exactamente un vecino más cercano, que es el punto de datos de entrenamiento más cercano al punto  que queremos hacer una predicción. Entonces la predicción es simplemente el resultado conocido para este punto de entrenamiento.  En la figura \ref{1-NN}   se  ilustra este proceso, para el caso de la clasificación en la dataset de falsificación de datos (forge dataset):\\

\begin{lstlisting}
In[10]:  mglearn.plots.plot_knn_classification(n_neighbors=1)

\end{lstlisting}


Aqui se añaden  tres  puntos nuevos que se muestra como estrella en la figura \ref{1-NN}.{}   Se ejecuta la técnica sobre data de entrenamiento con un vecino. La predicción del algoritmo con un vecino más cercano  (one-nearest-neighbor) es etiquetada por el color sobre los puntos estrella, como en la figura \ref{1-NN}. \\



\begin{figure}[h!]
	\centering
	\includegraphics[scale=0.50]{../Imagen_Machine/1-NN}
	\caption{La predicciones se hacen con el modelo de 1-NN sobre la data forge.}
	\label{1-NN}
\end{figure} 


También podemos considerar un número $k$  arbitrario vecinos, en lugar de considerar sólo un vecino más cercano,  de aquí  el nombre del algoritmo de k-vecinos más cercanos.
Cuando se considera más de un vecino, la técnica usa la frecuencia o votación para asignar la
etiqueta.  Esto significa que para cada punto de prueba, se cuenta cuántos vecinos pertenecen
a la clase 0 y cuántos vecinos pertenecen a la clase 1. Luego le asigna la clase con mayor frecuencia, en otras palabras, la clase mayoritaria entre los k-vecinos más cercanos. \\

El siguiente ejemplo (Figura \ref{3-NN}) emplea la técnica con 3-vecinos más cercanos:\\

\begin{lstlisting}
In[11]:  mglearn.plots.plot_knn_classification(n_neighbors=3)
\end{lstlisting}

\begin{figure}[h!]
	\centering
	\includegraphics[scale=0.50]{../Imagen_Machine/3-NN}
	\caption{Predicciones hechas por el modelo 3-NN en la data falsificación de datos (forge dataset)}
	\label{3-NN}
\end{figure} 


Nuevamente, la predicción se muestra con el color de la estrella. Puedes ver que la predicción para el  punto nuevo  en la parte superior izquierda no es la misma que resultó  cuando
se  usó  un único vecino (1-NN).\\


Esta ilustración es corresponde al problema de clasificación binaria, pero el método  puede ser  
aplicado a conjuntos de datos con cualquier número de clases. En el caso de más de dos clases, es semejante, se cuentan  cuántos los vecinos pertenecen a cada clase y nuevamente predicen la clase más frecuente.\\



Ahora vemos cómo se aplicar el algoritmo k-vecinos más cercanos usando  {\ttfamily scikit-learn}. 
Primero, dividimos los datos en el conjunto de entrenamiento y el conjunto de prueba, necesario para  evaluar el rendimiento del modelo, como se discutió en el Capítulo \ref{Cap1}:\\


\begin{lstlisting}
In[12]:  from sklearn.model_selection import train_test_split
         X, y = mglearn.datasets.make_forge()
         
         X_train, X_test, y_train, y_test = train_test_split(     
                 X, y, random_state=0)
         
\end{lstlisting}


Seguidamente importamos e instanciamos la clase. Esto es cuando se establecen los parámetros, como el número de vecinos a usar, que en este caso es  3.\\

\begin{lstlisting}
In[13]:  from sklearn.neighbors import KNeighborsClassifier
         clf = KNeighborsClassifier(n_neighbors=3)
         
\end{lstlisting}


Ahora, ajustamos el clasificador sobre el conjunto de entrenamiento. Para el {\ttfamily KNeighboursClassifier}, significa almacenar el conjunto de datos, para  calcular los vecinos durante la predicción:\\

\begin{lstlisting}
In[14]:  clf.fit(X_train, y_train)

\end{lstlisting}

Para hacer la predición sobre la data de prueba, invocamos el método predictor. Para cada punto en el conjunto de prueba, este calcula su vecino más cercano en la data de entrenamiento, y calcula la clase más común entre estos. \\


\begin{lstlisting}
In[15]:  print("Predicciones en el conjunto de prueba: {}".format(
                 clf.predict(X_test)))

Out[15]: 

Predicciones en el conjunto de prueba: [1 0 1 0 1 0 0]

\end{lstlisting}


Para evaluar qué tan bien generaliza el modelo, podemos usar el método { \ttfamily {score }} aplicado a la data de prueba junto con sus etiquetas:\\

\begin{lstlisting}
In[16]:  print("Exactitud en el conjunto de prueba: {:.2f}".format(
                 clf.score(X_test, y_test)))

Out[16]:

    Exactitud en el conjunto de prueba: 0.86
\end{lstlisting}

Este resultado indica que el modelo es proximadamente 86\% exacto, esto significa que el modelo predice la clase correctamente para un  86\%  de las muestras en el conjunto de prueba.\\

 
Si bien esta ilustración es del problema de clasificación binaria, este método puede ser
aplicado a conjuntos de datos con cualquier número de clases. En el caso de  más clases, se cuenta  cuántos
los vecinos pertenecen a cada clase y nuevamente predicen la clase más común.\\


\subsubsection{Analizando KNeighboursClassifier}

Para conjuntos de datos bidimensionales, también podemos ilustrar la predicción para todos los  posibles puntos de prueba en el plano $XY$. Coloreando el área del plano de acuerdo con la clase que sería
asignado a un punto en esta región. Esto nos permite ver el límite de decisión donde el algoritmo asigna las clases, el cual es la división o límite entre las clases, en este caso   la clase 0  y  la clase 1.\\


El siguiente código produce la visualización del límite de decisión para uno, tres y nueve vecinos como se ve en la figura \ref{limites}.\\

\begin{lstlisting}
In[17]:  fig, axes = plt.subplots(1, 3, figsize=(10, 3))

         for n_neighbors, ax in zip([1, 3, 9], axes):
             # the fit method returns the object self, so 
             # we can instantiate  and fit in one line
             clf = KNeighborsClassifier(n_neighbors=n_neighbors).fit(X, y)
             mglearn.plots.plot_2d_separator(clf, X, fill=True, eps=0.5, 
                                              ax=ax, alpha=.4)
             mglearn.discrete_scatter(X[:, 0], X[:, 1], y, ax=ax)
             ax.set_title("{} neighbor(s)".format(n_neighbors))
             ax.set_xlabel("feature 0")
             ax.set_ylabel("feature 1")
         axes[0].legend(loc=3)

\end{lstlisting}

\begin{figure}[h!]
	\centering
	\includegraphics[scale=0.70]{../Imagen_Machine/limites}
	\caption{Límites de decisión creados por el modelo de vecinos más cercanos para diferentes valores
		de k-vecinos}
	\label{limites}
\end{figure} 

Como se ve a la izquierda en la figura \ref{limites}, usando un sólo vecino resulta  un límite de decisión
que sigue muy cercono  los datos de entrenamiento.  Considerando cada vez más vecinos conduce a un límite de decisión más suave. Un límite más suave corresponde a un modelo simple. En otras palabras, usar pocos vecinos corresponde a un modelo de alta complejidad (como se muestra en el lado derecho de la Figura \ref{complejidad}), y utilizando  muchos más vecinos responde a un modelo de baja complejidad (como se muestra en el lado izquierdo de la Figura \ref{complejidad}).  Si se considera el caso extremo donde el número de vecinos es igual al número datos del conjunto de entrenamiento, cada punto del conjunto de prueba tendría exactamente los mismos vecinos (todos puntos de entrenamiento) y todas las predicciones serían las mismas: la clase más frecuente en el conjunto de entrenamiento.\\



Para analizar y visualizar la relación entre la complejidad del modelo y la generalización discutida anteriormente, desarrollamos una práctica sobre la data  real Cáncer de Mamas. Comencemos dividiendo la data en los dos conjuntos, el de entrenamiento y el de prueba. Luego evaluamos el rendimiento del conjunto  entrenamiento y  pruebas con diferentes números de vecinos. El resultado se muestra en la figura \ref{constraste_E_P}:\\

\begin{lstlisting}
In[18]:  from sklearn.datasets import load_breast_cancer

         cancer = load_breast_cancer()
         X_train, X_test, y_train, y_test = train_test_split(
             cancer.data, cancer.target, stratify=cancer.target,
             random_state=66)
         
         training_accuracy = []
         test_accuracy = []
         # try n_neighbors from 1 to 10
         neighbors_settings = range(1, 11)
         
         for n_neighbors in neighbors_settings:
         # build the model
             clf = KNeighborsClassifier(n_neighbors=n_neighbors)
             clf.fit(X_train, y_train)
             # record training set accuracy
             training_accuracy.append(clf.score(X_train, y_train))
             # record generalization accuracy
             test_accuracy.append(clf.score(X_test, y_test))
             
         
         plt.plot(neighbors_settings, training_accuracy,
                  label="training accuracy")
         plt.plot(neighbors_settings, test_accuracy, 
                  label="test accuracy")
         plt.ylabel("Accuracy")
         plt.xlabel("n_neighbors")
         plt.legend()
\end{lstlisting}

\begin{figure}[h!]
	\centering
	\includegraphics[scale=0.50]{../Imagen_Machine/constraste_E_P}
	\caption{Comparación de la exactitud en data entrenamiento y prueba con una función { \ttfamily   	 {}n\_neighbours}  }
	\label{constraste_E_P}
\end{figure} 


El gráfico muestra la exactitud en el conjunto de entrenamiento y  prueba en el eje $y$ contra la configuración del k-vecino sobre el eje $x$.  Los graficos del mundo real rara vez son suaves o muy suaves, sin embargo, podemos reconocer algunas de las características de sobreajuste y falta de ajuste (recuerde  que considerando pocos vecinos el modelo se hace más complejo). Considerando un único vecino más cercano, la predicción en el conjunto de entrenamiento es perfecta. Pero cuando  se consideran más vecinos, el modelo se vuelve más simple y la exactitud del entrenamiento disminuye. La exactitud del conjunto de prueba para  un solo vecino es menor que cuando se usan más vecinos, lo que indica que usar un solo vecino más cercano conduce a un modelo que es demasiado complejo. Por otro lado, cuando se consideran 10 vecinos, el modelo es demasiado simple, y el rendimiento es aún peor.   

El mejor rendimiento está en algún lugar intermedio, como 
alrededor de seis vecinos. Aún así, es bueno tener en cuenta la escala del gráfico. El peor rendimiento es alrededor del 88\%  de exactitud, lo que aún podría ser aceptable.\\


\subsubsection{Regresión con k-vecinos}

También hay una variante de regresión del algoritmo del  k-vecinos\index{regresión k-vecinos más cercano} más cercanos. Para ilustrar esta técnica se usará el conjunto de datos de Onda (wave). La técnica  se empieza a implementar  con un sólo vecino, y  se agregaron tres puntos de prueba representados como estrellas verdes en el eje $x$. La predicción usando un solo vecino que es el valor objetivo del vecino más cercano. Estos se muestran como estrellas azules en la Figura \ref{regresionPresicion}:\\

\begin{lstlisting}
In[19]:  mglearn.plots.plot_knn_regression(n_neighbors=1)

\end{lstlisting}

\begin{figure}[h!]
	\centering
	\includegraphics[scale=0.50]{../Imagen_Machine/regresionPresicion}
	\caption{Predicciones hechas por la regresión del vecino más cercano en el conjunto de datos de onda }
	\label{regresionPresicion}
\end{figure} 


Nuevamente, con esta técnica (regresión) también podemos usar más de un vecino, como el  caso de clasificación. Cuando se usa
múltiples vecinos, la predicción es el promedio o la media de los vecinos relevantes (Figura 2-9):\\

\begin{lstlisting}
In[20]:  mglearn.plots.plot_knn_regression(n_neighbors=3)

\end{lstlisting}

\begin{figure}[h!]
	\centering
	\includegraphics[scale=0.50]{../Imagen_Machine/Predicionregresion3}
	\caption{Predicciones hechas por la regresión de 3 vecinos en el conjunto de datos de onda}
	\label{Predicionregresion3}
\end{figure} 



El algoritmo del k-vecinos más cercanos para la regresión se implementa en la clase del regresor del  K-Neighbors  en {\ttfamily scikit-learn}. Se usa de manera similar a {\ttfamily KNeighboursClassifier}:\\


\begin{lstlisting}
In[21]:  from sklearn.neighbors import KNeighborsRegressor

         X, y = mglearn.datasets.make_wave(n_samples=40)
         
         # split the wave dataset into a training and a test set
         X_train, X_test, y_train, y_test = train_test_split(
                 X, y, random_state=0)

         # instantiate the model and set the number of neighbors 
         #  to consider to 3
         reg = KNeighborsRegressor(n_neighbors=3)
         # fit the model using the training data and training targets
         reg.fit(X_train, y_train)
\end{lstlisting}

\noindent y ahora  hacemos la predicción sobre el conjunto de prueba:\\ 

 
\begin{lstlisting}
In[22]:  print("Prediciones en el Conjuto de prueba:\n{}".format(
                 reg.predict(X_test)))

Out[22]:

   "Prediciones en el Conjuto de prueba:
   [-0.054 0.357 1.137 -1.894 -1.139 -1.631 0.357 0.912 -0.447 -1.139]
\end{lstlisting}


También podemos evaluar el modelo utilizando el método de score o puntuación, en el cual los regresores retornan  el puntaje de  $R^{2}$, también conocido como coeficiente de determinación, esta es una medidad
de la bondad de  predicción para un modelo de regresión, y da una puntuación entre 0
y 1. Un valor de 1 corresponde a una predicción perfecta, y un valor de 0 corresponde
a un modelo constante que solo predice la media de las respuestas del conjunto de entrenamiento {\ttfamily{ y\_train}}:\\


\begin{lstlisting}
In[23]:  print("R^2 en conjunto de Prueba: {:.2f}".format(
                 reg.score(X_test, y_test)))

Out[23]:
    R^2 en conjunto de Prueba: 0.83
   
\end{lstlisting}

Aquí, la puntuación es de 0,83, lo que indica un ajuste del modelo relativamente bueno.



\subsubsection{Analizando el KNeighborsRegressor}


Para la data unidimensional, podemos visualizar cómo se ven las predicciones para todos los valores 
posibles  de las características (Figura \ref{PredicionregresionMAS3}). Para hacer esto, creamos un conjunto  de prueba que consiste en muchos puntos sobre la línea:\\


\begin{lstlisting}
In[24]: fig, axes = plt.subplots(1, 3, figsize=(15, 4))
        # create 1,000 data points, evenly spaced between -3 and 3
        line = np.linspace(-3, 3, 1000).reshape(-1, 1)
        for n_neighbors, ax in zip([1, 3, 9], axes):
            # make predictions using 1, 3, or 9 neighbors
            reg = KNeighborsRegressor(n_neighbors=n_neighbors)
            reg.fit(X_train, y_train)
            ax.plot(line, reg.predict(line))
            ax.plot(X_train, y_train, '^', c=mglearn.cm2(0), markersize=8)
            ax.plot(X_test, y_test, 'v', c=mglearn.cm2(1), markersize=8)
			 
            ax.set_title(
               "{}neighbor(s)\n train score:{:.2f} test score:{:.2f}".format(
                  n_neighbors, reg.score(X_train, y_train),
                  reg.score(X_test, y_test)))
            ax.set_xlabel("Feature")
            ax.set_ylabel("Target")
        axes[0].legend(["Model predictions", "Training data/target",
                        "Test data/target"], loc="best")

\end{lstlisting}

\begin{figure}[h!]
	\centering
	\includegraphics[scale=0.50]{../Imagen_Machine/PredicionregresionMAS3}
	\caption{ Comparando predicciones hechas con regresión del vecino más cercano para diferentes 
		valores de $k.$}
	\label{PredicionregresionMAS3}
\end{figure} 


Como podemos ver en el gráfico, usando un sólo vecino,  cada punto en el conjunto de entrenamiento tiene una influencia obvia en las predicciones,  y los valores predichos pasan por todos los puntos de datos.   Esto lleva a una predicción muy inestable.   Considerando más  vecinos conduce a predicciones más suaves,    pero estos no  se ajustan tan bien  a los datos de entrenamiento.


\subsubsection{Fortalezas, debilidades y parámetros}


En principio,  el clasificador  {\ttfamily KNeighbours} tiene  dos parámetros importantes: el
número de vecinos y la distancia (métrica) entre puntos de la data. En la práctica,
se usa un  número pequeño de vecinos, como tres o cinco que generalmete funciona bien, pero debe
ajustarse este parámetro. Seleccionar la métrica de distancia correcta, por ahora es algo que va 
más allá del alcance de este libro. Por defecto, se usa la distancia euclidiana, que funciona
bien en muchos entornos.\\


Una de las fortalezas de k-NN es que el modelo es muy fácil de entender, y a menudo
ofrece un rendimiento razonable sin muchos ajustes. Este algoritmo es un buen método de referencia para probar antes de considerar técnicas más avanzadas.\\

 
La construccón del modelo del vecinos más cercano suele ser muy rápido, pero cuando el conjunto de entrenamiento es muy grande (ya sea en número de características o en número de muestras)  la predicción puede ser lenta.  Cuando se utiliza el algoritmo k-NN, es importante preprocesar sus datos (consulte el Capítulo \ref{Cap3}).  Este enfoque generalmente no funciona bien sobre data con muchas características (cientos o más), y lo hace particularmente mal en conjuntos de datos donde la mayoría de las características son cero  en su mayor parte de las veces (los llamados conjuntos de datos dispersos\index{conjuntos dispersos} o sparse\index{sparse}).\\



Así,  mientras  el algoritmo del k-vecinos más cercano es fácil de entender, en la práctica no es muy frecuente, debido a que la predicción es lenta y su incapacidad para manejar muchas características.
El método que discutimos a continuación supera estos inconvenientes.


\subsection{Modelos lineales}


Los modelos lineales\index{modelos lineales} son una clase de modelos usados ampliamente en la práctica y han sido extensamente estudiados  en las últimas décadas, sus raíces que se remontan a más de cien
años. Los modelos lineales hacen una predicción usando una función lineal, donde la variable de entrada son las características. 

A continuación explicamos en forma breve esta ecuación matemáticas, y que tal vez te sientas familiarizado con la misma,  pues es  la ecuación  matemática básica de una línea recta. 


\subsubsection{Modelos lineales para regresión}


Para la regresión, la fórmula general de predicción  para un modelo lineal tiene el siguiente ecuación:
\[ \hat{y}  = w [0] * x[0] + w[1] * x[1] + ... + w[p] * x[p] + b \]
Aquí,  los entes $ x[0], ...,x[p]$ denota las características (en este ejemplo, el número de características es $p+1$), $w$ y $b$  son los parámetros del modelo que  aprenden, y $\hat{y}$  es 
la predicción que hace el modelo.\\

 Para un conjunto de datos con una sola característica, su ecuación se reduce a siguiente expresión:\\
\[ \hat{y} = w [0] * x [0] + b \]

\noindent  Aquí, $w [0]$  es la pendiente y $b$ es la ordenada en el origen (intercepto o corte con $y$). Para más características, $w$ contiene las pendientes a lo largo de cada eje de características. Alternativamente, puede pensar en la respuesta pronosticada como una suma ponderada de las características de entrada, con ponderaciones (que pueden ser negativos o positivos) dadas por las entradas de $w$.\\

Con  la siguiente instrucción aplicado sobre la data Onda  (conjunto de datos unidimensional), los parámetros $ w[0] $  y  $ b $  aprenden  y  conducen a la siguiente línea recta (ver Figura \ref{LineaR1}):\\

\begin{lstlisting}
In[25]:  mglearn.plots.plot_linear_regression_wave()

Out[25]:

    w[0]: 0.393906 b: -0.031804

\end{lstlisting}

\begin{figure}[h!]
	\centering
	\includegraphics[scale=0.50]{../Imagen_Machine/LineaR1}
	\caption{Predicciones de un modelo lineal en el conjunto de datos de Onda}
	\label{LineaR1}
\end{figure} 



La pendiente  $ w[0],$ está al alrededor de 0.4, lo cual se puede confirmar 
visualizando el gráfico. La intersección es donde la línea de predicción corta el eje $y,$
está ligeramente por debajo de cero, lo que también se puede confirma en gráfico.\\

Los modelos de regresión lineal pueden caracterizarse como modelos de regresión para los cuales
la predicción es una línea cuando se trata de una sola característica, un plano  para dos características o un hiperplano para más de dos características  o alta dimensión (es decir, cuando se usan más características).\\

Si se compara las predicciones hechas por regresión lineal, con la técnica del  {\ttfamily KNeighboursRegressor} como se observa en  la figura \ref{PredicionregresionMAS3}, con la línea recta las predicciones son muy restringidas, y al parecer se pierden muchos de los detalles importantes de la data.  En cierto sentido, esto es cierto. Es una suposición fuerte (y poco  realista) que nuestro objetivo $y$,   es combinación lineal  de las características.  Para  datos con muchas características, los modelos lineales pueden ser muy potentes. En particular, si tiene más características que  datos de entrenamiento, cualquier objetivo $y$ puede ser perfectamente modelado como una función lineal (sobre el conjunto de entrenamiento). \\




Existen muchos modelos lineales diferentes para la regresión. La diferencia entre estos modelos radica en cómo  aprenden los parámetros $ w$  y  $b $   del modelo  en la data de entrenamiento, y cómo controlar la complejidad del modelo.

A continuación   revisaremos los  Modelos lineales más comunes  para la regresión.



 

\subsubsection{Regresión lineal}
% (también conocida como mínimos cuadrados ordinarios)

La regresión lineal o mínimos cuadrados ordinarios es el método más simple de regresión. La regresión lineal\index{regresión lineal} busca los parámetros $w$ y $b$ de tal manera que minimicen el error cuadrático medio  entre las predicciones $\hat{y},$  y los valores reales de la  regresión  $y,$ en el conjunto de entrenamiento. El error cuadrático medio es la suma de las diferencias al cuadrado entre las prediciones y los verdaderos valores. La regresión lineal no tiene parámetros a controlar, lo cual es un beneficio, pero pierde opción  de controlar la complejidad del modelo.\\

 
 
El código que se prensenta  a continuación  produce el modelo que se muestra en la Figura \ref{LineaR1}:\\

\begin{lstlisting}
In[26]:  from sklearn.linear_model import LinearRegression
         X, y = mglearn.datasets.make_wave(n_samples=60)
         X_train, X_test, y_train, y_test = train_test_split(
              X, y, random_state=42)
         
         lr = LinearRegression().fit(X_train, y_train)          

\end{lstlisting}

Los parámetros $w$ pendientes,  también  llamado ponderación o coeficiente, son almacenados en el atributo {\ttfamily lr.coef\_},  mientras que el intercepto $b$, es almacenado el atributo {\ttfamily intercept\_ }:\\


\begin{lstlisting}
In[27]:  print("lr.coef_: {}".format(lr.coef_))
         print("lr.intercept_: {}".format(lr.intercept_))

Out[27]:

    lr.coef_: [ 0.394]
    lr.intercept_: -0.031804343026759746

\end{lstlisting}
 
 
 
   

 
 \begin{minipage}[t]{0.2\textwidth}
 	\centering\raisebox{\dimexpr \topskip-\height}{%
 		
\includegraphics[width=\textwidth]{Imagen_Machine/Pajaro.png}}
 \end{minipage}%\hfill
 \begin{minipage}[t]{0.75\textwidth} 
 	{\scriptsize	
 	Observe el subrayado  final  de aspecto extraño (piso bajo después de coef e intercept) 		
 	al final  de {\ttfamily coef\_}  e  {\ttfamily intercepción\_}. El algoritmo {\ttfamily scikit-learn} siempre almacena cualquier   cosa que se  deriva de los datos de entrenamiento en atributos que terminan con un  subrayado final. Eso es para identificarlos de los parámetros que son establecidos por el usuario.}
 \end{minipage}\\


 El atributo intercepto es siempre un número float, mientras que el atributo {\ttfamily coef\_} es un arreglo NumPy con una entrada por característica. En el caso de la data Onda (wave) solo se tiene  una característica, así  el  coeficiente {\ttfamily lr.coef\_}   tiene  un sólo valor. \\


 Consideremos el rendimiento en el conjunto de entrenamiento y el de prueba: \\

\begin{lstlisting}
In[28]:  print("Training set score: {:.2f}".format(lr.score(
                  X_train, y_train)))
         print("Test set score: {:.2f}".format(lr.score(X_test, y_test)))

Out[28]:

   Training set score: 0.67
   Test set score: 0.66
   
\end{lstlisting}



Un $R^{2}$ cerca de 0.66 no es muy bueno, pero podemos ver que los puntajes de la data
entrenamiento y  prueba están muy cerca. Esto significa probablemente que tienen un subajuste, 
y no sobreajuste.  Para esta data unidimensional, existe poco peligro de sobreajuste, pues el modelo es muy simple (o restringido). Sin embargo, para data con alta dimensionalidad (data con gran número de características), el modelo de regresión es más potente, y existe alta probabilidad de sobreajuste. Veamos cómo se ejecuta  {\ttfamily LinearRegression} sobre una data más compleja, como la data  Boston Housing.
Recuerde que esta data  tiene 506 muestra y 105 características. Primero carguemos la data, y dividamos la data en dos, en data de entrenamiento y prueba. Luego construimos el modelo como antes:\\


\begin{lstlisting}
In[29]:  X, y = mglearn.datasets.load_extended_boston()
         X_train, X_test, y_train, y_test = train_test_split(
                 X, y, random_state=0)
         lr = LinearRegression().fit(X_train, y_train)
    
\end{lstlisting}     
         
Al comparar los puntajes del $R^{2}$ de la data entrenamiento con la de prueba, encontramos que la predicción en la data de entrenamiento es muy exacta (95\%), pero  en la data de prueba es aún más bajo (61\%):\\


\begin{lstlisting}
In[30]:  print("Training set score: {:.2f}".format(lr.score(
                 X_train, y_train)))
         print("Test set score: {:.2f}".format(lr.score(X_test, y_test)))

Out[30]:
 
    Training set score: 0.95
    Test set score: 0.61

\end{lstlisting}

Esta discrepancia de rendimiento en el conjunto de entrenamiento y  de prueba es una clara señal de sobreajuste; debido a esto debemos tratar de encontrar un modelo que nos permita controlar la complejidad. 


\subsubsection{Regression Ridge}


La regressión  Ridge\index{regressión  Ridge} es también un modelo de regresión,   por lo que la fórmula que  utiliza para las predicciones es misma que utilizan los modelos lineales bajo los mínimos cuadrados ordinarios. 
Aunque, en la regresión Ridge, los coeficientes $w$ están escogidos no sólo con el fin de que predigan bien en los datos de entrenamiento, sino  también para ajustar una restricción adicional; se trata que las magnitudes de los coeficientes sean lo más pequeño posible, en otras palabras, que todas las entradas de $w $ deberían estar próximas al cero.   Intuitivamente, esto quiere decir que cada característica debe tener poco efecto sobre el resultado en lo posible (esto significa que la pendiente debe ser pequeña), mientras  que el modelo siga prediciendo bien.  Esta restricción es un ejemplo de lo que es llamado regularización. La regularización\index{regularización} significa explícitamente restringir un modelo para evitar sobreajuste. El tipo  particular de regulación usada  por la regresión Ridge\index{regularización  $L2$} es conocida como Regularización  $L2$. \\

 La regresión Ridge está disponible  en  {\ttfamily scikit-learn}  como el modelo lineal  {\ttfamily linear\_model.Ridge}.   Veamos una aplicación de la regresión Ridge  sobre la data de Boston Housing y su rendimiento:\\


\begin{lstlisting}
In[31]:  from sklearn.linear_model import Ridge

         ridge = Ridge().fit(X_train, y_train)
         print("Training set score: {:.2f}".format(ridge.score(
                X_train, y_train)))
         print("Test set score: {:.2f}".format(ridge.score(
                X_test, y_test)))

Out[31]:
    
    Training set score: 0.89
    Test set score: 0.75

\end{lstlisting}
	
 
 Como se puede ver, el puntaje de la data de entrenamiento 0.89 es menor que la del modelo de {\ttfamily LinearRegression}  que resultó 0.95; mientras que el puntaje de la data de prueba es 0.75, resultado que supera el puntaje  del modelo {\ttfamily LinearRegression} con 0.61. Está situación es consistente con lo esperado; en el caso de regresión lineal con sobreajuste. Ridge es un modelo más restringido, así es que tenemos menos probabilidad de sobreajuste en el modelo.  Un modelo menos complejo quiere decir, menor rendimiento  sobre el conjunto de entrenamiento, pero mejora generalización. Como estamos  interesados sólo  en el rendimiento de generalización, deberíamos escoger el modelo Ridge en lugar del modelo regression lineal.\\
 

 El modelo  Ridge  hace una compensación (trade-off) entre la simplicidad del modelo (coeficientes cerca de cero) y su rendimiento sobre la data de entrenamiento. La mayor importancia del modelo es que coloca la simplicidad versus el rendimiento en la data entrenamiento, que puede ser especificada por el usuario, mediante el parámetro alfa. En el ejemplo anterior, se usó el parámetro {\ttfamily{alfa}} por defecto, esto es   {\ttfamily{alfa = 1}}. No hay razón para pensar qué esto siempre nos dará la mejor compensación; sin embargo, el ajuste óptimo del  {\ttfamily{alfa}} depende particularmente de la data que se use.  El incremento del {\ttfamily{alfa}} obliga a los coeficientes moverse más hacia el cero, lo cual decrese el rendimiento en la data de entrenamiento,  pero podría ayudar a la generalización. Por ejemplo:\\


\begin{lstlisting}
In[32]:  ridge10 = Ridge(alpha=10).fit(X_train, y_train)
         print("Training set score: {:.2f}".format(ridge10.score(
                X_train, y_train)))
         print("Test set score: {:.2f}".format(ridge10.score(
                X_test, y_test)))

Out[32]:
    
    Training set score: 0.79
    Test set score: 0.64

\end{lstlisting}

Disminuyendo el {\ttfamily{alfa}}, la  restricción es menor sobre los coeficientes, indicando que nos movemos a la derecha en de la Figura \ref{complejidad}. Para los valores muy pequeños de {\ttfamily{alfa}}, los coeficientes están restringidos del todo, y finalizamos con un modelo que se parece al modelo de regresión lineal (LinearRegression):\\


\begin{lstlisting}
In[33]:  ridge01 = Ridge(alpha=0.1).fit(X_train, y_train)
         print("Training set score: {:.2f}".format(
                 ridge01.score(X_train, y_train)))
         print("Test set score: {:.2f}".format(
                 ridge01.score(X_test, y_test)))

Out[33]:

    Training set score: 0.93
    Test set score: 0.77

\end{lstlisting}

En este caso el parámetro  {\ttfamily{alfa = 0.1}} parece trabajar bien. Podríamos probar disminuiyendo {\ttfamily{alfa}} un poco más para verificar si mejora la generalización. Por ahora, observe  cómo el parámetro {\ttfamily{alfa}} concuerda con la complejidad modelo, como se muestra en la Figura \ref{complejidad}. En el Capítulo \ref{Cap5} discutiremos el método para seleccionar los parámetros.\\ 


También podemos obtener una visión cualitativa  de cómo el parámetro {\ttfamily{alfa}} cambia el
modelo, observando el atributo {\ttfamily{coef\_}}  de modelos con diferentes valores de {\ttfamily{alfa}}. Un {\ttfamily{alfa}} alto significa un modelo más restringido; de esta forma, se espera que el  valor del {\ttfamily{coef\_}} tengan magnitud más pequeña para un valor alto de {\ttfamily{alfa}}, que para un valor bajo de {\ttfamily{alfa}}.  Esto se confirma en  el gráfico de la Figura \ref{alfa}:\\


\begin{lstlisting}
In[34]:  plt.plot(ridge.coef_, 's', label="Ridge alpha=1")
         plt.plot(ridge10.coef_, '^', label="Ridge alpha=10")
         plt.plot(ridge01.coef_, 'v', label="Ridge alpha=0.1")
         
         plt.plot(lr.coef_, 'o', label="LinearRegression")
         plt.xlabel("Coefficient index")
         plt.ylabel("Coefficient magnitude")
         plt.hlines(0, 0, len(lr.coef_))
         plt.ylim(-25, 25)
         plt.legend()

\end{lstlisting}



\begin{figure}[h!]
	\centering
	\includegraphics[scale=0.50]{../Imagen_Machine/alfa}
	\caption{Comparando las magnitudes de los coeficientes de la regresión de Ridge con los diferentes valores de alfa y la regresion lineal}
	\label{alfa}
\end{figure} 



Aquí, el eje $x$ enumera las entradas de {\ttfamily{coef\_:}}  $x = 0$  muestra el coeficiente asociado con la primera característica, $x = 1$ el coeficiente asociado con la segunda característica, y así sucesivamente hasta $x = 100$. El eje $y$ muestra los valores numéricos correspondientes a los
coeficientes. En este caso, la conclusión principal  es que para {\ttfamily{alfa = 10}}, los coeficientes se encuentran alrededor de $ {-}3$ y $3$. Los coeficientes del modelo Ridge con {\ttfamily{alfa = 1}} son un poco más grande. Y para  {\ttfamily{alfa = 0.1}} tienen una magnitud aún mayor, que el caso de la regresión lineal sin  regularización alguna (que sería alfa = 0), los coeficientes son tan grandes que están fuera del gráfico.\\


Otra forma de entender la influencia de la regularización es fijar un valor de alfa
y variar la cantidad de datos de entrenamiento disponibles. Con la data Boston Housing  ilustremos está situación; se ejecuta un proceso  sub-muestreo  en la data,  luego se evaluan las técnicas {\ttfamily LinearRegression} y {\ttfamily  Ridge(alpha = 1)}, incrementando los tamaño del conjunto de entrenamiento, el resultado se muestra en la figura \ref{RigiLineal}, los gráficos muestran el rendimiento del modelo en función del tamaño de datos, y se llaman curvas de aprendizaje:\\

\begin{lstlisting}
In[35]:  mglearn.plots.plot_ridge_n_samples()

\end{lstlisting}

\begin{figure}[h!]
	\centering
	\includegraphics[scale=0.50]{../Imagen_Machine/RigiLineal}
	\caption{Comparando las magnitudes de los coeficientes de la regresión de Ridge con los diferentes valores de {\ttfamily{alfa}} y la regresion lineal}
	\label{RigiLineal}
\end{figure} 


Como era de esperar, el puntaje (score)  en la data entrenamiento es más alto que en la data de prueba en todos  tamaños de muestra, tanto para Ridge como en regresión lineal.  Debido a que Ridge está regularizada, el puntaje de la data de entrenamiento con Ridge  es menor que el puntaje de la data de entrenamiento con la regresión lineal en todos los ámbitos. Sin embargo, el puntaje de la data de prueba en Ridge es mejor que la linenl, particularmente para subconjuntos  pequeños de datos. Para data con  menos de 400 puntos, la regresión lineal no es capaz de aprender nada.   Con más y  más datos  disponibles para el modelo, ambos modelos mejoran y la  regresión lineal se pone al día con la Ridgea al final (asintótico). La lección aquí es que con suficientes datos de entrenamiento, la regularización se vuelve menos importante; y con suficientes datos, la regresión Ridge  y la regresión lineal tendrán el mismo rendimiento (el hecho de que esto suceda aquí cuando se usa el conjunto de datos completo es solo por casualida). Otro aspecto interesante de la Figura \ref{RigiLineal} es la disminución del rendimiento de data de entrenamiento en la regresión lineal.  Si se agregan más datos,  disminuye la posibilidad de sobreajuste del modelo  o  se le hace más dificil al modelo memorizar los datos de entrenamiento.\\

   
\subsubsection{Lasso}

 Una alternativa a la técnica de regresión lineal de Ridge con regularización es Lasso\index{regresión Lasso}. Como sucede con la regresión Ridge,  Lasso usa también una restrinción sobre los coeficientes 
 llevandolos  cerca de cero, pero en una forma ligeramente diferente, la llamada regularización\index{regularización $L1$} $L1.$ La consecuencia de la regularización $L1$ es que cuando se usa la regresión Lasso, algunos coeficientes son exactamente cero. Esto significa que algunas características son  completamente  ignoradas por el modelo; y  puede  interpretarse como una forma de selección automático de las características.   Teniendo algunos coeficientes exactamente igual a cero el modelo se hace  más fácil de interpretar y además  puede revelar las características más importantes del modelo.\\
  
	Apliquemos la técnica de Lasso al conjunto de datos de Boston Housing:\\

\begin{lstlisting}

In[36]:  from sklearn.linear_model import Lasso

         lasso = Lasso().fit(X_train, y_train)
         print("Training set score: {:.2f}".format(
                lasso.score(X_train, y_train)))
         print("Test set score: {:.2f}".format(
                lasso.score(X_test, y_test)))
         print("Number of features used: {}".format(
                np.sum(lasso.coef_ != 0)))

Out[36]:

    Training set score: 0.29
    Test set score: 0.21
    Number of features used: 4
\end{lstlisting}

Como puede ver, el rendimiento de Lasso es muy malo, tanto  en el conjunto de entrenamiento como en el de prueba. Esto indica que esta subajustado, además solo usó 4 de las 105 características.
De manera similar a Ridge, Lasso también tiene un parámetro de regularización, {\ttfamily alfa}, que controla
 qué tan fuerte se llevan los coeficientes a cero.  En el ejemplo anterior, utilizamos el valor por defecto  {\ttfamily{alpha = 1.0}}. Para reducir la falta ajuste o subajuste,  disminuimos el  {\ttfamily alfa}  y probamos. Para disminuirlo  aumentamos el número máximo de iteraciones para ejecutarlo, configurando  el atributo  {\ttfamily max\_iter}:\\


\begin{lstlisting}
In[37]:  # Aumentando la configuracion predeterminada "max_iter", de lo
         # contrario, el modelo advertiria que debemos aumentar max_iter.
         lasso001 = Lasso(alpha=0.01, max_iter=100000).fit(
                        X_train, y_train)
         print("Training set score: {:.2f}".format(
                lasso001.score(X_train, y_train)))
         print("Test set score: {:.2f}".format(
                lasso001.score(X_test, y_test)))
         print("Number of features used: {}".format(
                np.sum(lasso001.coef_ != 0)))
Out[37]:

    Training set score: 0.90
    Test set score: 0.77
    Number of features used: 33

\end{lstlisting}



Un {\ttfamily alfa} bajo permite ajustar un modelo más complejo, el cual trabaja mejor sobre la data de entrenamiento  y prueba. El rendimiento es ligeramente mejor que la Ridge, y además usa $ 33$ características de las $105$. Esto hace el modelo más potente y fácil de entender.\\

Sin embargo, si configuramos un {\ttfamily alfa} demasiado bajo, nuevamente eliminamos el efecto de regularización y terminamos con sobreajuste, con un resultado similar a LinearRegression:\\

\begin{lstlisting}
In[38]:  lasso00001 = Lasso(alpha=0.0001, max_iter=100000).fit(
                          X_train, y_train)
         print("Training set score: {:.2f}".format(lasso00001.score(
                 X_train, y_train)))
         print("Test set score: {:.2f}".format(lasso00001.score(
                 X_test, y_test)))
         print("Number of features used: {}".format(
                np.sum(lasso00001.coef_ != 0)))

Out[38]:

    Training set score: 0.95
    Test set score: 0.64
    Number of features used: 94
\end{lstlisting}
           
   
 
 Nuevamente, podemos graficar los coeficientes de los diferentes modelos, similarmente como se obseva en la figura \ref{RigiLineal}. El resultado se muestra en la figura \ref{LasoRigi}.\\
 
 
\begin{lstlisting}
In[39]:  plt.plot(lasso.coef_, 's', label="Lasso alpha=1")
         plt.plot(lasso001.coef_, '^', label="Lasso alpha=0.01")
         plt.plot(lasso00001.coef_, 'v', label="Lasso alpha=0.0001")
         plt.plot(ridge01.coef_, 'o', label="Ridge alpha=0.1")
         plt.legend(ncol=2, loc=(0, 1.05))
         plt.ylim(-25, 25)
         plt.xlabel("Coefficient index")
         plt.ylabel("Coefficient magnitude")
\end{lstlisting}

\begin{figure}[h!]
	\centering
	\includegraphics[scale=0.50]{../Imagen_Machine/LasoRigi}
	\caption{Comparación de las magnitudes de los coeficientes para la regresión de Lazo con diferentes valores  alfa y regresión Ridge}
	\label{LasoRigi}
\end{figure} 


Para   {\ttfamily{alfa = 1}}, no solo vemos que la mayoría de los coeficientes son cero (que ya conocíamos), sino que que los coeficientes restantes también son pequeños en magnitud. Llevando el  {\ttfamily{alfa = 0.01}}, obtenemos la solución que se muestra como los puntos rojos, lo que causa que la mayoría de las características son exactamente cero. Usando alfa = 0.00001, obtenemos un modelo poco regularizado, con la mayoría de los coeficientes distintos de cero y de gran magnitud.  Por comparación, la mejor solución de Ridge se muestra en verde azulado. El modelo Ridge con  {\ttfamily{alfa = 0.1}} tiene similar rendimiento predictivo al modelo de Lazo con  {\ttfamily{alfa = 0.01}}, pero usando Ridge, todos los coeficientes  son distintos de cero.\\

En la práctica, la regresión de Ridge suele ser la primera opción entre estos dos modelos.
Sin embargo, si tiene una gran cantidad de características y espera que solo unas pocas sean
importante, Lasso podría ser una mejor opción. Del mismo modo, si desea tener un modelo que sea más fácil de interpretar, Lasso proporcionará un modelo que es más fácil para interpretar y entender, ya que seleccionará solo un subconjunto de las características de entrada.  

scikit-learn también proporciona la clase ElasticNet, que combina las penalizaciones de Lasso y Ridge.  En la práctica, esta combinación funciona mejor, aunque al costo de tener dos parámetros para ajustar: uno para la regularización $L1$ y otro para la regularización $L2$.
 
\subsubsection{Modelos lineales para clasificación binaria}\label{Mlineales}

Los modelos lineales también se utilizan ampliamente para la clasificación. Consideremos la clasificación  binaria\index{clasificación  binaria}  en primer lugar. En este caso,  una predicción  se hace  utilizando la siguiente fórmula:

\[ \hat{y} = w[0] * x[0] +w[1] * x[1] + ... + w[p] * x [p] + b > 0  \]

La fórmula se ve muy similar a la de la regresión lineal, pero en lugar de 
retornar solo la suma ponderada de las características, limitamos el valor predicho a cero.  Si la función es menor que cero, predecimos la clase –1; si es mayor que cero, 
predecimos la clase +1. Esta regla de predicción es común a todos los modelos lineales para clasificación\index{Modelos lineales para clasificación}.
 Una vez más, hay varias formas diferentes de encontrar los coeficientes $w$ y el intercepto $b$.\\


Para los modelos lineales por regresión la salida  $\hat{y}$, es una función lineal de las características: una línea, un plano o un hiperplano (para dimensiones superiores). Para modelos lineales por clasificación, el límite de decisión es una función lineal de la entrada. En otras palabras, un clasificador lineal (binario) es el clasificador  que separa dos clases usando una línea, un plano o un hiperplano. Se ilustrará este clasificador con ejemplos más adelante a fin de aclara la situación. \\

Existen muchos algoritmos de aprendizaje automático con modelos lineales. Todos estos algoritmos difieren en las siguientes dos formas:

\begin{itemize}
	\item La forma en que miden qué tan bien la  combinación particular de coeficientes 
	e intercepto, se ajusta a los datos de entrenamiento
	
	\item Y qué tipo de regularización emplea
\end{itemize}


Los algoritmos eligen diferentes formas para medir qué "tan bien se ajustan a la data de entrenamiento".   Por razones matemáticas no es posible ajustar los parámetros $w$ y $b$, para minimizar la cantidad de clasificaciones erróneas que producen los algoritmos. Para nuestros propósitos y muchas aplicaciones, las diferentes opciones para el primer item  de la lista anterior tiene poca importancia  o significancia, y se conocen como  funciones de pérdida.


Los dos algoritmos de clasificación lineal más comunes son la regresión logística\index{regresión logística},  implementada
en  {\ttfamily{linear\_model.LogisticRegression}} y  máquinas de soporte de vectores lineal\index{máquinas de soporte de vectores lineal} (SVM lineales), implementado en {\ttfamily{svm.LinearSVC}} (SVC significa máquina de vector de soporte para clasificación).  A pesar de su nombre, LogisticRegression  es un algoritmo de clasificación y no un
algoritmo de regresión, y no debe confundirse con LinearRegression.\\


A continuaciónn se aplica los modelos {\ttfamily LogisticRegression} y {\ttfamily LinearSVC} sobre la dataset de falsificación de datos (forge dataset), y se visualizar el límite de decisión como se presenta en los modelos lineales (Figura \ref{SVM}):


\begin{lstlisting}
In[40]:  from sklearn.linear_model import LogisticRegression
         from sklearn.svm import LinearSVC
         
         X, y = mglearn.datasets.make_forge()
         
         fig, axes = plt.subplots(1, 2, figsize=(10, 3))
         
         for model, ax in zip([LinearSVC(), LogisticRegression()], axes):
             clf = model.fit(X, y)
             mglearn.plots.plot_2d_separator(clf, X, fill=False, eps=0.5,
                                             ax=ax, alpha=.7)
             mglearn.discrete_scatter(X[:, 0], X[:, 1], y, ax=ax)
             ax.set_title("{}".format(clf.__class__.__name__))
             ax.set_xlabel("Feature 0")
             ax.set_ylabel("Feature 1")
         axes[0].legend()
\end{lstlisting}


\begin{figure}[h!]
	\centering
	\includegraphics[scale=0.65]{../Imagen_Machine/SVM}
	\caption{Límites de decisión de la regresión lineal SVM  y regresión logística en la data falsificación de datos (forge dataset) con los parámetros por defecto}
	\label{SVM}
\end{figure}


% Figura 2-15. Límites de decisión de la regresión lineal SVM  y regresión logística en el conjunto de datos  forge con los parámetros por defecto.\\


En la figura \ref{SVM}, la primera característica de la  data de falsificación (forge dataset)  está en el eje $x,$   y la  segunda característica en el eje $y$. Se observan los límites de decisión encontrados por
LinearSVC y LogisticRegression respectivamente como líneas rectas, separando el área  clasificada como clase 1 en la parte superior, y como clase 0 en la parte inferior. En  otras palabras, cualquier dato nuevo que se encuentre por encima de la línea negra se clasificará como clase 1 por el clasificador respectivo, mientras que cualquier punto que se encuentre debajo de la línea negra se clasificará  como clase 0.\\

Los dos modelos presentan límites de decisión similares. Observe que ambos modelos clasifican erróneamente dos  puntos. Por defecto, ambos modelos aplican una regularización $L2$, en la misma
forma en que Ridge lo hace para la regresión.\\

Para LogisticRegression y LinearSVC, el parámetro de compensación (trade-off) que determina 
la regulaión se llama $C$. Para valores  altos de $C$ corresponden a menor regularización.
En otras palabras, cuando se usa un valor alto para el parámetro $C$, LogisticRegression y LinearSVC intentan ajustar el conjunto de entrenamiento lo mejor posible, mientras que con valores bajos de $C$, los modelos ponen más énfasis en encontrar un coeficiente $w$ que esté cerca de cero.\\


Hay otro aspecto interesante de cómo actúa el parámetro $C$. Usando valores bajos de $C$  hará que los algoritmos intenten ajustarse a la "mayoría" de los puntos de datos, mientras usan  un valor más alto de $C$, enfatiza la importancia de cada punto  individual de los datos sea clasificado correctamente. \\

A continuación se ilustra con la técnica  LinearSVC (Figura \ref{LinearSVC}):\\

\begin{lstlisting}
In[41]:  mglearn.plots.plot_linear_svc_regularization()

\end{lstlisting}

\begin{figure}[h!]
	\centering
	\includegraphics[scale=0.65]{../Imagen_Machine/LinearSVC}
	\caption{Límites de decisión de un SVM lineal en la data falsificación de datos (forge dataset) para diferentes valores de C}
	\label{LinearSVC}
\end{figure}


 En el lado izquierdo de la figura \ref{LinearSVC}, tenemos una $C$ muy pequeña que corresponde a una gran regularización. La mayoría de los puntos en la clase 0 se observan en la parte superior, igualmente que la mayoría de los puntos de la clase 1 están en la parte el inferior. El modelo fuertemente regularizado selecciona una línea relativamente horizontal, clasificando erróneamente dos puntos. En el grafico central, C es ligeramente más alto, el modelo se enfoca más en las dos muestras mal clasificadas, inclinando el límite de decisión.  Finalmente, en el lado derecha, el valor muy alto de $C$, el modelo inclina aún más el límite de decisión, y ahora clasifica correctamente todos los puntos en la clase 0. Pero uno de los puntos en la clase 1 sigue siendo clasificado mal, ya que no es posible clasificar correctamente todos los puntos en este conjunto de datos usando una recta línea. El modelo ilustrado en el lado derecho se esfuerza por clasificar correctamente todos  puntos, pero podría no capturar bien el diseño general de las clases. En otras palabras,  es probable que este modelo esté sobreajustado.\\


De manera similar al caso de regresión, los modelos lineales para clasificación pueden parecer muy
restrictivo en espacios de baja dimensión, solo permitiendo límites de decisión que son
líneas rectas o planos.  Sin embargo,  en  altas dimensiones los modelos lineales para clasificación 
se vuelve muy poderoso, y la protección contra el sobreajuste se vuelve cada vez más importante al considerar más características.\\ 


Analicemos el algoritmo LinearLogistic con más cuidado con el conjunto de datos Cáncer de mama (Breast Cancer dataset):\\ 


\begin{lstlisting}
In[42]:  from sklearn.datasets import load_breast_cancer
         cancer = load_breast_cancer()
         X_train, X_test, y_train, y_test = train_test_split(cancer.data, 
              cancer.target, stratify=cancer.target, random_state=42)
         logreg = LogisticRegression().fit(X_train, y_train)
         print("Training set score: {:.3f}".format(
                logreg.score(X_train, y_train)))
         print("Test set score: {:.3f}".format(
                logreg.score(X_test, y_test)))
Out[42]:

    Training set score: 0.953
    Test set score: 0.958
\end{lstlisting}


El valor por defecto de $C$ es 1, y proporciona un buen rendimiento, con una exactitud del 95\% , tanto el conjunto de entrenamiento como en el de prueba. Pero como el rendimiento en el  conjunto entrenamiento y en el de pruebas son muy cercanos, es probable presente cierto  subajutes.  Aumentemos $C,$ para que se ajuste a una forma más flexible el modelo:\\


\begin{lstlisting}
In[43]:  logreg100 = LogisticRegression(C=100).fit(X_train, y_train)
         print("Training set score: {:.3f}".format(
                logreg100.score(X_train, y_train)))
         print("Test set score: {:.3f}".format(
                logreg100.score(X_test, y_test)))

Out[43]:

    Training set score: 0.972
    Test set score: 0.965
\end{lstlisting}

Usando un $C = 100,$ da como resultado una mayor exactitud del conjunto de entrenamiento y también un  ligero aumento en el de prueba,  confirmando nuestra intuición de que un modelo más complejo debería funcionar mejor. También podemos investigar qué sucede si utilizamos un modelo aún más regularizado que
el valor por defecto ($ C = 1$), configurando $C = 0.01$: \\

\begin{lstlisting}
In[44]:  logreg001 = LogisticRegression(C=0.01).fit(X_train, y_train)
         print("Training set score: {:.3f}".format(
                 logreg001.score(X_train, y_train)))
         print("Test set score: {:.3f}".format(
                 logreg001.score(X_test, y_test)))
Out[44]:

    Training set score: 0.934
    Test set score: 0.930
\end{lstlisting}

Como era de esperar, al moverse más a la izquierda a lo largo de la escala mostrada en la Figura \ref{complejidad}  resulta  un modelo  subajustado, tanto en el conjunto de entrenamiento y como en el de prueba, la exactitud disminuye en relación con los parámetros por defecto.\\


Finalmente, veamos los coeficientes aprendidos por los modelos con los tres conjuntos diferentes del parámetro de regularización $C$ (Figura \ref{Rlogistica}):\\

\begin{figure}[h!]
	\centering
	\includegraphics[scale=0.65]{../Imagen_Machine/Rlogistica}
	\caption{Coeficientes aprendidos por regresión logística en el conjunto de datos de Cáncer de Mama para
		diferentes valores de C}
	\label{Rlogistica}
\end{figure}

\begin{lstlisting}
In[45]:  plt.plot(logreg.coef_.T, 'o', label="C=1")
         plt.plot(logreg100.coef_.T, '^', label="C=100")
         plt.plot(logreg001.coef_.T, 'v', label="C=0.001")
         plt.xticks(range(
                cancer.data.shape[1]), cancer.feature_names, rotation=90)
         plt.hlines(0, 0, cancer.data.shape[1])
         plt.ylim(-5, 5)
         plt.xlabel("Coefficient index")
         plt.ylabel("Coefficient magnitude")
         plt.legend()

\end{lstlisting}

 


\begin{minipage}[t]{0.2\textwidth}
	\centering\raisebox{\dimexpr \topskip-\height}{%
		
\includegraphics[width=\textwidth]{Imagen_Machine/Escorpion.png}}
\end{minipage}%\hfill
\begin{minipage}[t]{0.75\textwidth} 
	{\scriptsize		
	Como la  {\ttfamily LogisticRegression} aplica una regularización $L2$ por defecto, el resultado parece similar  al producido por Ridge en la Figura \ref{RigiLineal}. Una regularización más fuerte obliga 
		a los coeficientes acercarse más hacia cero, aunque los coeficientes nunca llegan a ser 
		exactamente cero. Inspeccionando   el gráfico más de cerca, también podemos ver un efecto 
		interesante en el tercer  coeficiente,  para el   "{}perímetro medio"{}  \: (mean perimeter).   Para $C=100$ y $C=1$,
		el coeficiente es negativo, mientras que  para $C=0,001$, el coeficiente es positivo, con 
		una magnitud que es incluso mayor para  $C=1$. Interpretando un modelo como este, uno 
		podría  pensar que el coeficiente nos dice con qué clase se  asocia una característica. 
		Por ejemplo, uno podría pensar que un alto  valor en la característica "{}error de textura{}"{}   
		está relacionada con una muestra que es  "{}maligno."{}    \: Sin embargo, el cambio de signo en el 
		coeficiente de  "{}perímetro medio"{}  significa que dependiendo del modelo que miramos, un alto 
		{}"{}perímetro medio{}"{ }  podría ser tomado como  indicativo de "benigno"{ }  o indicativo de  "maligno". 	
	Esto ilustra que las interpretaciones de  coeficientes de modelos lineales siempre deben 	
	tomarse como gran cuidado.}
\end{minipage} \\


Si deseamos un modelo más interpretable, la regularización $L1$ podría ayudar, ya que límita al  modelo a utilizar sólo algunas características. \\

A continuación se presenta el gráfico de coeficientes y las precisiones de clasificación para tres regularizaciónes de $L1$,  0.001,  1  y  100  (Figure \ref{PenalRegreLogis}):\\




\begin{lstlisting}
In[46]:  for C, marker in zip([0.001, 1, 100], ['o', '^', 'v']):
             lr_l1 = LogisticRegression(C=C, penalty="l1").fit(
                                     X_train, y_train)
             print("Training accuracy of l1 logreg with C={:.3f}: {:.2f}"
                    .format(C, lr_l1.score(X_train, y_train)))
             print("Test accuracy of l1 logreg with C={:.3f}: {:.2f}".format(
                     C, lr_l1.score(X_test, y_test)))
             plt.plot(lr_l1.coef_.T, marker, label="C={:.3f}".format(C))
        
         plt.xticks(range(
                cancer.data.shape[1]), cancer.feature_names, rotation=90)
         plt.hlines(0, 0, cancer.data.shape[1])
         plt.xlabel("Coefficient index")
         plt.ylabel("Coefficient magnitude")
         
         plt.ylim(-5, 5)
         plt.legend(loc=3)

Out[46]:

	Training accuracy of l1 logreg with C=0.001: 0.91
	Test accuracy of l1 logreg with C=0.001: 0.92
	Training accuracy of l1 logreg with C=1.000: 0.96
	Test accuracy of l1 logreg with C=1.000: 0.96
	Training accuracy of l1 logreg with C=100.000: 0.99
	Test accuracy of l1 logreg with C=100.000: 0.98
\end{lstlisting}


Como pueden ver, hay muchos paralelismos entre los modelos lineales para la clasificación binaria y los modelos lineales para la regresión. Como en la regresión, la principal diferencia entre los modelos es el parámetro de penalización, que influye en la regularización y si el modelo utilizará todas las características disponibles o seleccionar sólo un subconjunto.\\


\begin{figure}[h!]
	\centering
	\includegraphics[scale=0.65]{../Imagen_Machine/PenalRegreLogis}
	\caption{Coeficientes aprendidos por regresión logística con penalización L1 en el Conjunto de datos sobre el cáncer para diferentes valores de C}
	\label{PenalRegreLogis}
\end{figure}




 \subsubsection{Modelos lineales para clasificación multiclases}
 %Linear models for multiclass classification

Muchos modelos de clasificación lineal son sólo para la clasificación binaria, y no se extienden naturalmente al caso multiclase (con la excepción de la regresión logística). Una técnica común para extender un algoritmo de clasificación binaria a un algoritmo de clasificación multiclase es el enfoque uno versus el resto (one-vs.-rest). En el enfoque one-vs.-rest, un modelo binario  es aprendido por cada clase que trata de separar esa clase de todas las otras clases, resultando  tantos modelos binarios como clases haya. Para hacer una predicción, todos los clasificadores binarios se ejecutan en un punto de prueba. El clasificador que tiene la puntuación (score) más alta en su única clase "gana", y esta  clase etiquetada  se devuelve como la predicción.\\

 

Tener un clasificador binario por cada clase, resulta también tener un vector de coeficientes $w$ y un intercepto $b$  para cada clase. La clase para la cual el resultado de la clasificación de la fórmula dada aquí, es la clase más alta a la que se le asigna la etiqueta:

\[  w [0] * x[0] + w[1] * x[1] + ... + w[p] * x[p] + b \]
 
 
Las matemáticas detrás de la regresión logística multiclase difieren un poco del enfoque uno versus el resto (one vs.-rest), pero también dan lugar a un vector de coeficiente e intercepto por clase, y se aplica el mismo método para predecir.\\

 
Ilustremos  el método one-vs.-rest para clasificar la data,  aplicado a un conjunto de datos  con tres clases. Es un conjunto de datos bidimensionales, donde cada clase viene dada por datos muestreados de una distribución gaussiana (ver Figura \ref{Toy} ): \\



\begin{lstlisting}
In[47]:  from sklearn.datasets import make_blobs

         X, y = make_blobs(random_state=42)
         mglearn.discrete_scatter(X[:, 0], X[:, 1], y)
         plt.xlabel("Feature 0")
         plt.ylabel("Feature 1")
         plt.legend(["Class 0", "Class 1", "Class 2"])
\end{lstlisting}



\begin{figure}[h!]
	\centering
	\includegraphics[scale=0.65]{../Imagen_Machine/Toy}
	\caption{Conjunto de datos bidimensional sobre juguetes que contiene tres clases}
	\label{Toy}
\end{figure}


Ahora, entrenemos  el classificador LinearSVC sobre la data:\\


\begin{lstlisting}
In[48]:  linear_svm = LinearSVC().fit(X, y)
         print("Coefficient shape: ", linear_svm.coef_.shape)
         print("Intercept shape: ", linear_svm.intercept_.shape)
Out[48]:

	Coefficient shape: (3, 2)
	Intercept shape: (3,)

\end{lstlisting}

Vemos que la forma del { \ttfamily coef\_} es (3, 2), lo que significa que cada fila de { \ttfamily coef\_}  contiene el vector de coeficientes para una de las tres clases y cada columna tiene el valor de coeficiente para una característica específica (hay dos en este conjunto de datos). El  { \ttfamily intercept\_}  es ahora una matriz unidimensional, almacenando los interceptos  para cada clase. \\

Visualizalicemos  las líneas de limites dadas por los tres clasificadores binarios (Figure \ref{3clases}):\\


\begin{lstlisting}
In[49]:  mglearn.discrete_scatter(X[:, 0], X[:, 1], y)
         line = np.linspace(-15, 15)
         for coef, intercept, color in zip(linear_svm.coef_, 
                                           linear_svm.intercept_,
                                           ['b', 'r', 'g']):
             plt.plot(line,
                      -(line * coef[0] + intercept) / coef[1], c=color)
         plt.ylim(-10, 15)
         plt.xlim(-10, 8)
         plt.xlabel("Feature 0")
         plt.ylabel("Feature 1")
         plt.legend(['Class 0', 'Class 1', 'Class 2','Line class 0', 
                     'Line class 1', 'Line class 2'], loc=(1.01, 0.3))
\end{lstlisting}


Se puede ver que todos los puntos pertenecientes a la clase 0 en los datos de entrenamiento están por encima de la línea correspondiente a la clase 0, lo que significa que están en el lado "{}clase 0{}"{}  de este clasificador binario. Los puntos de la clase 0 están por encima de la línea correspondiente a la clase 2, lo que significa que están clasificados como "{}resto" {}   por el clasificador binario de la clase 2.



\begin{figure}[h!]
	\centering
	\includegraphics[scale=0.60]{../Imagen_Machine/3clases}
	\caption{Límites de decisión aprendidos por los tres clasificadores uno contra uno.-rest}
	\label{3clases}
\end{figure}

 \noindent Los puntos pertenecientes a la clase 0 que  están a la izquierda de la línea correspondiente a la clase 1,  lo que significa el clasificador binario para la clase 1 también los clasifica como "resto".  Por lo tanto, cualquier punto en esta área será clasificada como clase 0 por el clasificador final (el resultado de la fórmula de clasificación de confianza  para el clasificador "{}clase 0"{} es mayor que cero, mientras que es menor que cero para las otras dos clases).\\ 
 

¿Pero qué pasa con el triángulo en el centro del gráfico?   Los tres clasificadores binarios clasifican los puntos allí como "resto." ¿A qué clase se le asignaría un punto? La respuesta es la que tiene el valor más alto para la fórmula de clasificación: la clase de la línea más cercana.\\


El siguiente ejemplo (Figura \ref{2D}) muestra la predicción  para todas las regiones del espacio 2D:\\


\begin{lstlisting}
In[50]:  mglearn.plots.plot_2d_classification(linear_svm, X, 
                                               fill=True, alpha=.7)
         mglearn.discrete_scatter(X[:, 0], X[:, 1], y)
         line = np.linspace(-15, 15)
         for coef, intercept, color in zip(linear_svm.coef_, 
                                           linear_svm.intercept_, 
                                           ['b', 'r', 'g']):
              plt.plot(line,
                       -(line * coef[0] + intercept) / coef[1], c=color)
         plt.legend(['Class 0', 'Class 1', 'Class 2', 'Line class 0',
                     'Line class 1','Line class 2'], loc=(1.01, 0.3))
         plt.xlabel("Feature 0")
         plt.ylabel("Feature 1")

\end{lstlisting}

\begin{figure}[h!]
	\centering
	\includegraphics[scale=0.60]{../Imagen_Machine/2D}
	\caption{Límites de decisión aprendidos por los tres clasificadores  one-vs.-rest}
	\label{2D}
\end{figure}



\subsubsection{Fortalezas, debilidades y parámetros}

 El principal parámetro de los modelos lineales es el parámetro de regularización, llamado {\ttfamily alfa}  en los modelos de regresión y $C$ en {\ttfamily LinearSVC} y {\ttfamily LogisticRegression}. Los valores grandes para alfa o pequeños para $C$ significan modelos simples. En particular para los modelos de regresión, ajustar estos es muy importante. Por lo general,  $C$ y {\ttfamily alfa}  se buscan en una escala logarítmica.   La otra decisión que debe tomar es si desea usar la regularización $L1$ o la regularización $L2$. Si asume que solo algunos de características son realmemte importantes, debe usar $L1$. De lo contrario, debe establecer por defecto a $L2$.   $L1$ también puede ser útil si la interpretabilidad del modelo es importante. Como $L1$ solo usará algunas características, es más fácil explicar qué características son importantes para el modelo, y cuáles son los efectos de estas características. 
 

 
Los modelos lineales son muy rápidos de entrenar y también rápidos para predecir. Escalan a conjuntos de datos muy grandes y funcionan bien con datos sparse. Si sus datos consisten en cientos de miles  o millones de muestras, es posible que desee investigar utilizando el solver =  "{}sag"{} opción en {\ttfamily LogisticRegression} y {\ttfamily Ridge}, que puede ser más rápido que el predeterminado en grandes conjuntos de datos.  Otras opciones son las clases {\ttfamily SGDClassifier} y  {\ttfamily SGDRegressor}, que implementa versiones aún más escalables que  los modelos lineales descritos aquí.\\

Otra fortaleza de los modelos lineales es que hacen que ellos son relativamente fácil de entender
cómo  hacen una predicción, mediante las fórmulas que vimos anteriormente para la regresión y la clasificación.  Pero si un conjunto de datos tiene características altamente correlacionadas; en estos casos, los coeficientes pueden ser difíciles de interpretar.\\

Generalmente los modelos lineales funcionan bien cuando el número de características es grande en comparación con el número de muestras.  También se utilizan a menudo en conjuntos de datos muy grandes, simplemente porque no es factible entrenar otros modelos.  Sin embargo, en espacios de menor dimensión, otros modelos podrían producir un mejor rendimiento de generalización. En  la página \pageref{kenrelpag} en la sección  "Máquinas de soporte de vectores kernelized"{},  trataremos algunos ejemplos en los que los modelos lineales fallan.







\vspace{-2cm}	
\hspace{-1cm}
{\fcolorbox{red}{white}{ 		
		\begin{minipage}[16cm]{13.5cm}    %   centrado, largo y ancho                {14cm} 
			\vspace{0.4cm}	
			
			
			\begin{center}
				{\bf{Método de encadenamiento}}
			\end{center}	
			
			
			El método de ajuste de todos los modelos en scikit-learn devuelve en el valor.  Lo que
			permite escribir código como el siguiente, que ya hemos utilizado extensamente en este capítulo:\\
			
			{ \ttfamily   	{\scriptsize{  	 
						
						In[51]:\\		
						{ \color{blue}
							
							\# instantiate model and fit it in one line \vspace{-0.2cm}
							
							logreg = LogisticRegression().fit(X\_train, y\_train) } \\ }}}
			
			
			Aquí, se usa  el valor  retornado del ajuste (del modelo) para luego asignar el modelo entrenado a la variable logreg. Esta concatenación del método (aquí \_\!\!\_init\_\!\!\_  y luego ajusta) se conoce como  método encadenamiento (method chaining). Otra aplicación común del encadenamiento del método en scikit-learn es encajar y predecir en línea:\\
			
			{\ttfamily   	{\scriptsize{  	 
						
						In[52]:\\		
						{ \color{blue}		
							
							logreg = LogisticRegression()  \vspace{-0.2cm}
							
							y\_pred = logreg.fit(X\_train, y\_train).predict(X\_test) } \\}}}
			
			
			
			Por último, incluso se puede hacer la instanciación del modelo, ajuste, y  predicción en línea:\\ 
			
			
			{\ttfamily   	{\scriptsize{  	 
						
						In[53]:\\		
						{ \color{blue}
							
							y\_pred = LogisticRegression().fit(X\_train, y\_train).predict(X\_test) } \\ }}}
			
			
			Aunque esta variante muy corta no es ideal. Están sucediendo muchas cosas en una sola línea, lo que podría dificultar la lectura del código. Además, el modelo ajustado de regresión logística no está almacenado en ninguna variable, por lo que no podemos inspeccionarlo o usarlo para predecir en cualquier otro data.			
		\end{minipage} 
}}\\







\subsection{Clasificadores  Naive   Bayes\index{naive bayes classifiers}} 


Los clasificadores Naive Bayes   (Ingenuos de Bayes) son una familia de clasificadores bastante similares a los  modelos lineales  discutidos en la sección anterior.  Sin embargo, tienden a ser aún más rápidos en el entrenamiento. El precio pagado por esta eficiencia es que los modelos de  naive de Bayes  a menudo proporcionan un rendimiento de generalización que es ligeramente má bajo (malo o peor) que el de los clasificadores lineales como {\ttfamily LinearSVC} y {\ttfamily LogisticRegression}.  La razón de que los modelos naives de Bayes   son tan eficientes se debe a que los parámetros aprenden sobre cada característica individualmente y recopilan las estadísticas simples por clase de cada característica. Hay tres tipos de clasificadores de naives de Bayes  implementados en {\ttfamily  scikit-learn}: 
 GaussianNB, BernoulliNB, y MultinomialNB. 
 
 GaussianNB se puede aplicar a 	cualquier dato continuo, mientras que BernoulliNB asume datos binarios y MultinomialNB asume datos de conteo absoluto (es decir, que cada característica representa un número entero de algo, por ejemlo, con qué frecuencia aparece una palabra en una frase).  BernoulliNB y MultinomialNB se utilizan principalmente en la clasificación de datos de texto. El clasificador BernoulliNB cuenta con qué frecuencia cada característica de cada clase no es cero. Esto se entiende más fácilmente con un ejemplo: \\
 
	
\begin{lstlisting}
In[54]:  X = np.array([[0, 1, 0, 1],         #  dato1
                       [1, 0, 1, 1],         #  dato2
                       [0, 0, 0, 1],         #  dato3
                       [1, 0, 1, 0]])        #  dato4   
         y = np.array([0, 1, 0, 1])
	 
\end{lstlisting}
Aquí, tenemos cuatro puntos de datos, con cuatro características binarias cada uno. Hay dos clases, 0 y 1. Para la clase 0 (el primer y tercer punto de datos), la primera característica es cero dos veces y no cero cero veces, la segunda característica es cero una vez y no cero una vez, y así sucesivamente. Estos mismos recuentos se calculan entonces para los puntos de datos de la segunda clase. Contando las entradas no cero por clase en esencia se ve así:\\

\begin{lstlisting}
In[55]:  counts = {}
         for label in np.unique(y):
             # iterate over each class
             # count (sum) entries of 1 per feature
             counts[label] = X[y == label].sum(axis=0)
         print("Feature counts:\n{}".format(counts))
Out[55]:

    Feature counts:
    {0: array([0, 1, 0, 2]), 1: array([2, 0, 2, 1])}

\end{lstlisting}
		

\noindent Los otros dos modelos naive Bayes, MultinomialNB y GaussianNB, son ligeramente diferentes en el tipos de estadísticas que  calculan.  MultinomialNB tiene en cuenta el valor promedio de cada característica por cada clase, mientras que GaussianNB almacena el valor promedio así como la desviación estándar de cada característica por cada clase.


Para hacer una predicción, se compara un punto de datos con las estadísticas de cada una de las clases,
y se predice la clase con más coincidencia.
 Curiosamente, para MultinomialNB y BernoulliNB, esto conduce a la misma  forma de predicción de los modelos lineales  (véase "Modelos lineales para la clasificación"{} en la página \pageref{Mlineales}).
Desafortunadamente, el {\ttfamily coef\_} para los modelos naives de Bayes   tiene un significado algo diferente que en los modelos lineales, en que {\ttfamily coef\_}  no es lo mismo que $w$.

\subsubsection{Fortalezas, debilidades y parámetros} 

 MultinomialNB y BernoulliNB tienen un único parámetro, {\ttfamily alfa}, que controla la complejidad del modelo. 
 La forma en que {\ttfamily alfa}  funciona consiste en que algoritmo añade a los datos  muchos puntos {\ttfamily alfa} de datos virtuales que tienen valores positivos para todas las características.  Esto resulta en un "suavizado"{}  de las estadísticas. Un {\ttfamily alfa} grande significa mayor suavizado, resultando  modelos menos complejos. El rendimiento del algoritmo es relativamente robusto al ajuste de {\ttfamily alfa}, lo que significa que el ajuste de {\ttfamily alfa} no es crítico para un buen rendimiento.  Sin embargo, por lo general, la afinación mejora un poco la exactitud.\\
 
El GaussianBN se utiliza principalmente en datos con dimensión muy grande, mientras que los otros dos  modelos se utilizan ampliamente para  recuento de datos sparse, como en el caso de  texto. El modelo MultinomialNB generalmente tiene mejores resultados que el  BernoulliNB, particularmente en conjuntos de datos con un número relativamente grande de características distintas de cero (es decir, documentos grandes).\\

Los  modelos naive Bayes   comparten muchas de las fortalezas y debilidades de los modelos lineales. Son muy rápidos para entrenar y predecir, y el procedimiento de entrenamiento es fácil entender. Los modelos funcionan muy bien con datos sparse de alta dimensión y son relativamente robusto a los parámetros. Los modelos naive Bayes  son excelentes modelos de referencia y   a menudo se usan en conjuntos de datos muy grandes, donde la capacitación o el entrenamiento para un modelo lineal puede tardar demasiado.

\subsection{Árboles de decisión\index{árboles de decisión}}

Los árboles de decisión son modelos ampliamente utilizados para tareas de clasificación y regresión. Esencialmente,  aprenden una jerarquía de preguntas if{\scriptsize{/}}else, que lo que lleva a una decisión. \\

Estas preguntas son similares a las que podrías hacer en un juego de 20 preguntas. Imagine que desea distinguir entre los siguientes cuatro animales: osos, halcones, pingüinos y delfines. Su objetivo es llegar a la respuesta correcta haciendo el menor número posible de preguntas. Podría empezar preguntando
 si el animal tiene plumas, una pregunta que reduce el problema a sólo dos  animales.
Si la respuesta es  "{}si"{}, puede hacer otra pregunta que podría ayudarlo a distinguir entre halcones y
pingüinos. Por ejemplo, podría preguntar si el animal puede volar. Si el animal no tiene plumas, sus posibles elecciones de animales son delfines y osos, y  tendrás que hacer una pregunta para distinguir entre estos dos animales, por ejemplo, preguntando si el animal tiene aletas. \\

Esta serie de preguntas puede expresarse como un árbol de decisión, como se muestra en la Figura \ref{Arbol1}.\\

\begin{lstlisting}
In[56]:  mglearn.plots.plot_animal_tree()
\end{lstlisting}

En la Figura \ref{Arbol1} se ilustra el método aplicado al problema de los cuatro animales, cada nodo en el árbol representa una pregunta o un nodo terminal (también llamado hoja) que contiene la respuesta. Los bordes conectan las respuestas a un pregunta con la siguiente pregunta que harías.\\

\begin{figure}[h!]
	\centering
	\includegraphics[scale=0.55]{../Imagen_Machine/Arbol1}
	\caption{Un árbol de decisión para distinguir entre varios animales}
	\label{Arbol1}
\end{figure}



En lenguaje de aprendizaje automático, construimos un modelo para distinguir entre cuatro clases de
animales (halcones, pingüinos, delfines y osos) que usan las tres características "tiene plumas"{}, "puede volar"{}  y  "tiene aletas"{}.   En lugar de construir estos modelos a mano, podemos aprender de ellos a partir de datos mediante el aprendizaje supervisado.


\subsubsection{Construyendo árboles de decisión}


Veamos el proceso de construcción\index{construcción de un árbol} de un árbol de decisión en un conjunto de datos  para una clasificación 2D, que se muestra en la Figura \ref{arbol2D}. El conjunto de datos consta de dos formas de media luna,  cada clase que consta de 75 puntos de datos. Nos referiremos a este conjunto de datos como {\ttfamily two\_moons}.\\

 Un árbol de decisión aprendiendo,  significa aprender la secuencia de preguntas if{\scriptsize/}else que nos lleva a la  respuesta verdadera de forma muy rápida. En el entorno de aprendizaje automático, estas preguntas son llamadas pruebas (que no deben confundirse con el conjunto de pruebas, que son los datos que usamos para probar  qué tan bien generaliza el modelo). Por lo general, los datos no vienen en forma de características binarias sí{\scriptsize/}no  como en el ejemplo animal, en su lugar,  se representa como características continuas como en el conjunto de datos 2D que se muestra en la Figura \ref{arbol2D}. Las pruebas o preguntas que se utilizan en los datos continuos tienen una  forma como:
 
                 "{}¿ Si la característica es mayor que el valor $a$? "{}. \\
 



\begin{figure}[h!]
	\centering
	\includegraphics[scale=0.60]{../Imagen_Machine/arbol2D}
	\caption{Conjunto de datos de dos lunas en el que se construirá el árbol de decisión}
	\label{arbol2D}
\end{figure}



Para construir un árbol, el algoritmo busca  todas las pruebas posibles y encuentra la que es
más informativa sobre la variable objetivo. La figura \ref{Arbol3} muestra la primera prueba que es
seleccionada. Dividiendo el conjunto de datos verticalmente en $x[1]=0.0596$  produce la mayor cantidad de información; esta es la mejor separación de los puntos de la clase 1 y de la clase 2.  El nodo superior, también es llamado  Raíz, y representa el conjunto de datos completo,  que en este caso consta de 75 puntos que pertenecen a la clase 0 y 75 puntos pertenecientes a la clase 1. La división se realiza probando si $x[1]<=0.0596$, indicado por una línea negra.  Si la prueba es verdadera, se asigna un punto al nodo izquierdo, que contiene 2  puntos que pertenecen a la clase 0 y 32 puntos que pertenecen a la clase 1.   De lo contrario el punto se asigna al nodo derecho, que contiene  48  puntos que pertenecen a la clase 0 y 18 puntos pertenecientes a la clase 1. Estos dos nodos corresponden a la parte superior y inferior de las regiones que se muestran en la Figura \ref{Arbol3}.  A pesar de que la primera división hizo un buen trabajo de separación clasificando las dos clases, la región inferior todavía contiene puntos que pertenecen a la clase 0, y la región superior todavía contiene puntos que pertenecen a la clase 1.  Podemos construir un  modelo más preciso  repitiendo el proceso de búsqueda de la mejor prueba en ambas regiones. La Figura \ref{Arbol4} muestra que próxima  división más informativa para la región de la izquierda y de la derecha se basa en $x[0]$.



\begin{figure}[h!]
	\centering
	\includegraphics[scale=0.50]{../Imagen_Machine/Arbol3}
	\caption{Límite de decisión del árbol con profundidad 1 (izquierda) y árbol correspondiente (derecha)}
	\label{Arbol3}
\end{figure}


\begin{figure}[h!]
	\centering
	\includegraphics[scale=0.50]{../Imagen_Machine/Arbol4}
	\caption{Límite de decisión del árbol con profundidad 2 (izquierda) y el árbol de decisión correspondiente (derecha)}
	\label{Arbol4}
\end{figure}


Este proceso recursivo produce un árbol binario de decisiones, con cada nodo que contiene una prueba.    Alternativamente, se puede pensar en cada prueba como la división de la parte de los datos que actualmente se está considerando a lo largo de un eje.  Esto produce una visión del algoritmo como la construcción de una partición jerárquica.    Como cada prueba concierne a una sola característica, las regiones de la partición resultante siempre tienen límites de eje paralelo.    La partición recursiva de los datos se repite hasta que cada región de la partición (cada hoja en el árbol de decisiones) sólo contiene un único valor objetivo (una sola clase o un único valor de regresión).    Una hoja del árbol que sólo contiene puntos de datos que comparten el mismo valor objetivo se llama hojas puras. La  partición final de este conjunto de datos se muestra en la Figura \ref{Arbol4}.\\



\begin{figure}[h!]
	\centering
	\includegraphics[scale=0.50]{../Imagen_Machine/Arbol5}
	\caption{Límite de decisión del árbol con profundidad 9 (izquierda) y parte del correspondiente
		árbol (derecha); el árbol completo es bastante grande y difícil de visualizar}
	\label{Arbol5}
\end{figure}


Una predicción para un nuevo punto de datos se hace chequeando en qué región de la partición del espacio de la característica se encuentra el punto, y después predice el objetivo por la mayoría (o el objetivo único en el caso de hojas puras) en esa región. La región se puede encontrar atravesando el árbol desde la raíz, yendo a izquierda o derecha, dependiendo de si la prueba se cumple o no. \\

También es posible usar árboles para tareas de regresión, usando exactamente la misma técnica. Para hacer una predicción, atravesamos el árbol basado en las pruebas en cada nodo y encontramos la hoja en la que cae el nuevo punto de datos. El resultado para este punto de datos es el objetivo medio de los puntos de entrenamiento en esta hoja.


\subsubsection{Controlar la complejidad de los árboles de decisión.}

Típicamente, construir un árbol como se describe aquí y continuar hasta que todas las hojas sean puras
conduce a modelos que son muy complejos y altamente ajustados a los datos de entrenamiento.  La presencia de hojas puras significa que el árbol es 100\% exacto en el conjunto de entrenamiento; cada punto de datos en el conjunto de entrenamiento está en una hoja que tiene la clase de mayoría correcta.  El sobreajuste se puede ver a la izquierda de la Figura \ref{Arbol5}. Donde se se observan las regiones determinadas para la clase 1 en medio de todos que pertenecen a la clase 0. Por otro lado, hay una pequeña franja predicha como clase 0 alrededor del punto que pertenece a la clase 1 a muy a la derecha. Esta no es la forma en que uno podría imaginar el límite de decisión para observar, y el límite de decisión se centra mucho en los puntos atípicos únicos que están lejos de los otros puntos de esa clase.\\



Hay dos estrategias comunes para evitar el sobreajuste: detener la creación del árbol temprano (también llamado pre-poda o poda previa, del inglés pre-pruning), o construir el árbol y luego  remover los nodos que contienen poca información (también llamado  poda posterior  o simplemente poda, del inglés post-poda). Los criterios posibles para la pre-poda   incluyen limitar la profundidad máxima del árbol, limitar el número máximo de hojas, o requerir un número mínimo de puntos en un nodo para seguir dividiéndolo. 

 En {\ttfamily scikit-learn} los árboles de decisión se implementan para regresión con {\ttfamily DecisionTreeRegressor}, y con {\ttfamily DecisionTreeClassifier} para clasificar; además, scikit-learn solo implementa la  pre-poda, y no post-poda. \\


El efecto de la pre-poda lo podemos ver con más detalle en el conjunto de datos de Cáncer de Mama (Breast Cancer).  La metodología se implementa como siempre, primero importamos el conjunto de datos y lo dividimos en entrenamiento y prueba, y luego se construye el modelo. Si se usa la  configuración por defecto, se  desarrolla completamente el árbol (el árbol crece  hasta que todas las hojas son puras). Y fijamos el {\ttfamily random\_state} en el árbol, que se usa para desempate internamente:\\


\begin{lstlisting}
In[58]:  from sklearn.tree import DecisionTreeClassifier
      
         cancer = load_breast_cancer()
         X_train, X_test, y_train, y_test = train_test_split(cancer.data,
              cancer.target, stratify=cancer.target, random_state=42)
         tree = DecisionTreeClassifier(random_state=0)
         tree.fit(X_train, y_train)
         print("Accuracy on training set: {:.3f}".format(tree.score(
                X_train, y_train)))
         print("Accuracy on test set: {:.3f}".format(tree.score(
                X_test, y_test)))

Out[58]:

      Accuracy on training set: 1.000
      Accuracy on test set: 0.937

\end{lstlisting}


Como era de esperar, la exactitud en el conjunto de entrenamiento es del 100\%, porque las hojas son puras, el árbol creció lo suficientemente profundo como para poder memorizar perfectamente todas las etiquetas en los datos de entrenamiento. La exactitud del conjunto de prueba es ligeramente menor a la del conjoto de entrenamiento, o mala y aun peor que en los modelos lineales que vimos anteriormente, que tenían alrededor de 95\% de exactitud.\\


Si no restringimos la profundidad de un árbol de decisión, el árbol puede volverse arbitrariamente profundo
y complejo. Por lo tanto, los árboles no podados son propensos al sobreajuste y en consecuencia no generalizan
bien (sobre nuevos datos).  Ahora apliquemos la poda previa al árbol, que limita el desarrollo
del árbol antes de alcanzar un ajuste perfecto con los datos de entrenamiento.  Una opción es dejar de construir el árbol después de que se haya alcanzado cierta profundidad. En este caso establecemos {\ttfamily max\_depth} = 4, que significa solo se pueden hacer cuatro preguntas consecutivas (ver Figuras \ref{Arbol3} y \ref{Arbol5}). Limitando la profundidad del árbol disminuye el sobreajuste, lo cual debe conducir a una menor exactitud en el conjunto de entrenamiento:\\


\begin{lstlisting}
In[59]:  tree = DecisionTreeClassifier(max_depth=4, random_state=0)
         tree.fit(X_train, y_train)
         print("Accuracy on training set: {:.3f}".format(tree.score(
                X_train, y_train)))
         print("Accuracy on test set: {:.3f}".format(tree.score(
                X_test, y_test)))

Out[59]:

    Accuracy on training set: 0.988
    Accuracy on test set: 0.951
\end{lstlisting}


\subsubsection{Analizando árboles de decisión}

Podemos visualizar el árbol empleado  la función {\ttfamily{export\_graphviz}} del módulo {\ttfamily .tree}  (árbol). Esta función  escribe un archivo  de texto en formato  .dot. Establecemos una opción para colorear los nodos a fin  reflejar la clase mayoritaria en cada nodo, y se dan los nombres de clase y características para que el árbol pueda etiquetarse correctamente:\\


\begin{lstlisting}
In[60]:  from sklearn.tree import export_graphviz
         export_graphviz(tree, out_file="tree.dot", 
                         class_names=["malignant", "benign"],
                         feature_names=cancer.feature_names, 
                         impurity=False, filled=True)

\end{lstlisting}

Podemos leer este archivo y visualizarlo, como se ve en la Figura \ref{ArbolCacer}, usando el modulo {\ttfamily graphviz} o  usar cualquier programa que pueda leer archivos .dot:\\

\begin{lstlisting}
In[61]:  import graphviz
         
         with open("tree.dot") as f:
             dot_graph = f.read()
         graphviz.Source(dot_graph)

\end{lstlisting}



\begin{figure}[h!]
	\centering
	\includegraphics[scale=0.50]{../Imagen_Machine/ArbolCacer}
	\caption{Límite de decisión del árbol con profundidad 9 (izquierda) y parte del correspondiente
		árbol (derecha); el árbol completo es bastante grande y difícil de visualizar}
	\label{ArbolCacer}
\end{figure}



La visualización del árbol proporciona una idea muy clara de cómo el algoritmo
hace predicciones y es un buen ejemplo de un algoritmo de aprendizaje automático que puede ser 
explicado fácilmente a los inexpertos.  Sin embargo, incluso con un árbol de profundidad, como se ve aquí,  el árbol puede volverse un poco abrumador. Árboles más profundos (una profundidad de 10 es poco común, pero puede ser) son aún más difíciles de comprender.

Un método  útil  para inspeccionar el árbol  con la finalidad  de averiguar qué ruta toma la mayoría de los datos, son las  {\ttfamily n\_samples}  (n\_muestras) que se muestran en cada nodo, como en la Figura \ref{ArbolCacer}  donde se ve claramente el  número de muestras en cada nodo, mientras que el valor da el número de muestras por clase. Siguiendo las ramas a la derecha, vemos que  worst radius $<=16.795$ (peor radio)  crea un nodo que contiene  8 benignos y 134  malignas.  El resto de este lado del árbol utiliza algunas distinciones más finas para dividir estas 8 muestras benignas restantes.   De las 142 muestras que quedaron  a la derecha en la división inicial, casi todos (132) terminan en la hoja de la derecha.\\

Tomando  a la izquierda en la raíz, para el worst radius $> 16.795$ (peor radio) terminamos con 25 malignos
y 259 muestras benignas. Casi todas las muestras benignas terminan en la segunda hoja, 
a la derecha, la mayoría de las otras hojas  contienen muy pocas muestras.\\


\subsubsection{Importancia de la característica en árboles} \label{ImpotCaract}

 En lugar de observar todo el árbol  que puede ser  agobiante,  podemos usar algunos propiedades 
 derivadas del resumen  de la ejecusión del árbol. El  resumen más común  utilizado es la importancia de la característica, que califica la importancia de cada característica por la decisión que toma un árbol.  Este es un número entre 0 y 1 para cada característica, donde 0 significa  "{}no utilizado en absoluto"{},  y 1 significa  "{}predice perfectamente el objetivo"{}, y la sumatoria  total de las importancias  de las características  es igual a 1:\\
  
\begin{lstlisting}
In[62]:  print("Feature importances:\n{}".format(tree.feature_importances_))

Out[62]:

    Feature importances:
    [ 0.     0.   0.     0.     0.     0.     0.     0.    0.   0.     0.01
      0.048  0.   0.     0.002  0.     0.     0.     0.    0.   0.727  0.046
      0.     0.   0.014  0.     0.018  0.122  0.012  0. ]
\end{lstlisting}


Podemos visualizar las importancias de las características de una manera similar a la forma en que visualizamos los coeficientes en el modelo lineal (Figure \ref{ImpotanciaCaract}):



\begin{lstlisting}
  
In[63]:  def plot_feature_importances_cancer(model):
             n_features = cancer.data.shape[1]
             plt.barh(range(n_features), model.feature_importances_, 
                      align='center')
             plt.yticks(np.arange(n_features), cancer.feature_names)
             plt.xlabel("Feature importance")
             plt.ylabel("Feature")
             plot_feature_importances_cancer(tree)
\end{lstlisting}



\begin{figure}[h!]
	\centering
	\includegraphics[scale=0.50]{../Imagen_Machine/ImpotanciaCaract}
	\caption{Importancia de características calculadas a partir de un árbol de decisión aprendido en el  conjunto de datos de Cáncer de Mama}
	\label{ImpotanciaCaract}
\end{figure}


Aquí vemos que la característica utilizada en la división superior es "{}worst radius"{}  (peor radio),  y resulta ser  la característica más importante. Esto confirma nuestra observación al analizar el árbol que en el primer nivel ya separa  bastante bien las dos clases. 

Sin embargo, si una característica aparece con  un puntaje  bajo  de importancia  en  {\ttfamily{ feature\_importance}},  no significa que esta característica no aporte información, sólo significa que la característica no fue seleccionada por el árbol, probablemente porque otra característica codifica la misma información.\\


En contraste con los coeficientes de los modelos lineales, la  importancia de características en este caso son siempre positivas, y no codifican de qué clase es indicativa una característica. El valor de  importancia de la  característica nos dicen que el "{}worst radius"{ }  es importante; pero no dice, si un radio alto es indicativo de una muestra que es benigna o maligna. De hecho, puede que no haya una relación tan simple entre características y clase, como se puede ver en el siguiente ejemplo (Figuras \ref{EtiquetArbol1}  y \ref{EtiquetArbol1} ):\\

 

\begin{lstlisting}
In[64]:  tree = mglearn.plots.plot_tree_not_monotone()
         display(tree)
Out[64]:

    Feature importances: [ 0. 1.]

\end{lstlisting}


\begin{figure}[h!]
	\centering
	\includegraphics[scale=0.45]{../Imagen_Machine/EtiquetArbol1}
	\caption{Conjunto de datos bidimensionales en el que la característica en el eje y tiene una relación no monótona con la etiqueta de clase, y los límites de decisión encontrados por un árbol de decisiones}
	\label{EtiquetArbol1}
\end{figure}



\begin{figure}[h!]
	\centering
	\includegraphics[scale=0.50]{../Imagen_Machine/EtiquetArbol2}
	\caption{Árbol de decisión aprendido en los datos mostrados en la Figura \ref{EtiquetArbol1}}
	\label{EtiquetArbol2}
\end{figure}



El gráfico \ref{EtiquetArbol1} muestra un conjunto de datos con dos características y dos clases. En este caso  toda la información está contenido en $X[1]$, y $X[0]$ no se utiliza en absoluto. Pero la relación entre $X[1]$ y la clase de salida no es monótona, lo que significa que no podemos decir que "{}un valor alto de $X[0]$ significa clase 0, y un valor bajo significa clase 1{}  o viceversa.\\

Mientras centramos nuestra discusión aquí en los árboles de decisión para la clasificación, todo lo dicho  es igualmente válido en árboles de decisión para la regresión, que se implementa en la
 {\ttfamily{DecisionTreeRegressor}}.  El uso y análisis de los árboles de regresión es muy similar al de
Árboles de clasificación. Sin embargo, hay una propiedad particular e  importante de tener presente al usar modelos basados en árboles para la regresión, se trata de que el  {\ttfamily{DecisionTreeRegressor}} (y todos los otros modelos de regresión basados en árboles) no es capaz de extrapolar, o hacer predicciones fuera del rango de los datos de entrenamiento.\\


Veamos  esto con más detalle, utilizando un conjunto de datos de los precios históricos de las memorias  para un ordenador (RAM). La Figura \ref{RAM} muestra el conjunto de datos, con la fecha en el eje $x$ y el precio de un megabyte de RAM en ese año en el eje $y$:\\


\begin{lstlisting}
In[65]:  import pandas as pd
         ram_prices = pd.read_csv("data/ram_price.csv")
         
         plt.semilogy(ram_prices.date, ram_prices.price)
         plt.xlabel("Year")
         plt.ylabel("Price in $/Mbyte")
\end{lstlisting}



\begin{figure}[h!]
	\centering
	\includegraphics[scale=0.50]{../Imagen_Machine/RAM}
	\caption{Evolución histórica del precio de la RAM, trazada en una escala logarítmica}
	\label{RAM}
\end{figure}


Se plantea  un pronóstico para los años posteriores a 2000, utilizando los datos históricos hasta ese 
punto, con la fecha como única  característica en el eje $x$.  Al re-escalar los precios, en este caso  $y$, mediante el logarítmo, la relación se hace relativamente lineal con algunas pequeñas variaciones,   y por lo tanto  será más fácil de predecir con los modelos lineales. Se costruyen y se comparan los modelos sobre la misma gráfica: mediante {\ttfamily DecisionTreeRegressor} y {\ttfamily LinearRegression}.  Después de entrenar  los modelos,  se hacen las predicciones sobre el conjunto  de prueba y se evalúa:\\


%   Esto no hace una diferencia para {\ttfamily DecisionTreeRegressor}, pero hace una gran diferencia para   {\ttfamily LinearRegression} ( discutiremos esto con más profundidad en el Capítulo \ref{Cap4}). 

\begin{lstlisting}
In[66]:  from sklearn.tree import DecisionTreeRegressor
         # use historical data to forecast prices after the year 2000
         data_train = ram_prices[ram_prices.date < 2000]
         data_test = ram_prices[ram_prices.date >= 2000]
         
         # predict prices based on date
         X_train = data_train.date[:, np.newaxis]
         # we use a log-transform to get a simpler relationship of data to target
         y_train = np.log(data_train.price)
         
         tree = DecisionTreeRegressor().fit(X_train, y_train)
         linear_reg = LinearRegression().fit(X_train, y_train)
         
         # predict on all data
         X_all = ram_prices.date[:, np.newaxis]
         
         pred_tree = tree.predict(X_all)
         pred_lr = linear_reg.predict(X_all)
         
         # undo log-transform
         price_tree = np.exp(pred_tree)
         price_lr = np.exp(pred_lr)
\end{lstlisting}



\begin{figure}[h!]
	\centering
	\includegraphics[scale=0.50]{../Imagen_Machine/RAM2}
	\caption{Comparación de predicciones hechas por un modelo lineal y las predicciones hechas
		por un árbol de regresión sobre los datos de precios de la memoria RAM}
	\label{RAM2}
\end{figure}


\begin{lstlisting}
In[67]:  plt.semilogy(data_train.date, data_train.price, 
                       label="Training data")
         plt.semilogy(data_test.date, data_test.price, label="Test data")
         plt.semilogy(ram_prices.date, price_tree, label="Tree prediction")
         plt.semilogy(ram_prices.date, price_lr, label="Linear prediction")
         plt.legend()
\end{lstlisting}


La diferencia entre los modelos es bastante sorprendente. El modelo lineal se aproxima a
los datos con una línea, como lo esperado.  Esta línea proporciona un pronóstico bastante bueno para los datos de prueba (los años posteriores a 2000), mientras pasa por alto algunas de las variaciones más finas en tanto el entrenamiento como los datos de prueba.  Por otro lado,  el modelo de árbol hace perfecto predicciones sobre los datos de entrenamiento; claro, no restringimos la complejidad del árbol, por lo que aprendió todo el conjunto de datos.  Sin embargo, una vez que dejamos el rango de la data entrenamiento, el modelo simplemente sigue prediciendo por el último punto conocido.  El árbol no tiene la capacidad para generar  "{}nuevas"{}  respuestas, fuera de lo que memorizó o aprendió en la data de entrenamiento. Esta deficiencia se presenta  en todos los modelos basados en árboles. \\

\subsubsection{Fortalezas, debilidades y parámetros}


Como se discutió anteriormente, los parámetros que controlan la complejidad del modelo de árboles de decisión son los parámetros como el de pre-poda,  que detienen la construcción del árbol antes de que esté completamente desarrollado. Por lo general, elegir una de las estrategias que controlan la complejidad del modelo (como {\ttfamily max\_depth,  max\_leaf\_nodes}  o {\ttfamily min\_samples\_leaf}): es suficiente para evitar el sobreajuste.\\

Los árboles de decisión tienen dos ventajas sobre muchos de los algoritmos que hemos discutido:
Primero el modelo resultante puede ser fácilmente visualizado y comprendido por inexpertos (al menos en árboles más pequeños).  Segundo los algoritmos son completamente invariate  a la escala de la data.  Como cada característica es procesada por separado,  las  divisiones posibles de la data no dependerá de la escala, no es necesario ningún proceso previo como la normalización o estandarización de características para los algoritmos del árbol de decisiones.   En particular, los árboles de decisión funcionan bien cuando
tiene características que están en escalas completamente diferentes, o una combinación de binario y continuas.\\


El principal inconveniente de los árboles de decisión es que incluso con el uso de la pre-poda,
tienden a sobreajustar y en consecuencia proporciona un rendimiento de generalización deficiente.  Por lo tanto, para eludir este problema se plantean los  métodos de conjuntos de áboles (Ensembles of Decision Trees). Generalmente en la mayoría aplicaciones  se usa estos  métodos,  en lugar de un simple  árbol de decisión. \\

A continuación se tratan algunos de los métodos de conjuntos de áboles. 


\subsection{Conjuntos de árboles de decisión \\
	(Ensembles of Decision Trees)}


Los métodos de conjunto  son técnicas  que combinan múltiples modelos de aprendizaje automático para crear modelos más potentes. Existen muchos modelos en la literatura de aprendizaje automático que pertenecen a esta categoría, pero hay dos modelos de conjunto que han demostrado ser eficaz en una amplia gama de conjuntos de datos para clasificación y regresión,  ambos usan árboles de decisión en sus bloques de construcción: bosques aleatorios (random forests) y  árboles de decisión potenciado  
o simplemente potenciación del gradiente (gradient boosted decision trees).  

\subsubsection{Bosques Aleatorios}


El pricipal inconveniente  de los árboles de decisión, es que tienden a generar sobreajuste en la  data de entrenamiento. Un forma de abordar este problema es con  bosques aleatorios\index{bosques aleatorios}. Esencialmente un bosque aleatorio es  una colección de árboles de decisión, donde cada árbol es ligeramente diferente de los demás. La idea de fondo de los bosques aleatorios es reducir el sobreajuste en los modelos de árbol de decisión, pues cada árbol seguro hace un buen trabajo de predicción, pero probablemente  sobreajustará  gran parte de los datos. Si construimos muchos  árboles, todos los cuales funcionan bien y se sobreajustan de diferentes maneras, podemos reducir la cantidad  del sobreajuste promediando sus resultados. Esta reducción en el sobreajuste, mientras se conserva  el poder predictivo de los árboles puede ser  demostrado usando matemáticas rigurosas (para mayor información revisar The Elements of Statistical Learning (Springer) by Trevor Hastie, Robert Tibshirani, and Jerome Friedman). \\

Para implementar esta estrategia, necesitamos construir muchos árboles de decisión.  Cada árbol debe hacer un trabajo aceptable de predicción del objetivo, y también debe ser diferente de los otros arboles. Los bosques aleatorios obtienen su nombre al incorporar la aleatoriedad en la construcción del árbol.
 
Para garantizar la diferencia entre los árboles, hay dos formas de construir los árboles que formarán los bosques: una consiste en  seleccionar diferentes conjuntos de puntos de datos para usarlos en la construción del árbol, y la otra forma es emplear una  selección aleatoria ponderada de  las características en cada nodo a dividir.  En la próxima sección se analiza  este proceso con más detalle.\\


\subsubsection{Construyendo Bosques  Aleatorios}


Para construir un modelo de bosque aleatorio, necesitamos decidir sobre el número de árboles a emplear. Esta acción se define en el parámetro {\ttfamily n\_estimators} de {\ttfamily RandomForestRegressor} o {\ttfamily RandomForestClassifier}.   Por ejemplo, decidimos  construir un bosques con 10 árboles, estos árboles serán construido completamente independiente  uno del otro, y el algoritmo hará  diferentes selecciones aleatorias  para cada árbol, asegurado que los árboles sean distintos.   
Para construir un árbol, primero tomamos lo que se llama una muestra de Bootstrap (muestra de arranque) en la data. El proceso consiste en extraer n\_muestras con reemplazo de la data (n\_muestras puntos de datos), es decir,  el conjunto creado que es tan grande como el conjunto de datos original, pero algunos puntos de datos se repetirán varias veces, y otros faltarán, aproximadamente un tercio. \\

Para ilustrar esta situación, supongamos que queremos crear un bootstrap de la lista ["{}a"{} , "{}b"{}, "{}c"{}, "{}d"{}].
 Una posible muestra de bootstrap sería ["{}b"{}, "{}d"{}, "{}d"{}, "{}c"{}]. Otro posible muestra sería ["{}d"{}, "{}a"{}, "{}d"{}, "{}c"{} ].\\



El paso siguiente es costruir un árbol de decisión basado en este conjunto de datos recién creado. Sin embargo, el algoritmo que describimos para el árbol de decisión está ligeramente modificado. En lugar de buscar la mejor prueba para cada nodo, en cada nodo el algoritmo selecciona aleatoriamente un subconjunto de las características, y busca la mejor prueba posible que involucre una de estas características. El número de características a seleccionar es controlado por el parámetro {\ttfamily max\_features}. Esta selección de un subconjunto de características se repite por separado en cada nodo, de modo que cada  nodo en un árbol puede tomar una decisión utilizando un subconjunto diferente de las características.\\


El bootstrap  permite que  cada árbol de decisión, en el bosque aleatorio empiece su construcción 
sobre un conjunto de datos ligeramente diferente.  Debido a la elección aleatoria de características en cada nodo, la dividisión  en cada árbol opera con un subconjunto diferente de características. Estos dos mecanismos juntos aseguran que todos los árboles en el bosque aleatorio sean diferentes.\\



Un parámetro crítico en este proceso es {\ttfamily{ max\_features}}. Observe, si se configura {\ttfamily max\_features}  para   {\ttfamily{n\_features}},  significa que cada división puede ver todas las características del conjunto de datos y no ejecutará la aleatoriedad en la selección de características (sin embargo, la aleatoriedad debida al bootstrap permanece).  Si configuramos {\ttfamily max\_features} igual a  1, significa que las divisiones no tienen ninguna opción en absoluto sobre qué característica probar, y sólo pueden buscar sobre diferentes umbrales para la característica que fué seleccionada  al azar. Por lo tanto, un alto {\ttfamily{max\_features}} significa que los árboles en el bosque aleatorio serán bastante similares, y serán capaz de ajustar los datos fácilmente, utilizando las características más distintivas. Un valor bajo para  {\ttfamily{max\_features}},  significa que los árboles en el bosque aleatorio serán bastante diferentes, y que cada árbol podria necesitar ser muy profundo para ajustarse bien a los datos.\\
 
 
Para hacer una predicción usando el bosque aleatorio, el algoritmo primero hace una predicción
por cada árbol en el bosque.  Para la  regresión,  promedia  estos resultados para obtener nuestro predicción final.  Y en el caso de la clasificación, se utiliza una estrategia de "{}votación suave"{} ( "{}soft voting"{}). Esto significa que cada algoritmo hace una predicción "{}suave"{}, proporcionando una probabilidad para cada etiqueta de salida posible. Se promedian las probabilidades predichas por todos los árboles y la clase con la mayor probabilidad  se predice.

\subsubsection{Analizando bosques Aleatorio} 

Para analizar esta técnica, construimos un bosque aleatorio con cinco árboles de decisión sobre la data  {\ttfamily two\_moons}:\\



\begin{lstlisting}
In[68]: from sklearn.ensemble import RandomForestClassifier
	from sklearn.datasets import make_moons
	
	X, y = make_moons(n_samples=100, noise=0.25, random_state=3)
	X_train, X_test, y_train, y_test = train_test_split(X, y, 
	       stratify=y, random_state=42)
	       
	forest = RandomForestClassifier(n_estimators=5, random_state=2)
	forest.fit(X_train, y_train)
\end{lstlisting}

Los árboles que se construyen como parte del bosque aleatorio se almacenan en el atributo {\ttfamily{estimador\_}}.  Los límites de decisión aprendidos por cada árbol, junto con su predicción  generada por el bosque se visualiza en la Figura \ref{Prome5Arbol}:\\


\begin{lstlisting}
In[69]: fig, axes = plt.subplots(2, 3, figsize=(20, 10))
	for i, (ax, tree) in enumerate(zip(axes.ravel(), 
	                                   forest.estimators_)):
        ax.set_title("Tree {}".format(i))
        mglearn.plots.plot_tree_partition(X_train,
    	                                  y_train, tree, ax=ax)
    	
	mglearn.plots.plot_2d_separator(forest, X_train,fill=True,
	                                ax=axes[-1, -1], alpha=.4)
	                                
	axes[-1, -1].set_title("Random Forest")
	mglearn.discrete_scatter(X_train[:, 0], X_train[:, 1], y_train)
\end{lstlisting}



	\begin{figure}[h!]
	\centering
	\includegraphics[scale=0.60]{../Imagen_Machine/Prome5Arbol}
	\caption{Límites de decisión encontrados por cinco árboles de decisión aleatorios y el  límite decisión obtenido promediando sus probabilidades predichas}
	\label{Prome5Arbol}
\end{figure}



Se puede ver claramente, que los límites de decisión aprendidos por los cinco árboles son bastante diferentes. Cada uno de ellos comete algunos errores, ya que algunos de los puntos de entrenamiento que se trazan aquí, no se incluyeron en los conjuntos de entrenamiento de los árboles, debido a la muestra de bootstrap.\\

El bosque aleatorio se sobreajusta menos que cualquiera de los árboles individuales, y proporciona un
límite de decisión mucho más intuitivo. En cualquier aplicación real, usaríamos muchos
más árboles (a menudo cientos o miles), lo que lleva a límites aún más suaves.\\

 

			


Como otro ejemplo, apliquemos un bosque aleatorio que consta de 100 árboles sobre la data Cáncer de Mama (Breast Cancer):\\

\begin{lstlisting}
In[70]: X_train, X_test, y_train, y_test = train_test_split(
	     cancer.data, cancer.target, random_state=0)
	forest = RandomForestClassifier(n_estimators=100, random_state=0)
	forest.fit(X_train, y_train)
	
	print("Accuracy on training set: {:.3f}".format(forest.score(X_train, y_train)))
	print("Accuracy on test set: {:.3f}".format(forest.score(X_test, y_test)))

Out[70]:

	Accuracy on training set: 1.000
	Accuracy on test set: 0.972

\end{lstlisting}



El bosque aleatorio nos da una exactitud del 97\%, mejor que los modelos lineales o un único árbol de decisión, sin ajustar ningún parámetro. Podríamos configurar el {\ttfamily max\_features}, o aplicar la  pre-poda como se hizo  para el árbol de decisión simple. Sin embargo, a menudo los parámetros por defecto en el bosque aleatorio ya funcionan bastante bien. De manera similar al árbol de decisión, el bosque aleatorio proporciona las características más importantes, que se calculan agregando  a los bosques la sentencia {\ttfamily feature\_importances\_}. Típicamente, las características de importancia proporcionadas por el bosque aleatorio son más confiables que las proporcionados por un sólo árbol. Observe la figura \ref{ImporTanciBosque}.


\begin{lstlisting}  % [frame=single]	

In[71]: plot_feature_importances_cancer(forest)

\end{lstlisting}



\begin{figure}[h!]
	\centering
	\includegraphics[scale=0.50]{../Imagen_Machine/ImporTanciBosque}
	\caption{Importancia de características calculadas a partir del bosque aleatorio sobre la data de cáncer de mama}
	\label{ImporTanciBosque}
\end{figure}

Como puede ver, el bosque aleatorio le da una importancia distinta de cero a  más características,
que un solo árbol. De manera similar al árbol de decisión simple, el bosque aleatorio también da
gran importancia a la característica de “worst radius”, pero en  elige  “worst perimeter”  ({}"peor perímetro"{})  como  la característica más informativa.



 La aleatoriedad en la construcción de los bosques aleatorios obliga al algoritmo a considerar más explicaciones posibles, el resultado es que el bosque aleatorio captura  mucho más información  de los datos, que un solo árbol.




\subsubsection{Fortalezas, debilidades y parámetros}

 Bosques aleatorios para regresión y clasificación estan actualmente entre los métodos de aprendizaje automático más utilizados. Son muy potentes, y a menudo funcionan bien sin grandes ajuste de los parámetros, y no requieren escalar la data.  \\
 
 
 
 
 Esencialmente, los bosques aleatorios comparten todos los beneficios de los árboles de decisión, al tiempo que compensan algunas de sus deficiencias.  Una razón para  utilizar árboles de decisión es cuando se necesita una representación compacta del proceso de toma de decisiones.  Básicamente es imposible interpretar decenas o cientos de árboles en detalle, y los árboles en bosques aleatorios tienden a ser más profundo que los árboles de decisión (debido al uso de subconjuntos de características).  Por lo tanto si necesita resumir la predicción haciendola de una manera visual a los inexpertos, un árbol de decisión simple podría ser una buena opción.   Al construir bosques aleatorios para data grandes podría consumir algún  tiempo significante; sin embargo, puede ser tratado en paralelo a través de múltiples núcleos del CPU dentro de un ordenador fácilmente;  si está utilizando un procesador multi-core (como casi todos los ordenadores modernos lo hacen), puede utilizar el parámetro {\ttfamily n\_jobs} para ajustar el número de núcleos a utilizar. Usando más núcleos de CPU resultará en velocidad lineal (usando dos núcleos, el entrenamiento del bosque aleatorio será el doble de rápido), pero especificando {\ttfamily n\_jobs} más grande que el número de núcleos no ayudará. Puede configurar {\ttfamily n\_jobs=-1}  para usar todos los núcleos de su computadora.\\



 Debe tener en cuenta que los bosques aleatorios, por su naturaleza, son aleatorios y configurar
 diferentes estados aleatorios ({\ttfamily random\_state}) o no configurar  puede cambiar drásticamente 
 el modelo que se construye. Cuantos más árboles haya en el bosque, más robusto será el modelo contra la elección del estado aleatorio.  Además, si quieres tener resultados reproducibles, es importante ajustar el random\_state.\\


Los bosques aleatorios  pueden presentar un  rendimiento bajo sobre datos con dimensiones muy altas y sparse, como el caso de  datos de texto. Para este tipo de datos, los modelos lineales podrían ser más apropiados. Los bosques aleatorios por lo general funcionan bien incluso en datas muy grandes, y  el entrenamiento puede ser fácilmente paralelizada sobre varios núcleos de CPU dentro de un ordenador potente. Sin embargo,  los bosques aleatorios requieren más memoria y son más lentos para entrenar y predecir que el modelo lineal.  Si el tiempo y la memoria son importantes en una aplicación, podría tener sentido usar un
modelo lineal en su lugar. \\


Los parámetros más importantes para ajustar el modelo son  {\ttfamily n\_estimators,  max\_features} y  las opciones de pre-poda como {\ttfamily max\_depth}. Para valores altos de {\ttfamily n\_estimators}  siempre resultará mejor, y promediar más árboles producirá un modelo de conjunto (ensembles) más robusto,  pués se  reduce el sobreajuste. Sin embargo, los  rendimientos en cuanto a que más árboles necesitan más memoria y más tiempo para entrenar, deben tenerse en cuenta. Una regla general para  construir es  "{}Para todos los que tenga tiempo-memoria{}".\\


Como se describió anteriormente, el {\ttfamily max\_features}  determina qué tan aleatorio es cada árbol;  un valor   bajo   {\ttfamily max\_features} reduce  el sobreajuste. En general, es una buena regla usar los valores predeterminados como:  {\ttfamily max\_features = sqrt(n\_features)} para la clasificación y  para la regresión   {\ttfamily max\_features = log2(n\_features)}.  

Generalmente agregando   {\ttfamily max\_features  o  max\_leaf\_nodes}  puede mejorar el rendimiento del modelo; y también puede reducir drásticamente el espacio y el tiempo del ordenador requeridos para el entrenamiento y la predicción.\\


%          https://sitiobigdata.com/


\subsubsection{Árboles de regresión con potenciación del gradiente \\ (Máquinas de refuerzo de gradiente)} 

% árboles de regresión potenciados


%         https://rpubs.com/Cristina_Gil/arboles_ensemble   bueno para leer gradientes y áboles
%  computacionalmente inviable
%   https://www.datahack.es/generacion-de-lenguaje-natural/
%   https://dlegorreta.wordpress.com/tag/naive-bayes/

El árbol de regresión  con potenciación del gradiente\index{potenciación del gradiente} o potenciación del gradiente simplemente (gradient boosted\index{gradient boosted}) es otro método de conjunto que  combina múltiples árboles de decisión para crear un modelo más potente. A pesar del  nombre de  "{}regresión", estos modelos  pueden  ser usados  para  regresión y clasificación.  En contraste  con el enfoque  de bosque aleatorio, el  gradient boosting (aumento de gradiente) funciona   para construir árboles en forma serial, donde cada árbol trata de corregir los errores del    anterior. Por defecto, no hay aleatorización en árboles de regresión gradient boosted; en su lugar, se usa una fuerte pre-poda. La  potenciación del gradiente   a generalmente utilizan árboles muy superficiales, con   profundidad de uno a cinco, lo que hace que el modelo sea más pequeño en términos  de memoria  y hace predicciones más rápidas.\\
 

La idea principal detrás de la potenciación del gradiente   (gradient boosting) es combinar muchos modelos 
simples (en este contexto es conocido como aprendizaje débil), como árboles poco profundos. 
Cada árbol sólo puede proporcionar buena predicciones en parte de los datos, por lo que se 
agregan más y más árboles en forma iterativa para mejorar el rendimiento.\\

Los árboles de regresión potenciados son con frecuencia las entradas ganadoras en las competencias de aprendizaje automático, y son ampliamente utilizados en la industria. Generalmente son un poco más sensibles a la  configuración de parámetros que bosques aleatorios, pero pueden proporcionar una mejor exactitud si los parámetros  están correctamente configurados.\\



Además de la pre-poda (poda previa) y el número de árboles,  hay otro  parámetro 
 importante  del  Gradient boosted   que es la tasa de aprendizaje  ({\ttfamily {learning\_rate}}), este parámetro controla qué   tan fuerte cada árbol trata de corregir los errores de los árboles anteriores. Un  valor alto de {\ttfamily {learning\_rate}}  significa que cada árbol puede hacer correcciones más 
 fuertes, permitiendo modelos  más complejas. También se pueden  añadir más árboles al conjunto, mediante el incremeto de  {\ttfamily n\_estimators}, pero también  aumentará la complejidad del modelo, ya que el 
 modelo tiene más posibilidades de corregir errores en el conjunto de entrenamiento. \\
 
  
 Acontinuación ilustramos la técnica {\ttfamily GradientBoostingClassifier}  con un ejemplo, aplicado 
sobre la data  cáncer de mama (Breast Cancer). Por defecto, se utilizan 100
árboles con profundidad máxima 3 y una tasa de aprendizaje de 0,1:	\\
 

\begin{lstlisting}
In[72]:  from sklearn.ensemble import GradientBoostingClassifier

	X_train, X_test, y_train, y_test = train_test_split(
	     cancer.data, cancer.target, random_state=0)
	
	gbrt = GradientBoostingClassifier(random_state=0)
	gbrt.fit(X_train, y_train)
	
	print("Accuracy on training set: {:.3f}".format(
	       gbrt.score(X_train, y_train)))
	print("Accuracy on test set: {:.3f}".format(
	       gbrt.score(X_test, y_test)))

Out[72]:
    Accuracy on training set: 1.000
    Accuracy on test set: 0.958
 
 \end{lstlisting}
Como la exactitud del conjunto de entrenamiento es del 100\%, es probable que tengamos un 
sobreajuste. Para reducir el sobreajuste, podríamos aplicar una pre-poda  más fuerte limitando la profundidad máxima o bajar la tasa de aprendizaje:\\


\begin{lstlisting}
In[73]: gbrt = GradientBoostingClassifier(random_state=0, max_depth=1)
	gbrt.fit(X_train, y_train)
	print("Accuracy on training set: {:.3f}".format(gbrt.score(
	       X_train, y_train)))
	print("Accuracy on test set: {:.3f}".format(gbrt.score(
	       X_test, y_test)))

Out[73]:
	Accuracy on training set: 0.991
	Accuracy on test set: 0.972

In[74]: gbrt = GradientBoostingClassifier(random_state=0, 
                                          learning_rate=0.01)
	gbrt.fit(X_train, y_train)
	print("Accuracy on training set: {:.3f}".format(gbrt.score(X_train,
	                                                           y_train)))
	print("Accuracy on test set: {:.3f}".format(gbrt.score(X_test,
	                                                       y_test)))

Out[74]:
	
	Accuracy on training set: 0.988
	Accuracy on test set: 0.965
\end{lstlisting}

Ambos métodos para disminuir la complejidad del modelo redujeron la exactitud del conjunto de
 entrenamiento, como se esperaba. En este caso, la reducción de la profundidad máxima de los
  árboles proprocinó una mejora  significativa en el modelo, mientras que la reducción de la tasa de 
 aprendizaje solo aumentó el rendimiento de generalización ligeramente. \\

En cuanto a los otros modelos basados en árbol de decisión, aquí también podemos visualizar la 
importancia de las características, para  obtener más información sobre el modelo 
(Figura \ref{ImportGradiente}). Como se usaron  100 árboles, no es práctico la inspección de los áboles, incluso  si son todos de profundidad 1:\\

\begin{lstlisting}
In[75]: gbrt = GradientBoostingClassifier(random_state=0, max_depth=1)
	gbrt.fit(X_train, y_train)
	plot_feature_importances_cancer(gbrt)
\end{lstlisting}

\begin{figure}[h!]
	\centering
	\includegraphics[scale=0.50]{../Imagen_Machine/ImportGradiente}
	\caption{Importancias de características calculadas a partir de un clasificador gradient boosted  que se ajusta a la data  cáncer de mama}
	\label{ImportGradiente}
\end{figure}


Podemos ver que la importancia de las características en gradient boosted,  son algo similar a la importancia de las características de los bosques aleatorios, aunque el gradient boosted ignoró por completo algunas de las características.\\

Dado que tanto en potenciación del gradiente   (gradient boosting)  como en bosques aleatorios se 
obtienen buenos resultados con datos similares, un enfoque común es probar primero en  bosques 
aleatorios, que funcionan con bastante solidez, muy robusto. Si los bosques aleatorios funcionan 
bien, pero el tiempo de predicción está en juego o es importante exprimir el último porcentaje de exactitud del modelo de aprendizaje automático, el aumento del gradiente generalmente ayuda.\\
calidad superior a la media

Si se quiere aplicar el gradient boosting a un problema de gran escala, tal vez  valer la pena
revisar el paquete {\ttfamily xgboost}\index{xgboost} y su interfaz Python, que al momento de la escritura
 es más rápido  que la implementación gradient boosting  en  {\ttfamily scikit-learn}   en muchos conjuntos de datos, y  {\ttfamily xgboost} es algunas veces es más fácil de ajustar.


\subsubsection{Fortalezas, debilidades y parámetros}

Los Árboles de decisión potenciados por gradiente o árboles de regresión potenciados se encuentran entre los modelos más potentes y ampliamente utilizados para el aprendizaje supervisado. Su principal inconveniente es que requieren un ajuste cuidadoso de los parámetros y pueden llevar mucho tiempo en su entrenamiento. De manera similar   otros modelos basados en árboles, el algoritmo funciona bien sin escalar y en una mezcla de  características binarias y continuas. Al igual que con otros modelos basados en árboles, por lo general no funciona bien  sobre datas  sparse de alta dimensión.\\

%  regresión del árbol de decisión potenciado

Los parámetros principales de los modelos de gradiente potenciado  son:  el número de árboles a usar en el modelo  {\ttfamily n\_estimators},  y  la taza de aprendizaje {\ttfamily learning\_rate} que controla el  grado en que  cada árbol correge los errores de los árboles anteriores.  Estos dos parámetros están altamente interconectados, ya que una  menor taza de aprendizaje {\ttfamily learning\_rate},  significa que se necesitan más árboles para construir un modelo de  complejidad similar. 

En contraste con los bosques aleatorios en donde un valor más alto  de \: {\ttfamily n\_estimators} es siempre mejor,  en el gradient boosting  un incremento de {\ttfamily n\_estimators} conduce a un modelo más complejo,  que puede conducir al sobreajuste. Una práctica común es ajustar el {\ttfamily n\_estimators}  dependiendo del   tiempo y la capacidad de memoria, y luego probar con diferentes  {\ttfamily learning\_rates}. \\	


Otro parámetro importante es la profundidad máxima  {\ttfamily max\_depth}, o como alternativa puede ser {\ttfamily max\_leaf\_nodes},  para reducir la  complejidad de cada árbol. Por lo general el parámetro {\ttfamily max\_depth} en los  modelos  gradient boosted,  se configuran con valores no más de cinco divisiones, es decir, valores bajos. \\
 
 
%  https://scikit-learn.org/stable/modules/model_evaluation.html

\subsection{Máquinas de soporte de vectores kernelized}\label{kenrelpag}

 %  máquinas de soporte vectorial  centralizadas,kernel\\

En  esta sección se estudia un modelo de aprendizaje  automatizado  supervisado,  llamado  Máquinas de soporte de vectores kernelized  (SVMs).  En la sección  dedicada a "{}Modelos lineales para clasificación"{}  en la página \pageref{Mlineales},  se introdujo el uso de máquinas de soporte de vectores  lineal (SVM) para la clasificación. Esta técnica de SVMs   es una extensión de SVM que permite modelos más complejos que definen otras curvas y no  simplemente rectas o hiperplanos  en el espacio de entrada. Más adelante se desarrolla un ejemplo donde resulta una elipse (figura \ref{LimitDesi3D}) como límite de decisión. 

Hay máquinas de soporte de vectores  para clasificación y regresión, sin embargo, en este texto nos limitaremos a la caso de clasificación, tal como se implementó en SVC. De manera análoga, los conceptos  se aplican al vector de soporte para regresión,   y se implementa en SVR.\\


La matemática detrás de las máquinas de soporte de vectores kernelized es un poco complicada y 
está más allá el alcance de este libro. Puede encontrar los detalles en el el Capítulo 1 del Libro: \href{https://web.stanford.edu/~hastie/ElemStatLearn//}{\color{vino} The Elements of Statistical Learning}
  de  Hastie,  Tibshirani, and Friedman’s. Sin embargo, 
trataremos  de darle sentido práctico a la idea del método.



\subsubsection{Modelos lineales y características no lineales}


Como se ve en la Figura \ref{SVM}, los modelos lineales pueden ser muy limitados en espacios
baja dimensión, como líneas rectas e hiperplanos que  tienen flexibilidad limitada. Una forma 
de hacer que un  modelo  un lineal sea más flexible es agregar más características, por ejemplo, agregar
interacciones o polinomios como características de entrada.\\


Ilustremos la idea con la data  sintética que se usó para visualizar la  "{}Importancia de la característica  en los árboles"{} en la página \pageref{ImpotCaract}  (Figura \ref{EtiquetArbol2}):

\begin{lstlisting}
In[76]: X, y = make_blobs(centers=4, random_state=8)
        y = y % 2
        mglearn.discrete_scatter(X[:, 0], X[:, 1], y)
        plt.xlabel("Feature 0")
        plt.ylabel("Feature 1")
\end{lstlisting}


\begin{figure}[h!]
	\centering
	\includegraphics[scale=0.50]{../Imagen_Machine/Conjunto2D}
	\caption{Data con dos clases, en el que las clases no son linealmente separables}
	\label{Conjunto2D}
\end{figure}

Un modelo lineal para  clasificar esta data, sólo puede separar los puntos usando una línea recta, y no será
capaz de hacer un buen trabajo en este conjunto de datos (ver Figura \ref{SVMLineal}):\\


\begin{lstlisting}
In[77]: from sklearn.svm import LinearSVC
	linear_svm = LinearSVC().fit(X, y)
	
	mglearn.plots.plot_2d_separator(linear_svm, X)
	mglearn.discrete_scatter(X[:, 0], X[:, 1], y)
	plt.xlabel("Feature 0")
	plt.ylabel("Feature 1")
\end{lstlisting}


Observe lo que sucede  al expandir el conjunto de características de entrada, por ejemplo, agregando también  {\ttfamily feature1 ** 2}, el cuadrado de la segunda característica, como una nueva característica.  Ahora, en lugar de representar cada punto de datos como un punto bidimensional (feature0, feature1), lo representamos como un punto tridimensional (feature0, feature1, feature1 ** 2). Es importante notar que la elección de  esta característica  se hizo únicamente con  fines  para la ilustración, esta elección no tiene alguna  particularida importante.   Ahora, esta nueva representación se ilustrada en la figura \ref{Tridimensional}, un diagrama de dispersión tridimensional:


\begin{figure}[h!]
	\centering
	\includegraphics[scale=0.50]{../Imagen_Machine/SVMLineal}
	\caption{Límite de decisión encontrado por un SVM lineal}
	\label{SVMLineal}
\end{figure}



\begin{lstlisting}
In[78]: # Anade el cuadrado de la primera caracteristica
	X_new = np.hstack([X, X[:, 1:] ** 2])    #  apila horizontal
	
	from mpl_toolkits.mplot3d import Axes3D, axes3d
	figure = plt.figure()
	# visualize in 3D
	ax = Axes3D(figure, elev=-152, azim=-26)
	# plot first all the points with y == 0, then all with y == 1
	mask = y == 0
	ax.scatter(X_new[mask, 0], X_new[mask, 1], X_new[mask, 2], c='b',
	           cmap=mglearn.cm2, s=60)
	ax.scatter(X_new[~mask, 0], X_new[~mask, 1], X_new[~mask, 2], 
	            c='r', marker='^', cmap=mglearn.cm2, s=60)
	ax.set_xlabel("feature0")
	ax.set_ylabel("feature1")
	ax.set_zlabel("feature1 ** 2")
\end{lstlisting}



En la nueva representación de la data, ahora es posible separar las dos clases usando un modelo 
lineal, un plano en tres dimensiones. Podemos confirmar esto ajustando un modelo lineal a la data aumentada (ver Figura \ref{HiperPlano}): 

\begin{lstlisting}
In[79]: linear_svm_3d = LinearSVC().fit(X_new, y)
	coef      =  linear_svm_3d.coef_.ravel() 
	intercept =  linear_svm_3d.intercept_
	# show linear decision boundary
	figure = plt.figure()
	ax = Axes3D(figure, elev=-152, azim=-26)
	xx = np.linspace(X_new[:, 0].min() - 2, X_new[:, 0].max() + 2, 50)
	yy = np.linspace(X_new[:, 1].min() - 2, X_new[:, 1].max() + 2, 50)
	
	XX, YY = np.meshgrid(xx, yy)
	ZZ = (coef[0] * XX + coef[1] * YY + intercept) / -coef[2]
	ax.plot_surface(XX, YY, ZZ, rstride=8, cstride=8, alpha=0.3)
	ax.scatter(X_new[mask, 0], X_new[mask, 1], X_new[mask, 2], 
	           c='b', cmap=mglearn.cm2, s=60)
	ax.scatter(X_new[~mask, 0], X_new[~mask, 1], X_new[~mask, 2],
	           c='r', marker='^', cmap=mglearn.cm2, s=60)

	ax.set_xlabel("feature0")
	ax.set_ylabel("feature1")
	ax.set_zlabel("feature0 ** 2")
\end{lstlisting}





\begin{figure}[h!]
	\centering
	\includegraphics[scale=0.50]{../Imagen_Machine/Tridimensional}
	\caption{Expansión del conjunto de datos de la Figura \ref{SVMLineal}, 
		creada añadiendo una tercera característica derivada de la característica 1.}
	\label{Tridimensional}
\end{figure}







\begin{figure}[h!]
	\centering
	\includegraphics[scale=0.50]{../Imagen_Machine/HiperPlano}
	\caption{Límite de decisión encontrado por un SVM lineal en el conjunto de datos tridimensionales expandidos.}
	\label{HiperPlano}
\end{figure}
 
 
  Luego de este cambio, en el modelo lineal SVM  en función de las dos características originales, el límite de decisión  ya no es lineal, ya no se comporta como una linea recta, sino más bien una elipse, como se puede ver en en el gráfico que se muestra en la Figura \ref{LimitDesi3D}:\\

\begin{lstlisting}
In[80]: ZZ = YY ** 2
	dec = linear_svm_3d.decision_function(np.c_[XX.ravel(),
	                                       YY.ravel(), ZZ.ravel()])
	plt.contourf(XX, YY, dec.reshape(XX.shape),
	             levels=[dec.min(), 0, dec.max()], cmap=mglearn.cm2, 
	             alpha=0.5)
	mglearn.discrete_scatter(X[:, 0], X[:, 1], y)
	plt.xlabel("Feature 0")
	plt.ylabel("Feature 1")
\end{lstlisting}


\begin{figure}[h!]
	\centering
	\includegraphics[scale=0.50]{../Imagen_Machine/LimitDesi3D}
	\caption{El límite de decisión de la Figura \ref{HiperPlano} en función de las dos características originales}
	\label{LimitDesi3D}
\end{figure}



\subsubsection{El truco del kernel}

La idea aquí es  añadir características no lineales a la data  para la representación,  
y de  esta forma se puede lograr que  los modelos lineales sean  más potentes.  Sin embargo, a menudo no 
sabemos qué características añadir, y la adición de muchas características (como todas las
 interacciones posibles en un espacio de características de 100 dimensiones) puede hacer  el costo de computación muy alto. Afortunadamente, hay un truco matemático inteligente que le permite el  aprendizaje a un clasificador en un espacio de dimensiones más altas,  sin ejecutar realmente los nuevos calculos. Este truco conoce como el truco del kernel; y trabaja calculando directamente la distancia (más precisamente, los productos escalares) de los puntos de datos para la representación de las características expandidas, sin calcular realmente la expansión.\\
 
 
 % Afortunadamente, hay un truco matemático inteligente que nos permite aprender un clasificador en un espacio de dimensiones superiores sin realmente computar la nueva, posiblemente muy grande representación
  
%  Luckily, there is   a clever mathematical trick that allows us to learn a classifier in a higher-dimensional   space without actually computing the new, possibly very large representation
  
  
 Hay dos maneras de mapear sus datos en un espacio de dimensiones superiores que se utilizan comúnmete con máquinas de soporte de vectores: una se refiere al kernel polinomial, que calcula todos los polinomios posibles hasta un cierto grado de las características originales (como {\ttfamily feature1**2 * feature2 ** 5}); y el otro es el kernel de la función de base radial   (radial basis function kernel,  RBF), también conocido como  el kernel Gaussiano.  El  kernel Gaussiano es un poco más difícil de explicar, ya que corresponde a un espacio de características de dimensiones infinitas. Una manera de explicar el  kernel Gaussian es considerar todos los polinomios posibles de todos los grados, pero la importancia de las características disminuyen para grados muy altos (expansión de Taylor).\\


En la práctica, los detalles matemáticos de fondo del  kernel  SVM no son tan importantes al análisis, sin embargo,  una SVM con un kernel RBF toma una decisión que se puede resumir fácilmente; para comprender esta idea explicamos en la  siguiente sección. 

\subsubsection{Comprendiendo el proceso de SVMs}

Durante el entrenamiento, el SVM aprende la importancia de cada uno de los puntos de datos del entrenamiento para representar el límite de decisión entre las dos clases; generalmente  sólo un subconjunto del entrenamiento es importartante para definir el límite de decisión: y son los datos que se encuentran en la frontera entre las clases. Estos se llaman soporte de vectores, que dan su nombre a la máquina de soporte de vectores.


 Para hacer una predicción de un nuevo punto, la distancia a cada uno de los  soporte de vectores es medida. Se toma una decisión de clasificación basada en las distancias al vector de soporte, y a la importancia de los soporte de vectores que fuerón aprendiendo durante el entrenamiento (almacenados en el atributo  {\ttfamily dual\_coef\_} de SVC).\\


La distancia entre puntos de la data es medido por el kernel Gaussiano: 


\begin{equation}\label{gausss}
 k_{rbf}(x_{1}, x_{2}) = exp (\gamma \|x_{1} - x_{2}\|^{2} ). 
 \end{equation} Aquí, $x_{1}$  y  $x_{2} $  son puntos de la data, $\| x_{1} - x_{2} \| $ denota distancia Euclidean, y $\gamma$ (gamma) es un parámetro que controla el ancho del kernel Gaussian.\\


 La Figura \ref{KernelGausiano1} muestra el resultado del entrenamiento de una máquina de vector de soporte en un conjunto de datos bidimensional con dos clases. El límite de decisión se muestra como la curva  en negro, y los vectores de suporte son puntos más grandes con el contorno ampliamante marcado. El siguiente código genera el gráfico de entrenanar  un SVM en la dataset de falsificación de datos (forge dataset):

\begin{lstlisting}
In[81]: from sklearn.svm import SVC
	X, y = mglearn.tools.make_handcrafted_dataset()
	svm = SVC(kernel='rbf', C=10, gamma=0.1).fit(X, y)
	mglearn.plots.plot_2d_separator(svm, X, eps=.5)
	mglearn.discrete_scatter(X[:, 0], X[:, 1], y)
	# plot support vectors
	sv = svm.support_vectors_
	# class labels of support vectors are given by the sign 
	# of the dual coefficients
	sv_labels = svm.dual_coef_.ravel() > 0
	mglearn.discrete_scatter(sv[:, 0], sv[:, 1], sv_labels, 
	                         s=15, markeredgewidth=3)
	plt.xlabel("Feature 0")
	plt.ylabel("Feature 1")	
\end{lstlisting}


\begin{figure}[h!]
	\centering
	\includegraphics[scale=0.50]{../Imagen_Machine/KernelGausiano1}
	\caption{Límite de decisión y soporte de vectores encontrados por un SVM con el kernel Gaussiano}
	\label{KernelGausiano1}
\end{figure}

En este caso, el SVM produce un límite de decisión  más suave y no lineal  (no es una línea recta). 
Aquí se ajustaron dos parámetros: el  $C = 10 $ y el  gamma = 0.1.  A  continuacón
 se trata con más detalle este ajuste.


\subsubsection{Ajuste de los parámetros de SVM}

El parámetro gamma es el que se muestra en la fórmula \ref{gausss} dada en la sección anterior, que controla el ancho del kernel Gaussiano; determina la escala de lo que significa para los puntos estar uno al lado del otro (estar cerca uno del otro).  Y  $C$  es un parámetro de regularización, similar al utilizado en los modelos lineales, y límita la importancia de cada punto  {\ttfamily dual\_coef\_}.\\

En la Figura \ref{gama2} podemos ver lo que sucede cuando variamos estos parámetros:\\

\begin{lstlisting}  % [frame=single]	

In[82]: fig, axes = plt.subplots(3, 3, figsize=(15, 10))
   
        for ax, C in zip(axes, [-1, 0, 3]):
            for a, gamma in zip(ax, range(-1, 2)):
            mglearn.plots.plot_svm(log_C=C, log_gamma=gamma, ax=a)
           
        axes[0, 0].legend(["class 0", "class 1", "sv class 0", "sv class 1"],
                           ncol=4, loc=(.9, 1.2))
  
\end{lstlisting}

 
 De izquierda a derecha, aumentamos el valor del parámetro gamma de 0,1 a 10. Un pequeño gamma significa un gran radio para el kernel Gaussiano, lo que significa que muchos puntos se consideran cercanos.
Esto se refleja en límites de decisión muy suaves a la izquierda, y límites que se centran más en puntos individuales más a la derecha.  Un valor bajo de gamma significa que el límite de decisión variará lentamente, lo que produce un modelo de baja complejidad, mientras que un valor alto de gamma produce un modelo más complejo. 
\begin{figure}[h!]
	\centering
	\includegraphics[scale=0.75]{../Imagen_Machine/gama2}
	\caption{Los límites de decisión y los vectores del soporte para los  diferentes  ajustes  de los  parámetros $C $ y  gamma}
	\label{gama2}
\end{figure} Y de arriba a abajo, aumentamos el parámetro C de 0.1 a 1000. Al igual que con los modelos lineales, un pequeño C significa un modelo muy restringido, donde cada punto de datos sólo puede tener una influencia muy limitada. Se puede ver que en la parte superior izquierda el límite de decisión parece casi lineal, con  puntos mal clasificados con poca o ninguna influencia sobre la línea. Incrementando $C$ como se muestra en la parte inferior derecha, permite que esos puntos tengan una mayor influencia sobre el modelo y hacen  del límite de decisión una curva suave  para clasificarlos correctamente.\\


Ahora aplicamos la técnica de SVM  con el kernel  gaussiano  al conjunto de datos de Cáncer de Mama (Breast Cancer). Con los parámetros por defecto,  $C=1 $  y gamma $= 1/n${\ttfamily \_features}:


\begin{lstlisting}
In[83]: X_train, X_test, y_train, y_test = train_test_split(
	    cancer.data, cancer.target, random_state=0)
	    
	svc = SVC()
	svc.fit(X_train, y_train)
	
	print("Accuracy on training set: {:.2f}".format(
	       svc.score(X_train, y_train)))
	print("Accuracy on test set: {:.2f}".format(
	       svc.score(X_test, y_test)))

Out[83]:

	Accuracy on training set: 1.00
	Accuracy on test set: 0.63

\end{lstlisting}


El modelo resultó con un puntaje de exactitud perfecta sobre el conjunto de entrenamiento, lo que indica un sobreajuste muy alto, mientras que la exactitud en el conjunto de prueba es de 63\%. Aunque los SVMs generalmente  funcionan bastante bien, estos son muy sensibles a la configuración de los parámetros y a la escala de los datos. En particular, requieren que todas las características varíen en una escala similar. Veamos los valores mínimos y máximos para cada característica, graficada en log-space (Figura \ref{RangosLog}):

\begin{lstlisting}
In[84]: plt.plot(X_train.min(axis=0), 'o', label="min")
	plt.plot(X_train.max(axis=0), '^', label="max")
	plt.legend(loc=4)
	plt.xlabel("Feature index")
	plt.ylabel("Feature magnitude")
	plt.yscale("log")

\end{lstlisting}


\begin{figure}[h!]
	\centering
	\includegraphics[scale=0.50]{../Imagen_Machine/RangosLog}
	\caption{Rangos de características para el conjunto de datos de cáncer de mama (el eje $y$ en escala logaritmica)}
	\label{RangosLog}
\end{figure}


A partir de esta gráfica, podemos determinar que las características en la data Cáncer de Mama tienen órdenes de magnitud completamente diferentes.  Esto puede ser un problema para otros modelos (como los modelos lineales), pero tiene efectos devastadores para el kernel SVM. Examinemos algunas maneras de tratar  este asunto.


\subsubsection{Preprocesamiento de datos para SVMs}


Una forma de resolver este problema es volver a reescalar  cada característica para que estén aproximadamente en la misma escala.  Un método común de re-escalar  para kernel SVMs, conciste en escalar los datos de tal manera que todas las características estén entre 0 y 1.  En el Capítulo \ref{Cap3} veremos cómo hacerlo usando el método de preprocesamiento Minmaxscaler, donde daremos más detalles. Por ahora, lo haremos esto "{}a mano":

\begin{lstlisting}  

In[85]: # compute the minimum value per feature on the training set
        min_on_training = X_train.min(axis=0)
        # compute the range of each feature (max - min) on the training set
        range_on_training = (X_train - min_on_training).max(axis=0)
   
        # subtract the min, and divide by range
        # afterward, min=0 and max=1 for each feature
        X_train_scaled = (X_train - min_on_training) / range_on_training
        print("Minimum for each feature\n{}".format(
               X_train_scaled.min(axis=0)))
        print("Maximum for each feature\n {}".format(
               X_train_scaled.max(axis=0)))

Out[85]:

    Minimum for each feature
    [ 0.  0.  0.  0.  0.  0.  0.  0.  0.  0.  0.  0.  0.  0.  0.  0.  0.  0.
      0.  0.  0.  0.  0.  0.  0.  0.  0.  0.  0.  0.]
    Maximum for each feature
    [ 1.  1.  1.  1.  1.  1.  1.  1.  1.  1.  1.  1.  1.  1.  1.  1.  1.  1.
      1.  1.  1.  1.  1.  1.  1.  1.  1.  1.  1.  1.]


In[86]: # use THE SAME transformation on the test set,
        # using min and range of the training set (see Chapter 3 for details)
        X_test_scaled = (X_test - min_on_training) / range_on_training
   
In[87]: svc = SVC()
        svc.fit(X_train_scaled, y_train)
   
        print("Accuracy on training set: {:.3f}".format(
               svc.score(X_train_scaled, y_train)))
        print("Accuracy on test set: {:.3f}".format(
               svc.score(X_test_scaled, y_test)))
   
Out[87]:
  
    Accuracy on training set: 0.948
    Accuracy on test set: 0.951


\end{lstlisting}



¡Escalando los datos hace una gran diferencia! Ahora nos encontramos en un régimen de subajuste, 
en el que el rendimiento del entrenamiento y el conjunto de prueba son muy similares, pero menos próximos al 100 \%. A partir de aquí, podemos tratar de aumentar ya sea C o gamma para adaptarse a un modelo más complejo. Por ejemplo:\\


\begin{lstlisting}
In[88]: svc = SVC(C=1000)
	svc.fit(X_train_scaled, y_train)
	
	print("Accuracy on training set: {:.3f}".format(
	       svc.score(X_train_scaled, y_train)))
	print("Accuracy on test set: {:.3f}".format(
	       svc.score(X_test_scaled, y_test)))

Out[88]:
	Accuracy on training set: 0.988
	Accuracy on test set: 0.972
\end{lstlisting}


\noindent Aquí, el aumento de C nos permite mejorar significativamente el modelo, lo que resulta en un 97.2\%
exactitud.


\subsubsection{Fortalezas, debilidades y parámetros}

Las máquinas de soporte de vectores Kernelized\index{soporte de vectores Kernelized} son modelos potentes y funcionan bien en una variedad
 de conjuntos de datos. Las SVMs permiten límites de decisión complejos, incluso si los datos solo tienen pocas características. Funcionan bien en datos de baja y alta dimensión (es decir, con pocas y muchas características), pero puede generar cierta dificultad con el número de muestras; correr un SVM con datos de tamaño hasta de 10.000 muestras podría funcionar bien, pero trabajar con conjuntos de datos de tamaño 100.000 o más, pueden convertirse en un desafío en términos de tiempo de ejecución y uso de memoria.\\


Otro inconveniente de los SVMs es que requieren un cuidadoso tratamiento previo (preprocesamiento)  de los datos y ajuste de los parámetros. Esta es la razón por la que hoy en día la mayoría de la gente utiliza modelos basados en árboles, como bosques aleatorios o potenciación de gradientes (que requieren poco o ningún proceso previo) en muchas aplicaciones.   Además, los modelos SVM son difíciles de inspeccionar; puede ser difícil entender por qué  hizo una predicción en particular, y podría ser difícil explicar el modelo a los inexpertos.\\

Aún así, podría valer la pena probar SVMs, especialmente si todas sus características estan  representadas
en unidades de medidas similares (por ejemplo, todas son intensidades de píxeles) y están en 
escala  similar.\\

Los parámetros importantes en el  kernel SVMs  son el parámetro de regularización C, la
elección del kernel, y los parámetros específicos del kernel.  Aunque principalmente se enfoca en el  kernel  kernel Gaussiano (RBF), en {\ttfamily scikit-learn}  hay otras opciones disponibles.  El  kernel RBF  tiene un sólo   que es el inverso del ancho del  kernel Gaussiano, el  parámetro el gamma.   Los  parámetro gamma y C controlan la complejidad del modelo, con valores grandes en cualquier de los dos, resulta un modelo más complejo. Por lo tanto, deben configurarse bien los dos los parámetros que  generalmente resulta altamente correlacionados; es decir,  deben cofigurarse juntos ($C$ y gamma).


\subsection{Redes Neuronales (Aprendizaje Profundo)}

Una familia de algoritmos conocidos como redes neuronales\index{redes neuronale}, se han visto recientemente en un renacimiento bajo el nombre de \: "{}Aprendizaje Profundo" \:(deep learnig). El aprendizaje profundo ofrece una  gran promesa en muchas aplicaciones de aprendizaje automático,  sus algoritmos   se adaptan generalmente a un caso de uso muy específico.  Aquí sólo discutiremos algunos métodos relativamente simples, como  los perceptores multicapa o multilayer para la clasificación y la regresión, que pueden servir como punto de partida para  comenzar a involucrarse con métodos de aprendizaje profundo, con mayor nivel.   Los perceptores Multilayer (MLPs)   también se conocen como \:  "{}Redes Neuronales de feed-forward"  {}\:(vanilla), o sólo redes neuronales. %    que se las alimentan adelante redes de nervios

% El perceptrón es una red neuronal artificial (RNA) formada por múltiples capas, de tal manera que tiene capacidad para resolver problemas que no son linealmente separables, lo cual es la principal limitación del perceptrón (también llamado perceptrón simple).


\subsubsection{El modelo de red neuronal}

Los perceptores Multilayer (MLPs)  pueden verse como generalización de modelos lineales que realizan múltiples etapas de procesamiento para llegar a una decisión.\\

Recuerde que la predicción para un regresor lineal esta dada como:

  \[  \hat{y} = w[0] * x[0] + w[1] * x[1] + ... + w[p] * x[p] + b \]

En palabras simple, $\hat{y}$  es una suma ponderada de las características de entrada $x[0]$ a $x[p]$, 
ponderada por los coeficientes aprendidos $w[0]$ a $w[p]$ .\\

En la Figura \ref{RegresLogis2-44} se ilustrar  gráficamente esta idea:

\begin{lstlisting}  % [frame=single]	

In[89]: display(mglearn.plots.plot_logistic_regression_graph())

\end{lstlisting}


\begin{figure}[h!]
	\centering
	\includegraphics[scale=0.50]{../Imagen_Machine/RegresLogis2-44}
	\caption{Visualización de la regresión logística, donde las características de entrada y las predicciones  se muestran como nodos, y los coeficientes son conexiones entre los nodos}
	\label{RegresLogis2-44}
\end{figure}


Aquí, cada nodo de la izquierda representa una característica de entrada, las líneas de conexión representan los coeficientes aprendidos, y el nodo de la derecha representa la salida, que es una suma ponderada de las entradas.\\ 


En un MLP, este proceso de cálculo de sumas ponderadas se repite varias veces, primero computando unidades ocultas que representan una etapa de procesamiento intermedia, y que luego  se combinan usando  nuevamente  sumas ponderadas para obtener el resultado final (Figura \ref{Layeoculas}):\\

\begin{lstlisting}  % [frame=single]	

In[90]: display(mglearn.plots.plot_single_hidden_layer_graph())

\end{lstlisting}

\begin{figure}[h!]
	\centering
	\includegraphics[scale=0.50]{../Imagen_Machine/Layeocultas}
	\caption{Ilustración de perceptrón multicapa con una única capa oculta}
	\label{Layeoculas}
\end{figure} Este modelo tiene muchos más coeficientes (también llamados pesos o ponderaciones) para aprender,  hay uno
entre cada entrada y cada unidad oculta (que forman la capa oculta), y uno entre cada unidad en la capa oculta y la salida.\\


Calcular una serie de sumas ponderadas es matemáticamente lo mismo que calcular sólo
una suma ponderada;    para que este modelo sea realmente más poderoso que un modelo lineal,
necesitamos un truco extra. Después de calcular una suma ponderada para cada unidad oculta, una función no lineal se aplica al resultado, generalmente la  rectificadora no lineal  (también conocido como unidad  rectificada  lineal o relu) o las tangetes hiperbólicos (tanh).   El resultado de esta función es entonces usado en la suma ponderada que calcula la salida  $ \hat{y}$.   Las dos  funciones se visualizan en la Figura \ref{Tangete}.    El relu corta valores por debajo de cero, mientras que la tanh se satura  hasta –1,  para valores de entrada bajos;  y    para valores de entrada altos el  $+1$. Otras funciones no lineales permite a la red neuronal aprender funciones mucho  más  complicadas,  de lo que  podría hacer un modelo lineal:\\


\begin{lstlisting}   
In[91]: line = np.linspace(-3, 3, 100)
        plt.plot(line, np.tanh(line), label="tanh")
        plt.plot(line, np.maximum(line, 0), label="relu")   
        plt.legend(loc="best")
        plt.xlabel("x")
        plt.ylabel("relu(x), tanh(x)")
	
\end{lstlisting}


\begin{figure}[h!]
	\centering
	\includegraphics[scale=0.50]{../Imagen_Machine/Tangete}
	\caption{La función de activación tangente hiperbólica y la función de activación lineal 
		rectificada.}
	\label{Tangete}
\end{figure}


Para  red neuronal representada en la Figura \ref{Layeoculas}, la fórmula completa para 
la cálculo $\hat{y}$  en el caso de regresión sería (cuando se utiliza  tanh, no linealidad):\\
{\scriptsize{
\[
\begin{array}{l c l}
	h[0] & = &  tanh (w [0, 0] * x [0] + w [1, 0] * x [1] + w [2, 0] * x [2] + w [3, 0] * x [3]) \\
	h[1] & = &  tanh (w [0, 0] * x [0] + w [1, 0] * x [1] + w [2, 0] * x [2] + w [3, 0] * x [3]) \\
	h[2] & = &  tanh (w [0, 0] * x [0] + w [1, 0] * x [1] + w [2, 0] * x [2] + w [3, 0] * x [3])\\
	\hat{y}& =& v [0] * h [0] + v [1] * h [1] + v [2] * h [2] 
\end{array} \\
\]
}}  Aquí, $w$ son los pesos entre la entrada $x$ y la capa oculta   $h$ (hidden layer), y  $ v $ son los
pesos entre la capa oculta $h$ y la salida $\hat{y}$. Los pesos $v$ y $w$ se aprenden
a partir de los datos, $x$ son las características de entrada, $\hat{y}$ es la salida calculada  y la $h$ son cálculos  intermedios. \\ 

Un parámetro importante que debe configurar el usuario es el número
de nodos en la capa oculta. Este puede ser tan pequeño como 10 para conjuntos de datos  muy pequeños o simples, y tan grande como 10.000 para data  muy complejas o muy grandes. También es posible agregar más
capas ocultas, como se muestra en la Figura \ref{Perceptron}:


\begin{lstlisting}  % [frame=single]	

In[92]: mglearn.plots.plot_two_hidden_layer_graph()

\end{lstlisting}


\begin{figure}[h!]
	\centering
	\includegraphics[scale=0.50]{../Imagen_Machine/Perceptron}
	\caption{Un perceptrón con dos capas ocultas, multicapa.}
	\label{Perceptron}
\end{figure}


Tener  redes neuronales grandes  compuestas de muchas de estas capas de cálculo computacinal es lo que inspiró el término "{}aprendizaje profundo".

\subsubsection{Ajuste de redes neuronales}

 Analicemos  el funcionamiento del perceptor multilayer MLP, aplicando el  {\ttfamily Mlpclassifier} sobre la dataset {\ttfamily two\_moons}  que se  usó anteriormente en este capítulo. Los resultados se muestran en la figura \ref{100ocultos}:\\


\begin{lstlisting}  % [frame=single]
	
In[93]: from sklearn.neural_network import MLPClassifier
	from sklearn.datasets import make_moons
   
   	 X, y = make_moons(n_samples=100, noise=0.25, random_state=3)
   
   	 X_train, X_test, y_train, y_test = train_test_split(X, y, 
   	                                                     stratify=y,
                                                         random_state=42)	
                                                 							   
   	 mlp = MLPClassifier(solve='lbfgs',
   	                      random_state=0).fit(X_train, y_train)
   	 mglearn.plots.plot_2d_separator(mlp, X_train, fill=True, alpha=.3)
   	 mglearn.discrete_scatter(X_train[:, 0], X_train[:, 1], y_train)
   	 plt.xlabel("Feature 0")
   	 plt.ylabel("Feature 1")
   
\end{lstlisting}



\begin{figure}[h!]
	\centering
	\includegraphics[scale=0.50]{../Imagen_Machine/100ocultos}
	\caption{Límite de decisión aprendido por una red neural con 100 unidades ocultas en el conjunto de datos \ttfamily{two\_moons}}
	\label{100ocultos}
\end{figure}

Como se puede ver, la red neuronal aprendió un límite de decisión relativamente suave y no lineal. Esto se logró empleando el  algoritmo {\ttfamily algorithm =  '\!lbfgs'}, que discurá más adelante.\\

Por defecto, el MLP usa 100 nodos ocultos, que en este caso es demasiado  para esta  pequeña  dataset.
Podemos reducir el número (el cual reduce la complejidad del modelo) y todavía podemos tener un buen resultado ( ver figura \ref{10nodos}):

lbfgs
\begin{lstlisting}	
In[94]:	mlp = MLPClassifier(solve='lbfgs', random_state=0,
                            hidden_layer_sizes=[10])
        mlp.fit(X_train, y_train)
        mglearn.plots.plot_2d_separator(mlp, X_train, fill=True, 
                                        alpha=.3)
        mglearn.discrete_scatter(X_train[:, 0], X_train[:, 1], y_train)
        plt.xlabel("Feature 0 ")
        plt.ylabel("Feature 1 ")

\end{lstlisting}	


\begin{figure}[h!]
	\centering
	\includegraphics[scale=0.50]{../Imagen_Machine/10nodos}
	\caption{Límite de decisión aprendido por una red neural con 10 unidades ocultas en la dataset  \ttfamily{two\_moons}}
	\label{10nodos}
\end{figure}


Con sólo 10 unidades ocultas, el límite de decisión parece algo más irregular.
La no linealidad por defecto es relu,  como se ve en la figura \ref{Tangete}.
Con una simple capa oculta, esto significa que la función  de decisión se compone de 10 segmentos de línea recta.   Si queremos un límite  de decisión más suave, podríamos añadir más unidades ocultas (Figura \ref{10nodos}), y agregando una segunda capa (Figura \ref{2nodos}), o usar el rectificador no lineal  tanh (Figura \ref{2nodosTagh}):
 

\begin{lstlisting}	

In[95]: # usando dos capas ocultas, con diez unidades cada una
        mlp = MLPClassifier(solve ='lbfgs', random_state=0,
   	                        hidden_layer_sizes=[10, 10])
        mlp.fit(X_train, y_train)
        mglearn.plots.plot_2d_separator(mlp, X_train, fill=True,
	                                    alpha=.3)
        mglearn.discrete_scatter(X_train[:, 0], X_train[:, 1], y_train)
        plt.xlabel("Feature 0")
        plt.ylabel("Feature 1")

\end{lstlisting}	


\begin{figure}[h!]
	\centering
	\includegraphics[scale=0.50]{../Imagen_Machine/2nodos}
	\caption{Límite de decisión aprendido con 2 capas ocultas con 10 unidades ocultas cada una, y con la activación de función de rectificación.}
	\label{2nodos}
\end{figure}

\begin{lstlisting}	

In[96]: # usando dos capas ocultas, con diez unidades cada una, ahora con 
        #  tanh (no lineal)       
        mlp = MLPClassifier(solve='lbfgs', activation='tanh',
                            random_state=0, hidden_layer_sizes=[10, 10])
        mlp.fit(X_train, y_train)
        mglearn.plots.plot_2d_separator(mlp, X_train, fill=True, alpha=.3)
        mglearn.discrete_scatter(X_train[:, 0], X_train[:, 1], y_train)
        plt.xlabel("Feature 0")
        plt.ylabel("Feature 1")

\end{lstlisting}

\begin{figure}[h!]
	\centering
	\includegraphics[scale=0.50]{../Imagen_Machine/2nodosTagh}
	\caption{Límite de decisión aprendido con 2 capas ocultas con 10 unidades ocultas cada una, con función de activación de tanh}
	\label{2nodosTagh}
\end{figure}

Por último, también podemos controlar la complejidad del modelo de  red neuronal  mediante el uso de una penalidad de tipo  $L2,$ para reducir los pesos hacia cero, como lo hicimos en la regresión Ridge y la clasificación lineal. El parámetro  del {\ttfamily Mlpclassifier} es  {\ttfamily alfa} (como en el modelo de regresión lineal),   por defecto está cofigurado con un valor muy bajo (poca regularización). La figura \ref{DiferenteHidden} muestra el efecto de diferentes valores de  {\ttfamily alfa} en el conjunto de datos {\ttfamily two\_moons}, utilizando dos capas ocultas de 10 o 100 unidades cada una:

\begin{lstlisting}  

In[97]: fig, axes = plt.subplots(2, 4, figsize=(20, 8))
        for axx, n_hidden_nodes in zip(axes, [10, 100]):
            for ax, alpha in zip(axx, [0.0001, 0.01, 0.1, 1]):
                 mlp = MLPClassifier(solve='lbfgs', random_state=0,
                                     hidden_layer_sizes=[n_hidden_nodes,
                                     n_hidden_nodes], alpha=alpha)
                 mlp.fit(X_train, y_train)
                 mglearn.plots.plot_2d_separator(mlp, X_train, fill=True, 
                                                 alpha=.3, ax=ax)
                 mglearn.discrete_scatter(X_train[:, 0], X_train[:, 1], 
                                          y_train, ax=ax)
                 ax.set_title("n_hidden=[{}, {}]\nalpha={:.4f}".format(
                               n_hidden_nodes, n_hidden_nodes, alpha))

\end{lstlisting}


\begin{figure}[h!]
	\centering
	\includegraphics[scale=0.50]{../Imagen_Machine/DiferenteHidden}
	\caption{Funciones de decisión para diferentes números de unidades ocultas y diferentes ajustes del parámetro alfa }
	\label{DiferenteHidden}
\end{figure}


Hay muchas formas para controlar la complejidad de una red neuronal: mediante el número de capas ocultas, por el número de unidades en cada capa oculta, y la regularización {\ttfamily alfa}.  Realmente hay  más, pero no serán tratamos aquí. \\

Una propiedad importante en redes neuronales es que sus pesos están colocados al azar
antes comenzar  el aprendizaje, y esta  aleatoriedad inicial afecta el modelo  aprendido. Esto significa  que  aún cuando usen exactamente los mismos parámetros, se pueden  obtener modelos muy diferentes al usar  semillas diferentes. Si las redes son grandes, y su complejidad está seleccionada  correctamente, ésto no debería afectar mucho la exactitud, pero esto es importante tenerlo  pendiente (en particular para redes más pequeñas). La figura \ref{RandonDif}  muestra varios modelos, todos aprendiendo con la misma configuración  de parámetros (efecto de semilla aleatoria).


\begin{lstlisting}  % [frame=single]	

In[98]: fig, axes = plt.subplots(2, 4, figsize=(20, 8))
        for i, ax in enumerate(axes.ravel()):
            mlp = MLPClassifier(algorithm='l-bfgs', random_state=i,
                                 hidden_layer_sizes=[100, 100])
            mlp.fit(X_train, y_train)
            mglearn.plots.plot_2d_separator(mlp, X_train, fill=True,
                                            alpha=.3, ax=ax)
            mglearn.discrete_scatter(X_train[:, 0], X_train[:, 1], 
                                      y_train, ax=ax)

\end{lstlisting}


\begin{figure}[h!]
	\centering
	\includegraphics[scale=0.50]{../Imagen_Machine/RandonDif}
	\caption{Funciones de decisión aprendidas con los mismos parámetros, pero con diferente iniciación aleatorias}
	\label{RandonDif}
\end{figure}



Ahora apliquemos el  {\ttfamily  Mlpclassifier} a la dataset  cáncer de mama (Breast Cancer), para obtener una mejor comprensión de las redes neuronales en datos del mundo real. y comencemos  con los parámetros 
por defecto: \\

\begin{lstlisting}  % [frame=single]	

In[99]: print("Cancer data per-feature maxima:\n{}".format(
                cancer.data.max(axis=0)))

Out[99]:

    Cancer data per-feature maxima:
    [ 28.110   39.280   188.500  2501.000  0.163   0.345    0.427
       0.201    0.304     0.097     2.873  4.885  21.980  542.200
       0.031    0.135     0.396     0.053  0.079   0.030   36.040
      49.540  251.200  4254.000     0.223  1.058   1.252    0.291
       0.664    0.207]
      
In[100]:  X_train, X_test, y_train, y_test = train_test_split(
               cancer.data, cancer.target, random_state=0)
       
          mlp = MLPClassifier(random_state=42)
          mlp.fit(X_train, y_train)
   
          print("Accuracy on training set: {:.2f}".format(
                 mlp.score(X_train, y_train)))
          print("Accuracy on test set: {:.2f}".format(
                 mlp.score(X_test, y_test)))

Out[100]:
   
   Accuracy on training set: 0.92
   Accuracy on test set: 0.90


\end{lstlisting}


La exactitud  del MLP es muy  bueno, pero no tan bueno como los otros modelos. Así como en el anterior ejemplo de SVC, esto es probablemente debido a la escala de los datos. Las redes neuronales también esperan que todas las características de entradas  varien en forma similar, e idealmente que tengan media cero, y varianza uno.

 Así que debemos re-escalar los datos a fin de que cumpla con estos requisitos. En en próximo ejemplo, se hace el re-escalado de datos a mano, pero en el Capítulo \ref{Cap3} introduciremos  el  StandardScaler para hacerlo automáticamente:\\

\begin{lstlisting} 
In[101]: # compute the mean value per feature on the training set
         mean_on_train = X_train.mean(axis=0)
         # compute the standard deviation of each feature on the training set
         std_on_train = X_train.std(axis=0)
   
         # subtract the mean, and scale by inverse standard deviation
         # afterward, mean=0 and std=1
         X_train_scaled = (X_train - mean_on_train) / std_on_train
         # use THE SAME transformation (using training mean and std)
         # on the test set
         X_test_scaled = (X_test - mean_on_train) / std_on_train
   
         mlp = MLPClassifier(random_state=0)
         mlp.fit(X_train_scaled, y_train)
   
         print("Accuracy on training set: {:.3f}".format(
                mlp.score(X_train_scaled, y_train)))
         print("Accuracy on test set: {:.3f}".format(
                mlp.score(X_test_scaled, y_test)))
  
Out[101]:
   
    Accuracy on training set: 0.991
    Accuracy on test set: 0.965

    ConvergenceWarning:
        Stochastic Optimizer: Maximum iterations reached and the
        optimization hasn't converged yet.


\end{lstlisting}

Los resultados han mejorado mucho después de re-escalar, y bastante competitiva. Sin embargo, Tenemos un 
advertencia del modelo, que nos dice que el máximo número de iteraciones ha sido alcanzado. Esto es parte del algoritmo  {\ttfamily adam }  para el aprendizaje del modelo, y nos dice que deberíamos aumentar el número de iteraciones:


\begin{lstlisting}  

In[102]: mlp = MLPClassifier(max_iter=1000, random_state=0)
         mlp.fit(X_train_scaled, y_train)

         print("Accuracy on training set: {:.3f}".format(
                mlp.score(X_train_scaled, y_train)))
         print("Accuracy on test set: {:.3f}".format(
	            mlp.score(X_test_scaled, y_test)))

Out[102]:

    Accuracy on training set: 0.995
    Accuracy on test set: 0.965
\end{lstlisting}


El incremento del número de iteraciones, sólo aumentó el rendimiento del conjunto de entrenamiento, y no
el rendimiento de generalización. Pero todavía, el modelo tiene un rendimiento bastante bueno.  Como hay 
alguna diferencia  entre el rendimiento en el conjunto de  entrenamiento y  prueba, podríamos intentar bajar la complejidad del modelo, para tener mejor rendimiento de generalización. Aquí, elegimos
aumentar el parámetro {\ttfamily alfa} (muy agresivamente, de  0.0001 al 1), para añadir una regularización más fuerte  a los pesos:


\begin{lstlisting}  

In[103]: mlp = MLPClassifier(max_iter=1000, alpha=1, random_state=0)
         mlp.fit(X_train_scaled, y_train)
   
         print("Accuracy on training set: {:.3f}".format(
                 mlp.score(X_train_scaled, y_train)))
         print("Accuracy on test set: {:.3f}".format(
                 mlp.score(X_test_scaled, y_test)))

Out[103]:
     
     Accuracy on training set: 0.988
     Accuracy on test set: 0.972


\end{lstlisting}

Esto conduce a un rendimiento a la par con los mejores modelos hasta ahora{\footnote{\scriptsize{Usted puede haber notado en este punto que muchos de los modelos de buen desempeño alcanzaron exactamente la misma exactitud de 0.972. Esto significa que todos los modelos cometen exactamente el mismo número de errores, que son cuatro. ¡Si comparas las predicciones reales, puedes incluso ver que cometen exactamente los mismos errores! Esto podría ser una consecuencia de que el conjunto de datos sea muy pequeño, o puede ser porque estos puntos son realmente diferentes del resto}}}. \\

Si bien es posible analizar lo que una red neuronal ha aprendido, esto suele ser mucho más complicado que analizar un modelo lineal o un modelo basado en árboles. Una forma de 
analizar lo que fue aprendido por el modelo  es observar los pesos en el modelo.   Puede ver un ejemplo de esto en la  \href{http://scikit-learn.org/stable/auto_examples/neural_networks/plot_mnist_filters.html}{\color{vino}      galería de ejemplo de scikit-learn}.  Para la  dataset de cáncer de mama, esto podría sea un poco difícil de entender. El siguiente gráfico  (figura \ref{ponderaciones}) muestra los pesos que se aprendieron conectando la entrada a la primera capa oculta. Las filas de este gráfico corresponde a las 30 características de entrada,  mientras las columnas  corresponden a las 100 unidades ocultas. Los colores claros representan valores positivos grandes, mientras los colores oscuros representan 
los valores negativos:\\

\begin{lstlisting}  % [frame=single]	

In[104]: plt.figure(figsize=(20, 5))
         plt.imshow(mlp.coefs_[0], interpolation='none', cmap='viridis')
         plt.yticks(range(30), cancer.feature_names)
         plt.xlabel("Columns in weight matrix")
         plt.ylabel("Input feature")
         plt.colorbar()

\end{lstlisting}


\begin{figure}[h!]
	\centering
	\includegraphics[scale=0.50]{../Imagen_Machine/ponderaciones}
	\caption{Mapa caliente de los pesos de la primera capa en una red neuronal aprendida en la dataset de Cáncer de Mama}
	\label{ponderaciones}
\end{figure}


Una posible inferencia que podemos hacer es que las características que tienen pesos muy pequeños para
todas las unidades ocultas son "menos importantes"{}   para el modelo. Podemos ver que   "{}mean smoothness"{}   y   "{}mean compactness"{},  además de las características encontradas entre  "{}smoothness error"{}  y  "{}fractal dimension error"{},  tienen pesos relativamente bajos en comparación con otras características. Esto podría significar que estas son características  son menos importantes  o posiblemente que no los representamos de una manera que la red neuronal podría utilizar.\\

También podríamos visualizar los pesos que conectan la capa oculta a la capa de salida, pero son aún más difíciles de interpretar.   Mientras que {\ttfamily MLPClassifier} y {\ttfamily MLPRegressor} proporcionan interfaces fáciles de usar para arquitecturas de redes neuronales más comunes, solo capturan un pequeño subconjunto de lo que es posible con redes neuronales.  Si estás interesado en trabajar con más flexibilidad
o modelos más grandes, te recomendamos que mires más allá de {\ttfamily scikit-learn,} aprende lo fantástico
 de las  librerías de  deep learning  (librerías de aprendizaje profundo) que están ahí libres. 

Para los usuarios de Python, lo  mejores librerías bien  establecidas son: {\ttfamily keras, lasagna},  y  {\ttfamily tensor-flow.  lasagna} se basa en la librería de {\ttfamily theano}, mientras {\ttfamily Keras} puede usar {\ttfamily tensor-flow  o  theano}.  Estas librerías proporcionan una interfaz mucho más flexible para construir redes neuronales y seguir el rápido progreso en la investigación del aprendizaje profundo. Todas las librerías más  populares de aprendizaje profundo también permiten el uso de unidades de procesamiento gráfico de alto rendimiento (graphics processing units  GPUs), que {\ttfamily scikit-learn} no soporta.  El uso de GPUs nos permite acelerar los cálculos por factores de 10x a 100x, y son esenciales para aplicar métodos de aprendizaje profundo en conjuntos de datos de gran escala.


\subsubsection{Fortalezas, debilidades y parámetros}



Las redes neuronales han resurgido como modelos de  última generación (vanguardia) en muchas 
aplicaciones de aprendizaje automático (machine learning). Una de sus principales ventajas es que pueden capturar información contenida en grandes cantidades de datos y construir modelos increíblemente complejos.
Dando suficiente tiempo de cálculo, datos y ajuste cuidadoso de los parámetros, las redes neuronales generalmente superan  otros algoritmos de aprendizaje automático (para tareas de clasificación y regresión).\\


Por otro lado, esto presenta algunos  inconvenientes o desventajas.
Las redes neuronales, particularmente las grandes y poderosas, a menudo tardan mucho tiempo en entrenarse, además requieren  de un  procesamiento previo a la data  muy cuidadoso, como hemos visto aquí.

De manera similar a los SVMs, funcionan mejor con data  "{}homogénea{},"{ }  donde todas las características tengan significados similares.  Para datos con  características muy diferentes, los modelos basados en árboles podrían funcionar mejor.  En nuestros experimentos, apenas descubrimos la cara sur de posibles formas de ajustar modelos de redes neuronales y cómo entrenarlos.   Configurar los parámetros de una red neuronal también es un arte en sí mismo. En los ejemplos tratados aquí,  apenas introducimos una pequeña parte de las formas posibles  de ajustar  modelos de redes neuronales y cómo entrenarlos.\\


\subsubsection{Estimación de la complejidad en las redes neuronales}

Los parámetros más importantes son el número de capas y el número de unidades ocultas por capa. Debe comenzar con una o dos capas ocultas, y posiblemente expandirse desde allí.  El número de nodos por capa oculta es a menudo similar al número de características de entrada, pero rara vez más alto que en el punto  medio  para  miles.\\

Una medida útil cuando se piensa en la complejidad del modelo de una red neuronal es el número de pesos o coeficientes que aprenden.  Si tiene una  dataset con  100 características para desarrollar una clasificación  binaria, y tiene 100 unidades ocultas, entonces hay 100 *100 = 10.000 pesos entre la entrada y la primera capa oculta. También hay 100 * 1 = 100  pesos entre la capa oculta y la capa de salida, para un total de alrededor de 10.100 pesos. Si agrega una segunda capa oculta con 100 unidades ocultas, allí
habrá otros 100 * 100 = 10.000 pesos de la primera capa oculta hasta la segunda
capa oculta, lo que resulta en un total de 20.100 pesos. 

Si en lugar de eso usas una capa con 1.000 unidades ocultas, tendras   100 * 1.000 = 100.000 pesos  aprendiendo  de la entrada hasta la capa oculta y los pesos de 1000 x 1 de  la capa oculta hasta la capa de salida, para un total de 101.000.  Si añade una segunda capa oculta, entonces se agregan  1.000 * 1.000 = 1.000.000  pesos, para un total  de 1.101.000; este modelo resulta  50 veces más grande que el modelo construido con dos capas ocultas de tamaño 100.\\

Una forma común de ajustar los parámetros en una red neuronal es crear primero una red
 suficientemente grande que genere  sobreajuste en el modelo, claro aquí asegura que la tarea pueda 
ser realmente aprendida por la red, y luego, sabiendo que  los datos de entrenamiento  pueden aprender, se reduce la red o se aumenta el parámetro  {\ttfamily alfa} para añadir regularización, que mejorará el rendimiento de generalización.\\

En estos ejemplos nos enfocamos más en la definición del modelo: el número de
capas y nodos por capa, la regularización y la no linealidad. Estos parámetros  definen el
modelo que queremos aprender. También está la cuestión de cómo aprende el modelo, o  
el algoritmo que se usa para el aprendizaje de los parámetros, el cual es configurado por los parámetro del algoritmo.

El algoritmo  tiene dos opciones fáciles de usar: Uno es el {\ttfamily "adam"{}},  que es el valor por defecto,  este funciona bien en la mayoría de las situaciones, pero es muy  sensible a la escala de los datos (por lo que es importante siempre re-escalar los datos a 0 media y varianza 1).   El otro es {\ttfamily "{}l-bfgs"{}},  el cual es bastante robusto pero puede hacer  un tiempo largo en modelos grandes o dataset  grandes. También está la opción más avanzada  {\ttfamily  "{}sgd"{}},  que es lo que muchos  investigadores del aprendizaje profundo emplean. La opción   {\ttfamily "{}sgd"{}}  viene con muchos parámetros adicionales, que necesitan ser  ajustados  para obtener mejores resultados.  Puede encontrar todos estos parámetros y  sus definiciones en la guía del usuario.  Al comenzar a trabajar con MLPs, le recomendamos  apegarse a {\ttfamily "adam"{}} y {\ttfamily "{}l-bfgs"{}}.\\


\begin{minipage}[t]{0.2\textwidth}
	\centering\raisebox{\dimexpr \topskip-\height}{%
		
\includegraphics[width=\textwidth]{Imagen_Machine/Pajaro.png}}
\end{minipage}%\hfill
\begin{minipage}[t]{0.75\textwidth} 
	{\scriptsize	
	{\bf{fit restablece un modelo (fit Resets a Model)}:}
		
		Una propiedad importante de los modelos de {\ttfamily scikit-learn} es que llamando {\ttfamily fit} siempre restablecerá todo lo que un modelo previamente ha aprendido.   Así que si construye un modelo en un conjunto de datos, y luego llama {\ttfamily fit} de nuevo sobre un conjunto de datos diferente, el modelo "{}olvidará"{}  todo lo que aprendió del primer conjunto de datos. Puede llamar a {\ttfamily fit} tanto  como desee en un modelo, y el resultado será el mismo que llamar {\ttfamily fit} en un "{}nuevo"{}  modelo.}
\end{minipage}

\subsubsection{Estimaciones de incertidumbre de los clasificadores}


Un tema importante y muy útil  del que aún no hemos hablado  en la interfaz de {\ttfamily scikit-learn}, es
la capacidad de los clasificadores de proveer estimaciones de incertidumbre de las predicciones.  A menudo estamos interesados no sólo  en la  clase que predice un clasificador para un cierto punto de prueba, sino también que tan seguro estamos de que la clase es la correcta. En la práctica, diferentes tipos de errores conducen a resultados muy diferentes en aplicaciones del mundo real. Imaginemos  una aplicación médica sobre pruebas sobre una enfermedad, por ejemplo el cáncer.  Hacer una predicción falsa positiva puede llevar a que el paciente se someta a  pruebas adicionales; mientras que una predicción falsa negativa puede llevar al paciente a consecuencias graves al no tratar la enfermedad a tiempo.  Este tema lo  trataremos  con mayor detalle en el Capítulo \ref{Cap6}.\\

 

Hay dos funciones en {\ttfamily scikit-learn} que se pueden utilizar para 
obtener estimaciones de la incertidumbre de los clasificadores, estas son:   {\ttfamily predict\_proba} y  {\ttfamily decision\_function}.   La  mayoría (pero no todos)  de los clasificadores tienen al menos uno de ellos, y muchos clasificadores tienen ambos. \\ 

Analicemos  lo que hacen estas dos funciones sobre la  dataset bidimensional sintético,  cuando se emplea 
 el  clasificador {\ttfamily GradientBoostingClassifier}, que tiene ambas funciones,  la  {\ttfamily decision\_function}  y la  {\ttfamily predict\_proba}:\\


\begin{lstlisting}  % [frame=single]	

In[105]: from sklearn.ensemble import GradientBoostingClassifier
         from sklearn.datasets import make_blobs, make_circles
         X, y = make_circles(noise=0.25, factor=0.5, random_state=1)
   
         # we rename the classes "blue" and "red" for illustration 
         # purposes
         y_named = np.array(["blue", "red"])[y]
   
        # we can call train_test_split with arbitrarily many arrays;
        # all will be split in a consistent manner
        X_train, X_test, y_train_named, y_test_named, y_train, y_test = \
             train_test_split(X, y_named, y, random_state=0)
   
        # build the gradient boosting model
        gbrt = GradientBoostingClassifier(random_state=0)
        gbrt.fit(X_train, y_train_named)


\end{lstlisting} 




\subsubsection*{La función de decisión}

En el caso de clasificación binaria, el valor de retorno de la función\index{función de decisión}   {\ttfamily decision\_function} tiene forma {\ttfamily (n\_samples,)}, y devuelve un número de punto flotante para cada muestra:\\

\begin{lstlisting}  % [frame=single]	

In[106]: print("X_test.shape: {}".format(X_test.shape))
         print("Decision function shape: {}".format(
                gbrt.decision_function(X_test).shape))

Out[106]:

    X_test.shape: (25, 2)
    Decision function shape: (25,)

\end{lstlisting} 


Este valor codifica con qué rigor  el modelo afirma  que un  dato pertenece a la
clase "positiva", en este caso clase 1. Los valores positivos indican una preferencia por la 
clase positiva, y los valores negativos indican una preferencia por la clase "negativa"{ }:\\

\begin{lstlisting}  % [frame=single]	

In[107]: # show the first few entries of decision_function
         print("Decision function:\n{}".format(
                 gbrt.decision_function(X_test)[:6]))
 
Out[107]:

    Decision function:
    [ 4.136 -1.683 -3.951 -3.626 4.29 3.662]
    
\end{lstlisting}


\noindent Podemos recuperar la predicción  sólo observando el signo de la función de decisión:\\
 
	\begin{lstlisting}  % [frame=single]	

In[108]: print("Thresholded decision function:\n{}".format(
         gbrt.decision_function(X_test) > 0))
         print("Predictions:\n{}".format(gbrt.predict(X_test)))

Out[108]:

    Thresholded decision function:
    [ True False False False True True False True True True False True
      True False True False False False True True True True True False
      False]

    Predictions:
    ['red' 'blue' 'blue' 'blue' 'red' 'red' 'blue' 'red' 'red' 'red' 'blue'
     'red' 'red' 'blue' 'red' 'blue' 'blue' 'blue' 'red' 'red' 'red' 'red'
     'red' 'blue' 'blue']

\end{lstlisting} 


Para la clasificación binaria, la clase "negativa"{ }   es siempre la primera entrada del atributo {\ttfamily clases\_}, y la clase "positiva"{ }  es la segunda entrada. Entonces si quieres recuperar completamente la salida de predicción, debe hacer uso del atributo {\ttfamily classes\_}:\\

	\begin{lstlisting}  % [frame=single]	

In[109]: # make the boolean True/False into 0 and 1
         greater_zero = (gbrt.decision_function(X_test) > 0).astype(int)
         # use 0 and 1 as indices into classes_
         pred = gbrt.classes_[greater_zero]
         # pred is the same as the output of gbrt.predict
         print("pred is equal to predictions: {}".format(
                np.all(pred == gbrt.predict(X_test))))

Out[109]:

    pred is equal to predictions: True
\end{lstlisting}

El rango de {\ttfamily decision\_function} puede ser arbitrario, y depende de los datos y los parámetros del modelo:


\begin{lstlisting}  % [frame=single]	

In[110]: decision_function = gbrt.decision_function(X_test)
         print("Decision function minimum: {:.2f} maximum: {:.2f}".format(
                np.min(decision_function), np.max(decision_function)))

Out[110]:

    Decision function minimum: -7.69 maximum: 4.29

\end{lstlisting}

Este escalado arbitrario hace que la salida de {\ttfamily decision\_function} sea un poco difícil para interpretar.\\

En el ejemplo siguiente graficamos  la función {\ttfamily decision\_function} para todos los puntos en el plano 2D, usando una codificación de color  junto a al límite de decisión, como vimos anteriormente. Los  datos  de entrenamiento se presentan como círculos y  de  prueba como  triángulos (Figure \ref{Toy2}):\\

\begin{lstlisting} 

In[111]: fig, axes = plt.subplots(1, 2, figsize=(13, 5))
         mglearn.tools.plot_2d_separator(gbrt, X, 
                                         ax=axes[0], alpha=.4,
                                         fill=True, cm=mglearn.cm2)
         scores_image = mglearn.tools.plot_2d_scores(gbrt, X,
                                                     ax=axes[1],
                                                     alpha=.4, 
                                                     cm=mglearn.ReBl)
                                               
          for ax in axes:
              # plot training and test points
              mglearn.discrete_scatter(X_test[:, 0], X_test[:, 1], 
                                       y_test, markers='^', ax=ax)
              mglearn.discrete_scatter(X_train[:, 0], X_train[:, 1],
                                        y_train, markers='o', ax=ax)
              ax.set_xlabel("Feature 0")
              ax.set_ylabel("Feature 1")
          cbar = plt.colorbar(scores_image, ax=axes.tolist())
          axes[0].legend(["Test class 0", "Test class 1", 
                           "Train class 0","Train class 1"],
                            ncol=4, loc=(.1, 1.1))

\end{lstlisting}



\begin{figure}[h!]
	\centering
	\includegraphics[scale=0.50]{../Imagen_Machine/Toy2}
	\caption{Límite de decisión (izquierda) y función de decisión (derecha) para un modelo de gradient boosting en la dataset de Toy bidimensional.}
	\label{Toy2}
\end{figure}

Codificar no solo el resultado predicho sino también la certeza del clasificador proporciona información adicional. Sin embargo, en esta visualización, es difícil distinguir el límite entre las dos clases.



\subsubsection{Probabilidades de predicción}


La salida de {\ttfamily predic\_proba} es una probabilidad\index{probabilidades de predicción} para cada clase, y  generalmente es más fácil
de entender  que la salida de {\ttfamily decision\_function}. Éste siempre tiene la forma de {\ttfamily (n\_samples, 2)} para clasificación binaria:\\

\begin{lstlisting}

In[112]: print("Shape of probabilities: {}".format(
                 gbrt.predict\_proba(X\_test).shape))

Out[112]:

   Shape of probabilities: (25, 2)

\end{lstlisting}


La primera entrada en cada fila es la probabilidad estimada de la primera clase, y la segunda
 entrada es la probabilidad estimada de la segunda clase. Debido a que es una probabilidad,
la salida de {\ttfamily predic\_proba} siempre está entre 0 y 1, y la suma de las entradas
para ambas clases es siempre 1:\\

\begin{lstlisting}

In[113]: # show the first few entries of predict_proba
         print("Predicted probabilities:\n{}".format(
                gbrt.predict_proba(X_test[:6])))


Out[113]:

   Predicted probabilities:
   [[ 0.016 0.984]
    [ 0.843 0.157]
    [ 0.981 0.019]
    [ 0.974 0.026]
    [ 0.014 0.986]
    [ 0.025 0.975]]
    

\end{lstlisting} 



Debido a que las probabilidades para las dos clases suman 1, exactamente una de las dos  clases
debería estar por encima del 50\% de certeza. Esa clase es la que se predice (las probabilidades son números flotante, puede suceder aunque es poco probable que ambos sean exactamente 0.500; sin embargo, si es el caso, la predicción se hace al azar). Puedes ver en el resultado anterior que el clasificador es relativamente seguro para la mayoría puntos. El grado de incertidumbre reflejada  realmente  depende del modelo y de los parámetros.   Un modelo con  mayor  sobreajuste tiende a hacer predicciones más seguras, incluso si podrían estar erradas. Un modelo con menos complejidad generalmente tiene más   incertidumbre en sus predicciones.    Un modelo se llama calibrado\index{modelo calibrado } si la incertidumbre reporta  realmente  las coincidencias  con lo correcto: esto es en un modelo calibrado, un una predición  hecha con un 70\% de certeza sería correcta el 70\% de las veces.\\ 
 

En el siguiente ejemplo (Figura \ref{Toy3} ) nuevamente mostramos el límite de decisión en el
conjunto de datos, junto a las probabilidades de clase para la clase 1:\\


\begin{lstlisting}  % [frame=single]	

In[114]: fig, axes = plt.subplots(1, 2, figsize=(13, 5))
   
         mglearn.tools.plot_2d_separator(gbrt, X, ax=axes[0], 
                                         alpha=.4, fill=True, 
                                         cm=mglearn.cm2)
         scores_image = mglearn.tools.plot_2d_scores(gbrt, X, ax=axes[1], 
                                                     alpha=.5, cm=mglearn.ReBl,
                                                     function='predict_proba')
       
         for ax in axes:
             # plot training and test points
             mglearn.discrete_scatter(X_test[:, 0], X_test[:, 1],
                                      y_test, markers='^', ax=ax)
             mglearn.discrete_scatter(X_train[:, 0], X_train[:, 1], 
                                      y_train, markers='o', ax=ax)
             ax.set_xlabel("Feature 0")
             ax.set_ylabel("Feature 1")
      cbar = plt.colorbar(scores_image, ax=axes.tolist())
      axes[0].legend(["Test class 0", "Test class 1", "Train class 0",
                      "Train class 1"], ncol=4, loc=(.1, 1.1))

\end{lstlisting} 


\begin{figure}[h!]
	\centering
	\includegraphics[scale=0.50]{../Imagen_Machine/Toy3}
	\caption{Límite de decisión (izquierda) y predición de probabilidades para el modelo gradient 
		boosting mostrado en la Figura \ref{Toy2}}
	\label{Toy3}
\end{figure}





 Los límites en este gráfico están mucho más bien definidos, y las pequeñas áreas de incertidumbre son claramente visibles.\\
 
  El \href{http://bit.ly/2cqCYx6}{\color{vino} sitio web de {\ttfamily scikit-learn}}  tiene una gran comparación de muchos modelos y cómo se ven sus estimaciones de incertidumbre. Hemos reproducido esto en la Figura \ref{Incerti}, y le recomendamos seguir a través del ejemplo.

\begin{figure}[h!]
	\centering
	\includegraphics[scale=0.55]{../Imagen_Machine/Incerti}
	\caption{Comparación de varios clasificadores en scikit-learn sobre la data de sintéticos (imagen cortesía de  \href{http://scikit-learn.org\%29/}{\color{vino} http://scikit-learn.org})}
	\label{Incerti}
\end{figure}


\subsubsection*{Incertidumbre en la clasificación multiclase}

Hasta ahora, sólo hemos hablado de estimaciones de incertidumbre en la clasificación binaria. Pero los métodos {\ttfamily decision\_function}  y  {\ttfamily predict\_proba}  también funcionan en el entorno multiclase. \\

Ilustremos con la data  Iris, que es una data  de clasificación en tres clases:



\begin{lstlisting}  % [frame=single]	


In[115]: from sklearn.datasets import load_iris
   
         iris = load_iris()
         X_train, X_test, y_train, y_test = train_test_split(
               iris.data, iris.target, random_state=42)
       
         gbrt = GradientBoostingClassifier(learning_rate=0.01,
                                            random_state=0)
         gbrt.fit(X_train, y_train)
   
In[116]: print("Decision function shape: {}".format(
                gbrt.decision_function(X_test).shape))
         # plot the first few entries of the decision function
         print("Decision function:\n{}".format(
                 gbrt.decision_function(X_test)[:6, :]))
 
Out[116]:
    
    Decision function shape: (38, 3)
    Decision function:
    [[-0.529 1.466 -0.504]
     [ 1.512 -0.496 -0.503]
     [-0.524 -0.468 1.52 ]
     [-0.529 1.466 -0.504]
     [-0.531 1.282 0.215]
     [ 1.512 -0.496 -0.503]]

\end{lstlisting} 


En el caso multiclase, la {\ttfamily decision\_function}  tiene la forma {\ttfamily  (n\_samples, n\_classes)} y cada columna proporciona una "{}puntuación de certeza"{ }  para cada clase, donde una gran puntuación  significa que una clase es más probable y una puntuación pequeña significa que la clase es menos probable. Usted puede recuperar las predicciones de estas puntuaciones mediante la búsqueda de la entrada máxima para cada punto de datos:\\


\begin{lstlisting}  % [frame=single]	

In[117]: print("Argmax of decision function:\n{}".format(
                np.argmax(gbrt.decision_function(X_test), axis=1)))
         print("Predictions:\n{}".format(gbrt.predict(X_test)))

Out[117]:
   
    Argmax of decision function:
    [1 0 2 1 1 0 1 2 1 1 2 0 0 0 0 1 2 1 1 2 0 2 0 2 2 2 2 2 0 0 0 0 1 0 0 2 1 0]
    Predictions:
    [1 0 2 1 1 0 1 2 1 1 2 0 0 0 0 1 2 1 1 2 0 2 0 2 2 2 2 2 0 0 0 0 1 0 0 2 1 0]


\end{lstlisting} 


La salida de {\ttfamily predict\_proba} tiene la misma forma, {\ttfamily(n\_samples, n\_classes)}. Una vez más, las probabilidades de las clases posibles para cada punto de datos suma a 1:



\begin{lstlisting}  % [frame=single]	

In[118]: # show the first few entries of predict_proba
         print("Predicted probabilities:\n{}".format(
                  gbrt.predict_proba(X_test)[:6]))
         # show that sums across rows are one
         print("Sums: {}".format(
                 gbrt.predict_proba(X_test)[:6].sum(axis=1)))

Out[118]:

    Predicted probabilities:
    [[ 0.107 0.784 0.109]
     [ 0.789 0.106 0.105]
     [ 0.102 0.108 0.789]
     [ 0.107 0.784 0.109]
     [ 0.108 0.663 0.228]
     [ 0.789 0.106 0.105]]
Sums: [ 1. 1. 1. 1. 1. 1.]

\end{lstlisting} 

Podemos recuperar de nuevo las predicciones mediante el cálculo del {\ttfamily argmax}  de {\ttfamily predict\_proba}:


\begin{lstlisting}  % [frame=single]	


In[119]: print("Argmax of predicted probabilities:\n{}".format(
         np.argmax(gbrt.predict_proba(X_test), axis=1)))
         print("Predictions:\n{}".format(gbrt.predict(X_test)))

Out[119]:
   
    Argmax of predicted probabilities:
    [1 0 2 1 1 0 1 2 1 1 2 0 0 0 0 1 2 1 1 2 0 2 0 2 2 2 2 2 0 0 0 0 1 0 
     0 2 1 0]
    Predictions:
    [1 0 2 1 1 0 1 2 1 1 2 0 0 0 0 1 2 1 1 2 0 2 0 2 2 2 2 2 0 0 0 0 1 0 
     0 2 1 0]

\end{lstlisting} 


Para resumir, la función {\ttfamily predict\_proba}  y  {\ttfamily\_function} siempre tienen forma {\ttfamily (n\_sam ples, n\_classes)}, aparte de {\ttfamily decision\_function} en el caso especial  binario, {\ttfamily decision\_function} sólo tiene una columna, que corresponde a la clase  "positive"{}    
{\ttfamily classes\_[1] }. Esto es principalmente por razones históricas.\\

 Se puede recuperar la predicción cuando hay {\ttfamily n\_classes} (muchas columnas) mediante el cálculo del {\ttfamily argmax} a través de las columnas. Sin embargo, tenga cuidado, si sus clases son strings (cadenas de texto), o usa números enteros, pero no son consecutivos y a partir de 0. Si desea comparar los resultados obtenidos con el predict con los resultados obtenidos a través de {\ttfamily decision\_function} o pre {\ttfamily dict\_proba}, asegúrese de usar el atributo {\ttfamily classes\_} del clasificador para obtener los nombres de clase reales:


\begin{lstlisting}  % [frame=single]	

In[120]: logreg = LogisticRegression()
         # represent each target by its class name in the iris dataset
         named_target = iris.target_names[y_train]
         logreg.fit(X_train, named_target)
         print("unique classes in training data: {}".format(
                logreg.classes_))
         print("predictions: {}".format(logreg.predict(X_test)[:10]))
         argmax_dec_func = np.argmax(logreg.decision_function
                                     X_test), axis=1)
         print("argmax of decision function: {}".format(
                argmax_dec_func[:10]))
         print("argmax combined with classes_: {}".format(
                logreg.classes_[argmax_dec_func][:10]))
   
Out[120]:

    unique classes in training data: ['setosa' 'versicolor' 'virginica']
    predictions: ['versicolor' 'setosa' 'virginica' 'versicolor' 
     'versicolor' 'setosa' 'versicolor' 'virginica' 'versicolor'
     'versicolor']
    argmax of decision function: [1 0 2 1 1 0 1 2 1 1]
    argmax combined with classes_: ['versicolor' 'setosa' 'virginica' 
     'versicolor' 'versicolor' 'setosa' 'versicolor' 'virginica' 'versicolor'
     'versicolor']

\end{lstlisting} 



\subsubsection{Resumen y perspectivas}\ref{Pagina127} 

Comenzamos este capítulo con una discusión sobre la complejidad del modelo, luego discutimos 
generalización, o aprendizaje de un modelo que puede rendir (desenpeño) bien con datos nuevos, nunca antes vistos. Esto nos llevó a los conceptos de falta de ajuste, que describe un modelo que no puede
capturar las variaciones presentes de los datos de entrenamiento,  y el sobreajuste que describe un
modelo que se enfoca demasiado en los datos de entrenamiento y no puede generalizar bien con datos nuevos.

Luego discutimos una amplia gama de modelos de aprendizaje automático (machine learning) para la clasificación y regresión, con sus ventajas y desventajas, y cómo controlar 
complejidad  del modelo para cada uno de ellos.  Vimos que para muchos de los algoritmos, configurar los  parámetros correctos  es muy importante para un buen rendimiento.   Algunos de los algoritmos también son sensibles a cómo representamos los datos de entrada, y en particular a la escala de las características. Por lo tanto, aplicar a ciegas  un algoritmo a un conjunto de datos sin comprender los supuestos que hace el modelo y los significados de la configuración de los parámetros rara vez conducen a un modelo exacto.\\


Este capítulo contiene mucha información sobre los algoritmos, y no es necesario que recuerdes todos estos detalles para los capítulos.  Sin embargo, algún conocimiento de los modelos descritos aquí, y cuáles usar en una situación específica es importante para aplicar con éxito el aprendizaje automático en la práctica. Aquí hay un resumen rápido de cuándo usar cada modelo\index{cuando usar el modelo}:\\



{\bf Vecinos más cercanos }

Para dataset pequeños, buenos como línea de base, fáciles de explicar.\\

{\bf Modelos lineales}

Va  como un primer algoritmo para probar, bueno para conjuntos de datos muy grandes,  y para datos de muy alta dimensión.\\


{\bf Naives Bayes}

 Solo para clasificación. Incluso más rápido que los modelos lineales, bueno para datos muy grandes.
conjuntos y datos de alta dimensión. A menudo menos exacto que los modelos lineales.\\


{ \bf Árboles de decisión}

Muy rápido, no necesita escalar los datos, puede visualizarse y explicarse fácilmente.\\

{\bf Bosques aleatorios}

Casi siempre se desempeña mejor que un solo árbol de decisión, muy robusto y potente. No es necesario escalar los datos. No es bueno para data sparse de muy alta dimensión.\\


{\bf Gradiente potenciado}

A menudo un poco más exacto  que los bosques aleatorios. Más lento para entrenar pero más rápido para
predecir que los bosques aleatorios, y menos  memoria. Necesita más ajuste de parámetros
 que bosques aleatorios.\\
 
{\bf Máquinas de Soporte de vectores}

Potente para  dataset de tamaño medio de características  y  con un significado similar. Exige
re-escalar los datos, sensible a los parámetros.\\

{\bf Redes neuronales}

Puede construir modelos muy complejos, particularmente para datas muy  grandes. Sensible al re-escalado
de los datos y a la elección de parámetros. Los modelos grandes necesitan mucho tiempo para
entrenar.\\



Cuando se trabaja con una data nueva es una buena idea comenzar con un modelo simple, como el lineal, el  clasificador  naives Bayes    o  con vecinos más cercanos, y ver
qué tan lejos se puede llegar con ellos. Después de comprender más sobre los datos, puede considerar
pasar a un algoritmo que pueda construir modelos más complejos, como bosques aleatorios,
gradiente potenciado, SVMs,  o redes neuronales.\\


Con esta información dada Usted debe  estar en capacidad  de  aplicar, ajustar y
analizar los modelos que discutimos aquí. En este capítulo, nos centramos en el caso de la clasificación binaria, por  ser más fácil de entender. La mayoría de los algoritmos presentandos
tiene variantes de clasificación y regresión,  y todos los algoritmos de  clasificación  admiten la clasificación binaria y multiclase. 
Recomendamos aplicar cualquiera de estos algoritmos a las datas  integrados en {\ttfamily  scikit-learn}, como {\ttfamily boston\_housing} o {\ttfamily diabetes} para regresión, o el conjunto de datos de dígitos para clasificación multiclase. \\

Tratar  los algoritmos vistos  en  diferentes datas  le dará una mejor percepción a cerca de cuánto tiempo  necesitan para entrenar, qué tan fácil es analizar los modelos,  y las opciones  posibles  para la representación de los datos,  en general adquirirá mayor destrezas en el uso y aplicación de los modelos.\\


También  se analizó  las consecuencias  que generan los  ajustes  de los parámetro  de  los  algoritmos;  este tema se tratará con más énfasis   en el   Capítulo \ref{Cap6}, donde  se   estudian algunas técnicas  que permiten encontrar   y ajustar   automáticamente  los mejores valores para los  parámetros del modelo.\\ 


En el próximo capítulo  será dedicado al  aprendizaje no supervisado,   y se introduce   con más detalle  el pre-procesamiento  de datos  necesario para algunos  algoritmos  muy  sensibles  a la escala  de los datos,  como redes neuronales  y  SVM.   

